\chapter{Energy and finite-energy potentials}
\index{Monge--Ampere energy!BFJ}\index{finite-energy potential!BFJ}
\label{chap:bfj-energy}

The Monge--Amp\`ere operator is nonlinear, but it has a primitive.  This
simple fact reorganizes the theory: comparison becomes concavity, weak
compactness can be refined by one scalar functional, and unbounded potentials
can be admitted exactly as long as their energy does not fall to
\(-\infty\).

\section{Orientation}
\label{sec:ch26-orientation}

\subsection{Why the operator has a primitive, and what that buys}
\label{sec:ch26-why}

\Cref{chap:bfj-ma} produced, for each \(n\)-tuple of bounded
\(\omega\)-psh functions, a positive measure of total mass \(V\).  The
resulting map
\[
 \varphi\longmapsto\MA_\omega(\varphi)
 =(\omega+\ddc\varphi)^n
\]
is nonlinear of degree \(n\), and by itself it does not organize the space of
potentials at all.  What organizes it is the following observation, which is
purely formal at model level and becomes a theorem as soon as the potentials
are allowed to be singular.

Fix a segment \(\varphi_t=\varphi+t(\psi-\varphi)\) of psh functions.  Then
\begin{equation}
\label{eq:ch26-motivating-variation}
 \frac{\rd}{\rd t}\int\varphi_t\,\MA_\omega(\varphi_t)
 \neq\int(\psi-\varphi)\,\MA_\omega(\varphi_t),
\end{equation}
because differentiating the integrand and the measure both contribute; but
the \emph{averaged} functional
\[
 E_\omega(\varphi)
 =\frac1{n+1}\sum_{j=0}^n\int\varphi\,
   \omega_\varphi^j\wedge\omega^{n-j}
\]
does have derivative \(\int(\psi-\varphi)\MA_\omega(\varphi_t)\), by an exact
cancellation produced by integration by parts.  In other words, the
Monge--Amp\`ere operator is the gradient of a single scalar functional.  The
combinatorial weight \(\tfrac1{n+1}\) is not cosmetic: it is exactly what
makes the cancellation work, and every sign and constant in
\Cref{part:bfj} is calibrated to it in \Cref{sec:appJ-energy}.

Three consequences follow immediately, and they are the reason this chapter
sits where it does.

\emph{Comparison becomes concavity.}  The second derivative of
\(t\mapsto E_\omega(\varphi_t)\) is a Dirichlet-type integral, negative by
the relative Hodge index theorem \Cref{thm:ch20-mixed-hodge}.  Consequently
\(E_\omega\) is concave, and the two supporting-line inequalities that
concavity supplies are precisely the sandwich estimates that replace, in the
global theory, the local comparison principles of \Cref{chap:bfj-capacity}.

\emph{Weak compactness acquires a scale.}  The space
\(\PSH(X,\omega)/\mathbf R\) is compact
(\Cref{thm:bfj-psh-compactness}), but the compactness is soft: nothing in it
distinguishes a potential with a mild singularity from one whose
Monge--Amp\`ere mass has escaped entirely into a polar set.  Energy is the
single real-valued functional that makes that distinction, and its normalized
superlevel sets \eqref{eq:normalized-energy-superlevel} are the compact sets
on which the variational method of
\Cref{chap:bfj-differentiability,chap:bfj-finite-energy-measures} operates.

\emph{Unbounded potentials become admissible.}  The bounded theory of
\Cref{chap:bfj-ma} cannot see a potential with a pole.  Truncation
\(\varphi\mapsto\max\{\varphi,-t\}\) always produces bounded approximants, but
their Monge--Amp\`ere measures need not converge to anything useful: mass may
run off to the pole set and be lost.  Finiteness of energy is exactly the
condition under which it does not, and the extension of the mixed
Monge--Amp\`ere product to that class, proved in
\Cref{sec:ch26-mixed-products}, is the technical heart of the chapter.

\subsection{The payoff, stated in advance}
\label{sec:ch26-payoff}

Six statements proved below are used throughout
\Crefrange{chap:bfj-differentiability}{chap:bfj-examples}, and the reader may
wish to see them before the constructions.

\begin{enumerate}
\item\label{itm:ch26-payoff:cocycle}
 \emph{The energy difference of two potentials is a single explicit sum of
 mixed integrals} (\Cref{prop:energy-cocycle},
 \Cref{thm:ch26-mixed-energy-identities}).  All later variational arguments
 are manipulations of that sum; no other formula for \(E_\omega\) is needed.
\item\label{itm:ch26-payoff:concave}
 \emph{Energy is increasing and concave, with an explicit negative second
 variation} (\Cref{cor:energy-monotone-concave}).  This is the analytic face
 of the relative Hodge index inequality
 \eqref{eq:ch20-mixed-hodge}, and it is the reason the variational problem of
 \Cref{chap:bfj-finite-energy-measures} is a concave maximization problem
 rather than a critical-point problem.
\item\label{itm:ch26-payoff:usc}
 \emph{Energy is upper semicontinuous for the weak topology and continuous
 along decreasing nets} (\Cref{thm:extended-energy-basic}).  These are two
 different statements; \Cref{ex:ch26-usc-strict} exhibits a sequence for
 which the first holds strictly and the second does not apply.
\item\label{itm:ch26-payoff:full-mass}
 \emph{A potential of finite energy has full Monge--Amp\`ere mass, with a
 quantitative tail bound} (\Cref{prop:e1-integrability-tail}).  The tail bound
 \eqref{eq:e1-mass-tail} is what makes every later truncation argument
 uniform.
\item\label{itm:ch26-payoff:mixed}
 \emph{Mixed Monge--Amp\`ere products extend from bounded potentials to
 \(\cE^1\), keeping total mass \(V\)}
 (\Cref{thm:e1-mixed-products}).  The extension is a theorem with a
 quantitative proof (\Cref{lem:ch26-mixed-tail}), not a definition by fiat,
 and it is the statement that the delimitation table of
 \Cref{sec:ch16-delimitation-table} refers to.
\item\label{itm:ch26-payoff:compact}
 \emph{Normalized energy superlevels are weakly compact}
 (\Cref{prop:e1-compact-superlevels}), in net form
 (\Cref{cor:ch26-net-superlevel}).
\end{enumerate}

A seventh statement is negative and is worth flagging with the others.  The
natural translation-invariant gauge built from energy and the maximum
operation does \emph{not} satisfy the triangle inequality:
\Cref{ex:ch26-triangle-fails} violates it by a factor of two on a metric
graph with two vertices.  The strong topology of
\Cref{sec:ch26-strong} is therefore introduced as a topology, not as a
metric, and no completeness statement is smuggled in with the word
``strong''.

\subsection{Standing hypotheses and notation}
\label{sec:ch26-hypotheses}

The hypotheses are those of \eqref{eq:ch20-standing}: \(K\) is complete and
discretely valued with residue field of characteristic zero, \(X\) is a
smooth connected projective \(K\)-variety of dimension \(n\geq1\), and models
are normal, flat and projective over \(R=K^\circ\).  We fix once and for all a
closed model form
\[
 \omega\in\mathcal Z^{1,1}(X)
\]
which is \emph{semipositive} and whose generic numerical class \(\{\omega\}\)
is ample, and we write
\begin{equation}
\label{eq:ch26-standing-notation}
 V:=(\{\omega\}^n)>0,
 \qquad
 \omega_\varphi:=\omega+\ddc\varphi .
\end{equation}

\begin{remark}[The reference form is called \(\theta\) in
\Cref{chap:bfj-model-forms,chap:bfj-ma}]
\label{rem:ch26-omega-theta}
\Crefrange{chap:bfj-model-forms}{chap:bfj-ma} write \(\theta\) for the
reference closed model form, whereas \Cref{chap:bfj-capacity} and the present
chapter write \(\omega\).  The two letters denote the same kind of object;
the change is made where the form ceases to be a variable of the discussion
and becomes a fixed background.  Every statement quoted from
\Cref{chap:bfj-model-forms,chap:bfj-ma} is applied with \(\theta=\omega\), and
the reader who is checking a normalization should perform that substitution
before comparing signs.  By \Cref{prop:ch20-positive-representative} one may
in addition assume, after replacing \(\omega\) by \(\omega+\ddc\psi\) and
every potential \(\varphi\) by \(\varphi-\psi\), that \(\omega\) is
\(\mathcal X\)-positive on some model; this normalization is used only where
it is explicitly invoked, and by
\Cref{cor:ch21-reference-independence} it changes no energy difference.
\end{remark}

Energy is normalized by \(E_\omega(0)=0\), that is, relative to the reference
form itself.  Since all Monge--Amp\`ere measures in sight have total mass
\(V\), translation by a constant satisfies
\begin{equation}
\label{eq:bfj-energy-translation}
                    E_\omega(\varphi+c)=E_\omega(\varphi)+cV .
\end{equation}
This factor of \(V\) is easy to lose when reading a source normalized to
probability mass, and it is the first of the two unit tests of
\eqref{eq:appJ-energy-translation}; the second is the first-variation formula
\eqref{eq:appJ-energy-variation}.

Two further pieces of notation are fixed here and used without comment.  For
\(\varphi\in\PSH(X,\omega)\) and \(t>0\) we write
\begin{equation}
\label{eq:ch26-truncation}
 \varphi^{\langle t\rangle}:=\max\{\varphi,-t\},
\end{equation}
the \emph{truncation of \(\varphi\) at level \(-t\)}; it is
\(\omega\)-psh by \Cref{thm:ch21-stability}\,\cref{itm:ch21-stab:max} together
with \Cref{prop:ch21-constants}\,\cref{itm:ch21-const:const}, and it is bounded
as soon as \(\varphi\) is bounded above, which we arrange by translating so
that \(\varphi\leq0\).  For a multi-index we write
\(\MA_\omega(\varphi^{(j)},\psi^{(n-j)})\) for the mixed measure with
\(j\) entries \(\varphi\) and \(n-j\) entries \(\psi\), consistently with
\eqref{eq:appJ-energy}.

\subsection{Four Monge--Amp\`ere products, kept apart}
\label{sec:ch26-four-products}

The governing convention of this book is that model, continuous, bounded and
finite-energy Monge--Amp\`ere products are never conflated.  This chapter is
where the last of the four is constructed, so we restate the hierarchy once,
in the form in which it will be used.

\begin{center}
\small
\begin{tabular}{p{.19\textwidth}|p{.24\textwidth}|p{.24\textwidth}|p{.24\textwidth}}
\toprule
\textbf{Product}&\textbf{Domain}&\textbf{Definition}&\textbf{Constructed in}\\
\midrule
model
 &\(\omega\)-psh model functions
 &vertical intersection numbers on one model
 &\Cref{def:mixed-model-measure}, \eqref{eq:model-mixed-measure}\\
continuous
 &continuous semipositive metrics
 &uniform limit of model products
 &\Cref{thm:bfj-cld-ma-common-locus}\\
bounded
 &bounded \(\omega\)-psh functions
 &unique monotone extension
 &\Cref{thm:bounded-ma-extension}\\
finite energy
 &\(\cE^1(X,\omega)\), possibly unbounded
 &increasing limit of restricted truncations
 &\Cref{def:nonpluripolar-ma}, \Cref{thm:e1-mixed-products}\\
\bottomrule
\end{tabular}
\end{center}

Each row is a genuinely different construction, and each arrow between rows is
a theorem.  The passage from the model row to the continuous row is uniform
approximation; from the continuous row to the bounded row is the monotone
extension of \Cref{thm:bounded-ma-extension}; and from the bounded row to the
finite-energy row is \Cref{thm:e1-mixed-products} of this chapter, whose proof
occupies \Cref{sec:ch26-mixed-products}.  Every numbered statement below says
which of the four products it concerns, and \Cref{sec:ch26-scope} records
where each hypothesis is used.

\begin{warning}{This chapter is one side of a boundary drawn elsewhere}
The relation between the products of this chapter and the local
Bedford--Taylor products of the Chambert--Loir--Ducros calculus is delimited,
feature by feature and with a table, in
\Cref{sec:ch16-local-vs-global,sec:ch16-delimitation-table}.  We do not redraw
that boundary here.  Two points from it are used repeatedly and are worth
recalling: the CLD product is local and its total mass is not an intersection
number, whereas the products of this chapter are global and have total mass
\(V\) by construction; and the CLD product may carry mass on a polar set
(\Cref{ex:ch16-polar-failure}), whereas the finite-energy product of
\Cref{def:nonpluripolar-ma} discards that mass by definition.  Conflating the
two would produce a wrong total mass in \Cref{chap:compare-ma}.
\end{warning}

\subsection{Extended-real protocol}
\label{sec:ch26-extended-real}

Potentials in this chapter take the value \(-\infty\), energies take the value
\(-\infty\), and the dual energies of
\Cref{chap:bfj-finite-energy-measures} take the value \(+\infty\).  The
arithmetic of these values is fixed once and for all by the ledger of
\Cref{sec:extended-arithmetic}, summarized in
\Cref{tab:extended-real-ledger}, and we use it rather than improvising.  Four
entries are in constant use, and it is worth saying which.

\begin{enumerate}
\item\label{itm:ch26-ledger:add}
 \emph{Addition of a finite constant} is an order isomorphism of \(\Rbar\), so
 \(\varphi+c\) is defined on the polar locus and
 \eqref{eq:bfj-energy-translation} makes sense for every
 \(\varphi\in\PSH(X,\omega)\).  This is entry~(2) of the ledger.
\item\label{itm:ch26-ledger:diff}
 \emph{Differences are never formed on a common pole set.}  The expression
 \(\psi-\varphi\) is undefined at a point where both potentials equal
 \(-\infty\), by entry~(4).  Every integral below involving that difference is
 obtained either from bounded truncations or from a monotone identity in which
 no subtraction occurs; the reduction is carried out explicitly each time.
 \Cref{rem:forbidden-cancellations} lists the three cancellations that fail.
\item\label{itm:ch26-ledger:integral}
 \emph{Integrals} are formed only when one of \(f^\pm\) has finite integral,
 by entry~(7).  In particular, if \(\varphi\leq0\) is Borel and \(\mu\) is a
 finite positive measure, then \(\int\varphi\,\rd\mu\in[-\infty,0]\) is always
 defined, and
 \(\int\varphi^{\langle t\rangle}\rd\mu\searrow\int\varphi\,\rd\mu\) by
 \Cref{lem:integration-ledger}\,(2).  This is the mechanism behind every
 truncation limit in the chapter.
\item\label{itm:ch26-ledger:null}
 \emph{A null set carries no mass}, even where the integrand is \(-\infty\),
 by \Cref{lem:integration-ledger}\,(3).  This is what allows
 \(\int(-\varphi)\MA_\omega(\varphi)\) to be finite for
 \(\varphi\in\cE^1\) although \(\varphi=-\infty\) somewhere.
\end{enumerate}

\subsection{What is proved here and what is imported}
\label{sec:ch26-imported}

Everything in \Crefrange{sec:ch26-model-energy}{sec:ch26-strong} is proved,
with the following exceptions, each quoted with its hypotheses at the point of
use.

\begin{enumerate}
\item\label{itm:ch26-import:hodge}
 The relative Hodge index inequality \eqref{eq:ch20-mixed-hodge} and its
 Cauchy--Schwarz form \eqref{eq:ch20-mixed-cauchy-schwarz}, proved in
 \Cref{chap:bfj-model-forms}.  This is the one place where the sign of the
 second variation comes from, and it is an algebraic input.
\item\label{itm:ch26-import:bounded}
 The bounded monotone extension \Cref{thm:bounded-ma-extension} and its
 corollaries \Cref{cor:bounded-ma-multilinearity},
 \Cref{prop:bounded-ma-no-polar-mass}, proved in \Cref{chap:bfj-ma}.
\item\label{itm:ch26-import:capacity}
 The bounded maximum principle \eqref{eq:bounded-maximum-principle}, the
 bounded comparison principle \eqref{eq:bounded-comparison-principle}, the
 capacity comparison \eqref{eq:capacity-comparison} and the domination
 statement of \Cref{prop:bounded-capacity-domination}, proved in
 \Cref{chap:bfj-capacity}.
\item\label{itm:ch26-import:compactness}
 Compactness of \(\PSH_0(X,\omega)\) (\Cref{thm:bfj-psh-compactness}), its net
 form \Cref{cor:ch22-net-compactness}, and decreasing regularization
 \Cref{thm:bfj-decreasing-regularization}, proved in
 \Cref{chap:bfj-compactness}.  Residue characteristic zero enters the theory
 through the last of these and nowhere else in this chapter.
\item\label{itm:ch26-import:convex}
 For the toric computations only: the Alexandrov measure
 \Cref{thm:appD-alexandrov-measure}, the Legendre pushforward formula
 \eqref{eq:appD-legendre-pushforward}, and the toric comparison
 \Cref{thm:toric-bfj-alexandrov}.
\end{enumerate}

No result of \Cref{part:cld} is used.  In particular the finite-energy product
constructed here is not asserted to agree with any local Bedford--Taylor
product; the loci on which such an agreement is proved are exactly those of
\Cref{sec:ch16-transfer} and \Cref{thm:bfj-cld-ma-common-locus}.

\subsection{Plan}
\label{sec:ch26-plan}

\Cref{sec:ch26-model-energy} defines the energy of a model function, identifies
it with a mixed vertical intersection sum, and proves the cocycle formula,
monotonicity and concavity; it then transports all of this to bounded
potentials.  \Cref{sec:ch26-mixed} isolates the mixed-energy functionals
\(I\) and \(J\), proves the identities relating them by the integration-by-parts
chain, and records exactly at which regularity each integration by parts is
legitimate.  \Cref{sec:ch26-extension} extends \(E_\omega\) to arbitrary psh
functions and proves upper semicontinuity and monotone continuity;
\Cref{sec:ch26-nonpluripolar} constructs the nonpluripolar measure and proves
the finite-energy tail estimates.  \Cref{sec:ch26-mixed-products} proves the
extension of the mixed products to \(\cE^1\).  \Cref{sec:ch26-e1-gauge}
studies the order and gauge structure of \(\cE^1\), including the failure of
the triangle inequality, and \Cref{sec:ch26-strong} treats strong convergence
and compactness of normalized sublevels.  \Cref{sec:ch26-examples} carries six
computations to explicit numbers, \Cref{sec:ch26-scope} is the scope table, and
\Cref{sec:ch26-dictionary} is the chapter dictionary.

\section{Energy on the model and bounded loci}
\label{sec:ch26-model-energy}

\subsection{The definition and its intersection-theoretic meaning}
\label{sec:ch26-model-energy-def}

\begin{definition}[Monge--Amp\`ere energy]
\label{def:model-energy}
For a \(\omega\)-psh model function \(\varphi\), set
\begin{equation}
\label{eq:model-energy}
 E_\omega(\varphi)
 :=\frac1{n+1}\sum_{j=0}^n
   \int_{X^{\An}}\varphi\,
   \omega_\varphi^j\wedge\omega^{n-j}.
\end{equation}
\end{definition}

The integrand is a model function and each measure
\(\omega_\varphi^j\wedge\omega^{n-j}\) is a positive atomic measure of mass
\(V\), by \Cref{prop:ch20-measure-positivity}; so the sum is a finite sum of
finite numbers and nothing has to be justified about its convergence.  The
normalization \(E_\omega(0)=0\) is immediate, and
\eqref{eq:bfj-energy-translation} follows because each of the \(n+1\)
summands changes by \(cV/(n+1)\).

Formula \eqref{eq:model-energy} is a sum of \(n+1\) integrals.  The next
proposition records the equivalent vertical intersection formula.  Every
term has an actual vertical-divisor factor; no absolute
\((n+1)\)-fold intersection of closed model forms is used.

\begin{proposition}
\label{prop:ch26-energy-intersection}
\index{Monge--Ampere energy!as intersection number}
Let \(\varphi\) be an \(\omega\)-psh model function.  Then, with the pairing
of \Cref{prop:ch20-intersection-pairing},
\begin{equation}
\label{eq:ch26-energy-intersection}
 E_\omega(\varphi)
 =\frac1{n+1}\sum_{j=0}^n
   \bigl(D_\varphi\cdot\omega_\varphi^{\,j}\cdot
                         \omega^{\,n-j}\bigr),
\end{equation}
where \(D_\varphi\) is the vertical divisor representing \(\varphi\) on a
common model.
More generally, for \(\omega\)-psh model functions \(\varphi,\psi\),
\begin{equation}
\label{eq:ch26-energy-difference-intersection}
 E_\omega(\psi)-E_\omega(\varphi)
 =\frac1{n+1}\sum_{j=0}^n
   \bigl(D_{\psi-\varphi}\cdot\omega_\psi^{\,j}\cdot
                         \omega_\varphi^{\,n-j}\bigr).
\end{equation}
\end{proposition}

\begin{proof}
On a common determining model,
\eqref{eq:ch20-pairing-as-integral} gives
\[
 \int_{X^{\An}}\varphi\,
       \omega_\varphi^{\,j}\wedge\omega^{\,n-j}
 =\bigl(D_\varphi\cdot\omega_\varphi^{\,j}
                         \cdot\omega^{\,n-j}\bigr).
\]
Summing this identity and using \eqref{eq:model-energy} proves
\eqref{eq:ch26-energy-intersection}.

For the difference formula, put \(u:=\psi-\varphi\) and
\(\varphi_t:=\varphi+tu\), so that
\(\omega_t:=\omega_{\varphi_t}=\omega_\varphi+t\ddc u\).  We first prove
the formula analytically, keeping \(u\) itself rather than only its
\(\ddc\)-class; this point includes constant potentials.  Differentiating
the finite sum in \eqref{eq:model-energy} and using integration by parts on
a common determining model gives
\[
\begin{aligned}
 (n+1)\frac{\rd}{\rd t}E_\omega(\varphi_t)
 ={}&\sum_{j=0}^n\int u\,\omega_t^{\,j}\wedge\omega^{\,n-j}\\
 &+\sum_{j=1}^n j\int \varphi_t\,\ddc u\wedge
             \omega_t^{\,j-1}\wedge\omega^{\,n-j}\\
 ={}&\sum_{j=0}^n\int u\,\omega_t^{\,j}\wedge\omega^{\,n-j}\\
 &+\sum_{j=1}^n j\int u\,(\omega_t-\omega)\wedge
             \omega_t^{\,j-1}\wedge\omega^{\,n-j}\\
 ={}&(n+1)\int u\,\omega_t^{\,n}.
\end{aligned}
\]
The last equality is telescoping: if
\(I_j=\int u\,\omega_t^j\wedge\omega^{n-j}\), the two sums are
\(\sum_{j=0}^n I_j+\sum_{j=1}^n j(I_j-I_{j-1})=(n+1)I_n\).
Integrating from \(0\) to \(1\), expanding
\(\omega_t^n=((1-t)\omega_\varphi+t\omega_\psi)^n\), and using
\[
 \binom nk\int_0^1t^k(1-t)^{n-k}\,\rd t=\frac1{n+1}
 \qquad(0\leq k\leq n),
\]
we obtain
\[
 E_\omega(\psi)-E_\omega(\varphi)
 =\frac1{n+1}\sum_{j=0}^n
   \int u\,\omega_\psi^{\,j}\wedge\omega_\varphi^{\,n-j}.
\]
Finally \eqref{eq:ch20-pairing-as-integral} changes each summand into the
corresponding term with the actual vertical divisor \(D_u\), proving
\eqref{eq:ch26-energy-difference-intersection}.  Independence of the common
model is \Cref{prop:ch20-intersection-pairing}.
\end{proof}

\begin{remark}[Two readings of \eqref{eq:ch26-energy-intersection}]
\label{rem:ch26-energy-intersection-reading}
Read from left to right, \eqref{eq:ch26-energy-intersection} says that the
analytic functional \(E_\omega\) is computed by a finite sum of vertical
intersection numbers.  Read from right to left, the mixed vertical sum is a
\emph{primitive} of the Chambert--Loir measure: its first variation leaves
only \((D\cdot\omega_\varphi^n)\).  This is exactly the energy functional
introduced in \Cref{cor:ch20-concavity}, whose second variation has the
relative Hodge-index sign.  It is important not to abbreviate the formula as
a difference of two absolute self-intersections: neither
\((\omega_\varphi^{n+1})\) nor \((\omega^{n+1})\) is defined for closed
model forms over a DVR.  Only their vertical-factor mixed difference is.
\end{remark}

\begin{example}[Normalization audit on a two-vertex graph]
\label{ex:ch26-normalization-audit}
Return to \Cref{ex:two-vertex-ma}: \(n=1\), \(\Gamma=[x_0,x_1]\) of length
one, \(\omega=m_0\delta_{x_0}+m_1\delta_{x_1}\), \(V=m_0+m_1\).  A model
function with vertex values \(a,b\) has
\(\ddc\varphi=(b-a)\delta_{x_0}+(a-b)\delta_{x_1}\), whence
\[
 \int\varphi\,\omega=m_0a+m_1b,
 \qquad
 \int\varphi\,\ddc\varphi=a(b-a)+b(a-b)=-(a-b)^2 .
\]
Therefore
\begin{equation}
\label{eq:ch26-two-vertex-energy}
 E_\omega(\varphi)
 =\tfrac12\Bigl(\int\varphi\,\omega+\int\varphi\,\omega_\varphi\Bigr)
 =m_0a+m_1b-\tfrac12(a-b)^2 .
\end{equation}
Three unit tests are now available.  Replacing \((a,b)\) by \((a+c,b+c)\)
adds \(c(m_0+m_1)=cV\), confirming
\eqref{eq:bfj-energy-translation}.  Taking \(a=b=0\) gives \(0\), confirming
\(E_\omega(0)=0\).  Taking \(a=-m_1\), \(b=0\) gives
\(-m_0m_1-\tfrac12m_1^2\), which is \eqref{eq:two-component-energy}.  Finally,
the sign of the quadratic term is the sign of
\eqref{eq:model-dirichlet-negativity}: it is negative, and it is what makes
\eqref{eq:ch26-two-vertex-energy} concave.
\end{example}

\subsection{The Dirichlet form and the vertical intersection form}
\label{sec:ch26-dirichlet-vertical}

The quadratic term of \eqref{eq:ch26-two-vertex-energy} is not an accident of
the two-vertex graph.  It is the value of a bilinear form on model functions
which, on a curve, \emph{is} the intersection form on vertical divisors of
an integral model, and which in Thuillier's potential theory on Berkovich
curves is the fundamental object: the non-Archimedean Dirichlet form.  We
record the identification because \eqref{eq:ch26-energy-intersection}
expresses the energy as a mixed vertical \((n+1)\)-fold intersection sum, and the case
\(n=1\) of that statement is exactly the classical relation between energy
and the vertical intersection pairing.

\begin{proposition}
\label{prop:ch26-dirichlet-vertical}
\index{Dirichlet form!non-Archimedean}
Fix a regular model \(\cX\), a relatively ample class \(\mathcal A\) on it,
and model functions \(f=\varphi_D\), \(g=\varphi_F\) determined on \(\cX\),
where \(D,F\) are vertical \(\mathbf R\)-Cartier divisors.  Put
\begin{equation}
\label{eq:ch26-dirichlet-form}
 \langle f,g\rangle_{\mathrm D}
 :=-\int_{X^{\An}}f\,\ddc g\wedge\mathcal A^{n-1}.
\end{equation}
Then:
\begin{enumerate}
\item\label{itm:ch26-dir:sym}
 \(\langle\cdot,\cdot\rangle_{\mathrm D}\) is symmetric, and
 \[
  \langle f,g\rangle_{\mathrm D}
  =-\bigl(D\cdot F\cdot\mathcal A^{n-1}\bigr)
  =\langle D,F\rangle_{\mathcal A},
 \]
 the model Dirichlet pairing \eqref{eq:ch20-dirichlet-pairing};
\item\label{itm:ch26-dir:pos}
 \(\langle f,f\rangle_{\mathrm D}\geq0\), with equality exactly when \(D\)
 is numerically proportional to \(\cX_s\);
\item\label{itm:ch26-dir:curve}
 if \(n=1\) no auxiliary class is needed and
 \(\langle f,g\rangle_{\mathrm D}=-(D\cdot F)_{\cX}\) is minus the
 intersection form on the vertical divisors of \(\cX\);
\item\label{itm:ch26-dir:energy}
 if \(n=1\) and \(f\) is \(\omega\)-psh, then with
 \(\mu_0=\MA_\omega(0)\),
 \begin{equation}
 \label{eq:ch26-energy-dirichlet}
  E_\omega(f)=\int_{X^{\An}}f\,\rd\mu_0
              -\tfrac12\langle f,f\rangle_{\mathrm D}.
 \end{equation}
\end{enumerate}
\end{proposition}

\begin{proof}
\cref{itm:ch26-dir:sym}  Both equalities are
\Cref{cor:ch20-integration-by-parts}, \eqref{eq:ch20-ibp}, applied with
\(\eta_1=\cdots=\eta_{n-1}=\mathcal A\); symmetry is the symmetry of the
middle term there, and the second equality is the definition
\eqref{eq:ch20-dirichlet-pairing}.  (For \(n\geq2\) this is also
\Cref{prop:model-integration-by-parts}, \eqref{eq:model-ibp}, which is the
same identity proved analytically.)

\cref{itm:ch26-dir:pos}  A relatively ample class is semipositive
(\Cref{lem:ch20-special-fibre-test}\,\cref{itm:ch20-sft:cone}), so
\eqref{eq:model-dirichlet-negativity} gives
\(\int f\,\ddc f\wedge\mathcal A^{n-1}\leq0\), that is
\(\langle f,f\rangle_{\mathrm D}\geq0\).  By
\cref{itm:ch26-dir:sym} the left-hand side is
\(-(D^2\cdot\mathcal A^{n-1})\), and the equality clause is that of
\Cref{prop:ch20-relative-hodge}.

\cref{itm:ch26-dir:curve}  For \(n=1\) the product
\(\mathcal A^{n-1}\) is empty and \eqref{eq:ch20-ibp} reads
\(\int f\,\ddc g=(D\cdot F)_{\cX}\).

\cref{itm:ch26-dir:energy}  With \(n=1\), \eqref{eq:model-energy} is
\[
 E_\omega(f)=\tfrac12\Bigl(\int f\,\omega+\int f\,\omega_f\Bigr)
 =\int f\,\rd\mu_0+\tfrac12\int f\,\ddc f ,
\]
because \(\omega_f=\omega+\ddc f\) and \(\int f\,\omega=\int f\,\rd\mu_0\) by
\(\mu_0=\MA_\omega(0)=\omega\).  The last integral is
\(-\langle f,f\rangle_{\mathrm D}\).
\end{proof}

\begin{remark}[Thuillier's Dirichlet form, and what \(n=1\) buys]
\label{rem:ch26-thuillier-dirichlet}
\index{Thuillier, Amaury!Dirichlet form}
On a connected proper smooth strictly \(K\)-analytic curve, Thuillier defines
the symmetric positive bilinear form
\[
 \mathrm A^0_c(X)\times \mathrm A^0_c(X)\longrightarrow\mathbf R,
 \qquad
 (\varphi,\psi)\longmapsto-\int_X\varphi\,\ddc\psi
\]
on compactly supported smooth functions, and says explicitly that it
generalizes, for a discrete valuation, the intersection form on the group of
vertical divisors on integral models of \(X\) \cite[\S3.5]{Thu05}.
\Cref{prop:ch26-dirichlet-vertical}\,\cref{itm:ch26-dir:curve} is that
statement, proved here in the book's normalizations: the sign is such that
the Dirichlet form is \emph{minus} the vertical intersection form, and its
positivity \cref{itm:ch26-dir:pos} is the relative Hodge index inequality
\eqref{eq:ch20-relative-hodge}, that is, Zariski's negativity of vertical
self-intersections.  Two consequences are worth naming.

\emph{The energy is literally quadratic in dimension one.}
Formula \eqref{eq:ch26-energy-dirichlet} exhibits \(E_\omega\) on a curve as
an affine functional minus half a positive quadratic form; this is the same
content as \eqref{eq:ch26-graph-energy}, where the quadratic form appears as
the graph Dirichlet integral \(\int_\Gamma|u'|^2\rd s\), and as
\Cref{cor:ch08-dirichlet-energy}, where it appears as
\(-\int g\,\rd\Delta_\Gamma g\).  For \(n\geq2\) the corresponding statement
fails: \eqref{eq:ch26-energy-intersection} is a polynomial of degree
\(n+1\), and no bilinear form controls it.  This is the structural reason
why \(\cE^1\) carries a genuine Hilbert-space theory in dimension one
(\Cref{sec:ch26-thuillier}) while in general it carries only the gauge of
\Cref{def:ch26-gauge}, for which the triangle inequality fails
(\Cref{ex:ch26-triangle-fails}).

\emph{Normalization.}  Thuillier's \(\ddc\) is the sum-of-outgoing-slopes
Laplacian on every skeleton and satisfies
\(\ddc\log|f|=\delta_{\operatorname{div}(f)}\); this is the convention of
\eqref{eq:graph-laplacian} and \Cref{thm:ch08-slope-formula}, so
\eqref{eq:ch26-dirichlet-form} is his form verbatim, with no rescaling.  He
presents it as the analogue of
\(\frac{\rmi}{\pi}\int_M\partial\varphi\wedge\bar\partial\psi\) on a compact
Riemann surface, which corresponds to
\(\ddc=\frac{\rmi}{\pi}\partial\bar\partial\), twice the book's complex
operator \eqref{eq:appJ-ddc}.  The factor \(2\) therefore appears only in
the archimedean half of the analogy; it is discussed once, in
\Cref{rem:ch25-thuillier-normalization}, and is never applied to a
non-Archimedean formula.
\end{remark}

\subsection{The cocycle formula}
\label{sec:ch26-cocycle}

The next proposition is the computational engine of the whole chapter.  We
first isolate the telescoping identity on which its proof rests, because the
same identity will be used again in \Cref{sec:ch26-mixed} with a different
weighting.

\begin{lemma}
\label{lem:ch26-telescope}
Let \(\varphi,\psi\) be \(\omega\)-psh model functions and put
\(u:=\psi-\varphi\), a model function.  For \(0\leq j\leq n\),
\begin{equation}
\label{eq:ch26-telescope}
 \int u\,\omega_\psi^{\,j}\wedge\omega_\varphi^{\,n-j}
 -\int u\,\omega_\psi^{\,j-1}\wedge\omega_\varphi^{\,n-j+1}
 =\int u\,\ddc u\wedge
   \omega_\psi^{\,j-1}\wedge\omega_\varphi^{\,n-j}
\end{equation}
for \(1\leq j\leq n\), and the right-hand side is \(\leq0\).
\end{lemma}

\begin{proof}
Since \(\omega_\psi-\omega_\varphi=\ddc(\psi-\varphi)=\ddc u\), the left-hand
side of \eqref{eq:ch26-telescope} equals
\[
 \int u\,(\omega_\psi-\omega_\varphi)\wedge
   \omega_\psi^{\,j-1}\wedge\omega_\varphi^{\,n-j}
 =\int u\,\ddc u\wedge
   \omega_\psi^{\,j-1}\wedge\omega_\varphi^{\,n-j},
\]
all products being taken on a common determining model, which exists by
\Cref{cor:ch20-limit-structure}.  The sign is
\eqref{eq:model-dirichlet-negativity}, that is, the relative Hodge index
inequality \eqref{eq:ch20-mixed-hodge} applied to the vertical divisor
representing \(\ddc u\) and to the \(n-1\) relatively nef classes
\(\omega_\psi^{\,j-1}\wedge\omega_\varphi^{\,n-j}\).  Relative nefness of
those classes is \Cref{prop:ch20-semipositive-well-defined}.
\end{proof}

\begin{proposition}\label{prop:energy-cocycle}
For psh model functions \(\varphi,\psi\),
\begin{equation}
\label{eq:energy-cocycle}
 E_\omega(\psi)-E_\omega(\varphi)
 =\frac1{n+1}\sum_{j=0}^n
  \int(\psi-\varphi)\,
  \omega_\varphi^j\wedge\omega_\psi^{n-j}.
\end{equation}
If \(\varphi_t=\varphi+t(\psi-\varphi)\), then
\begin{equation}
\label{eq:energy-first-variation}
 \frac{\rd}{\rd t}E_\omega(\varphi_t)
 =\int(\psi-\varphi)\,\MA_\omega(\varphi_t),
\end{equation}
and
\begin{equation}
\label{eq:energy-second-variation}
 \frac{\rd^2}{\rd t^2}E_\omega(\varphi_t)
 =n\int(\psi-\varphi)\,\ddc(\psi-\varphi)
       \wedge\omega_{\varphi_t}^{n-1}\leq0.
\end{equation}
\end{proposition}

\begin{proof}
Write \(u=\psi-\varphi\).

\emph{Step 1: the cocycle.}  By
\eqref{eq:ch26-energy-difference-intersection} and
\eqref{eq:ch20-pairing-as-integral}, on a common determining model,
\begin{equation}
\label{eq:ch26-cocycle-target}
 E_\omega(\psi)-E_\omega(\varphi)
 =\frac1{n+1}\sum_{j=0}^n\int u\,
   \omega_\psi^{\,j}\wedge\omega_\varphi^{\,n-j}.
\end{equation}
Here the first slot is the actual divisor representing \(u\), not merely
the numerical class \(\ddc u\).  Reindexing \(j\mapsto n-j\) in
\eqref{eq:ch26-cocycle-target} gives \eqref{eq:energy-cocycle}.

\emph{Step 2: the first variation.}  For \(t\) in the interval on which
\(\varphi_t\) is \(\omega\)-psh --- which contains \([0,1]\), the psh model
functions being a convex set --- apply \eqref{eq:energy-cocycle} to the pair
\(\varphi_t,\varphi_{t+h}\).  Since
\(\varphi_{t+h}-\varphi_t=hu\) and
\(\omega_{\varphi_{t+h}}=\omega_{\varphi_t}+h\,\ddc u\),
\[
 \frac{E_\omega(\varphi_{t+h})-E_\omega(\varphi_t)}{h}
 =\frac1{n+1}\sum_{j=0}^n\int u\,
   \omega_{\varphi_t}^{\,j}\wedge
   \bigl(\omega_{\varphi_t}+h\,\ddc u\bigr)^{\,n-j}.
\]
The right-hand side is a polynomial in \(h\) with value
\(\int u\,\MA_\omega(\varphi_t)\) at \(h=0\), all \(n+1\) summands then being
equal.  Letting \(h\to0\) gives \eqref{eq:energy-first-variation}.
Equivalently, \eqref{eq:ch26-energy-taylor} below exhibits
\(t\mapsto E_\omega(\varphi_t)\) as a polynomial of degree at most
\(n+1\); each coefficient is a vertical pairing with first slot \(u\).

\emph{Step 3: the second variation.}  Differentiating
\eqref{eq:energy-first-variation} once more, the only \(t\)-dependence is in
the measure, and
\(\frac{\rd}{\rd t}\omega_{\varphi_t}^{\,n}=n\,\ddc u\wedge
\omega_{\varphi_t}^{\,n-1}\).  This gives the identity in
\eqref{eq:energy-second-variation}.  Its sign is
\Cref{lem:ch26-telescope}, that is,
\eqref{eq:model-dirichlet-negativity}, applied with the \(n-1\) relatively nef
classes \(\omega_{\varphi_t}\).
\end{proof}

\begin{corollary}
\label{cor:ch26-energy-polynomial}
With the notation of \Cref{prop:energy-cocycle}, the function
\(t\mapsto E_\omega(\varphi_t)\) is a polynomial of degree at most \(n+1\) on
the largest interval \(J\) on which \(\varphi_t\) is \(\omega\)-psh, and
\begin{equation}
\label{eq:ch26-energy-taylor}
 E_\omega(\varphi_t)
 =E_\omega(\varphi)
  +\sum_{k=1}^{n+1}\frac{t^{k}}{k}\binom{n}{k-1}
   \int u\,(\ddc u)^{k-1}\wedge\omega_\varphi^{\,n-k+1},
 \qquad u=\psi-\varphi .
\end{equation}
In particular \([0,1]\subseteq J\), and \(E_\omega\) is concave on \(J\).
\end{corollary}

\begin{proof}
Expanding \eqref{eq:energy-first-variation} by multilinearity,
\[
 \frac{\rd}{\rd t}E_\omega(\varphi_t)
 =\int u\,(\omega_\varphi+t\,\ddc u)^{\,n}
 =\sum_{k=1}^{n+1}t^{\,k-1}\binom{n}{k-1}
   \int u\,(\ddc u)^{k-1}\wedge\omega_\varphi^{\,n-k+1}.
\]
Integrating from \(0\) to \(t\) gives \eqref{eq:ch26-energy-taylor}.  That
\([0,1]\subseteq J\) is convexity of \(\PSH(X,\omega)\), which is
\Cref{thm:ch21-stability}\,\cref{itm:ch21-stab:sum}.  Concavity is
\eqref{eq:energy-second-variation}.
\end{proof}

\begin{corollary}\label{cor:energy-monotone-concave}
Energy is increasing and concave on the convex set of psh model functions.
If \(\varphi\leq\psi\), then
\begin{equation}
\label{eq:energy-sandwich}
 \int(\psi-\varphi)\MA_\omega(\psi)
 \leq E_\omega(\psi)-E_\omega(\varphi)
 \leq\int(\psi-\varphi)\MA_\omega(\varphi).
\end{equation}
\end{corollary}

\begin{proof}
All mixed measures in \eqref{eq:energy-cocycle} are positive, proving
monotonicity. Concavity follows from
\eqref{eq:energy-second-variation}. The two inequalities are the supporting
line inequalities for a concave differentiable function on the segment
joining \(\varphi\) and \(\psi\).
\end{proof}

The proof just given is three lines because concavity has already been
established; it is worth spelling out the mechanism, since the same argument
recurs in \Cref{sec:ch26-mixed} and in
\Cref{chap:bfj-differentiability}.  Set
\(h(t)=E_\omega(\varphi_t)\), a concave polynomial on \([0,1]\) with
\(h(0)=E_\omega(\varphi)\) and \(h(1)=E_\omega(\psi)\).  Concavity gives
\[
 h'(1)\leq h(1)-h(0)\leq h'(0),
\]
and \eqref{eq:energy-first-variation} identifies \(h'(0)\) with
\(\int(\psi-\varphi)\MA_\omega(\varphi)\) and \(h'(1)\) with
\(\int(\psi-\varphi)\MA_\omega(\psi)\).  No hypothesis \(\varphi\leq\psi\) is
needed for \eqref{eq:energy-sandwich}: it holds for every pair of psh model
functions, and monotonicity is the special case in which the integrand is
nonnegative.

\begin{remark}[Where the ordering hypothesis is and is not needed]
\label{rem:ch26-sandwich-scope}
Only the \emph{sign} conclusion \(E_\omega(\varphi)\leq E_\omega(\psi)\)
requires \(\varphi\leq\psi\); the sandwich
\eqref{eq:energy-sandwich} is an unconditional consequence of concavity, and
it is used in that unconditional form in the proof of
\Cref{prop:bfj-euler-lagrange}.  This distinction matters because the
right-hand integral in \eqref{eq:energy-sandwich} may then be negative.
\end{remark}

\subsection{Bounded potentials}
\label{sec:ch26-bounded-energy}

Everything above was an identity between intersection numbers on one model.
We now transport it to bounded \(\omega\)-psh functions.  The transport is
not automatic: \eqref{eq:model-energy} contains a \emph{varying} psh test
function, so what is required is the weighted monotone continuity
\eqref{eq:weighted-monotone-ma}, and not merely weak convergence of the
measures.

\begin{proposition}\label{prop:bounded-energy-identities}
\Cref{def:model-energy} and
\eqref{eq:energy-cocycle}--\eqref{eq:energy-sandwich} extend to bounded
\(\omega\)-psh functions. They are continuous under decreasing bounded nets.
\end{proposition}

\begin{proof}
Let \(\varphi\) be a bounded \(\omega\)-psh function.  By
\Cref{thm:bfj-decreasing-regularization} choose a decreasing net
\(\varphi_\alpha\searrow\varphi\) of \(\omega\)-psh model functions; after
adding constants we may assume \(-M\leq\varphi\leq\varphi_\alpha\leq0\) for
some \(M\), because the net eventually lies below \(\sup\varphi+1\) and one
may replace it by \(\min\{\varphi_\alpha,\sup\varphi+1\}\), still a
decreasing net of model functions by
\Cref{prop:ch20-model-max} and \Cref{prop:ch21-constants}.

\emph{Step 1: existence of the limit.}  Each summand of
\eqref{eq:model-energy} for \(\varphi_\alpha\) is of the form
\(\int h^\alpha\MA_\omega(\varphi_1^\alpha,\ldots,\varphi_n^\alpha)\) with
\(h^\alpha=\varphi_\alpha\searrow\varphi\) and every
\(\varphi_i^\alpha\in\{\varphi_\alpha,0\}\).  By
\eqref{eq:weighted-monotone-ma} each summand converges, so
\(E_\omega(\varphi_\alpha)\) converges; we define \(E_\omega(\varphi)\) to be
the limit and observe that it is given by \eqref{eq:model-energy}, the
integrals now being taken against the bounded mixed products of
\Cref{thm:bounded-ma-extension}.

\emph{Step 2: independence of the net.}  Let \(\varphi_\beta'\searrow\varphi\)
be a second decreasing net of \(\omega\)-psh model functions.  Fix
\(\alpha\) and \(\varepsilon>0\).  The sets
\[
 F_\beta:=\{\varphi_\beta'\geq\varphi_\alpha+\varepsilon\}\subset X^{\An}
\]
are closed, because \(\varphi_\beta'\) is upper semicontinuous and
\(\varphi_\alpha\) is continuous, and they decrease with \(\beta\); their
intersection is contained in
\(\{\varphi\geq\varphi_\alpha+\varepsilon\}=\varnothing\).  Compactness of
\(X^{\An}\) therefore makes one of them empty, that is,
\(\varphi_\beta'\leq\varphi_\alpha+\varepsilon\) for some \(\beta\).
Monotonicity of \(E_\omega\) at model level
(\Cref{cor:energy-monotone-concave}) and
\eqref{eq:bfj-energy-translation} give
\(\lim_{\beta'}E_\omega(\varphi_{\beta'}')\leq
E_\omega(\varphi_\alpha)+\varepsilon V\); taking the limit in \(\alpha\) and
then \(\varepsilon\downarrow0\), and repeating with the two nets exchanged,
shows that the two limits agree.

\emph{Step 3: the identities.}  Fix bounded \(\varphi,\psi\) and decreasing
model nets \(\varphi_\alpha\searrow\varphi\),
\(\psi_\beta\searrow\psi\), indexed by the product of the two directed sets.
Every integral occurring in
\eqref{eq:energy-cocycle}--\eqref{eq:energy-sandwich} is again of the form
\(\int h\,\MA_\omega(\chi_1,\ldots,\chi_n)\) with \(h\) a difference of two
members of decreasing bounded nets and each \(\chi_i\) a member of such a
net.  Splitting \(h=\psi_\beta-\varphi_\alpha\) into its two terms and
applying \eqref{eq:weighted-monotone-ma} to each, every such integral
converges to the corresponding integral for \(\varphi,\psi\).  Since the
identities hold along the net, they hold in the limit.  For
\eqref{eq:energy-first-variation} and
\eqref{eq:energy-second-variation}, note that
\(\varphi_{t,\alpha\beta}=\varphi_\alpha+t(\psi_\beta-\varphi_\alpha)\)
decreases to \(\varphi_t\) for each fixed \(t\in[0,1]\), and that
\eqref{eq:ch26-energy-taylor} exhibits both sides as polynomials in \(t\) of
degree at most \(n+1\) whose coefficients converge; convergence of the
coefficients of a polynomial of bounded degree implies convergence of the
derivatives.

\emph{Step 4: decreasing bounded nets.}  If \(\varphi_\gamma\searrow\varphi\)
is a decreasing net of \emph{bounded} \(\omega\)-psh functions with bounded
limit, then choosing model nets below each \(\varphi_\gamma\) and passing to
the diagonal directed set produces a decreasing model net converging to
\(\varphi\); by Step~2 the value of \(E_\omega(\varphi)\) computed from it is
the correct one, and monotonicity gives
\(E_\omega(\varphi_\gamma)\searrow E_\omega(\varphi)\).
\end{proof}

\begin{proposition}
\label{prop:ch26-energy-lipschitz}
For bounded \(\omega\)-psh \(\varphi,\psi\),
\begin{equation}
\label{eq:ch26-energy-lipschitz}
 \bigl|E_\omega(\psi)-E_\omega(\varphi)\bigr|
 \leq V\,\lVert\psi-\varphi\rVert_\infty .
\end{equation}
In particular \(E_\omega\) is continuous for uniform convergence of bounded
\(\omega\)-psh functions, and \(E_\omega\) is finite on that locus.
\end{proposition}

\begin{proof}
Put \(c=\lVert\psi-\varphi\rVert_\infty\), so \(\varphi-c\leq\psi\leq\varphi+c\).
Monotonicity and \eqref{eq:bfj-energy-translation} give
\[
 E_\omega(\varphi)-cV=E_\omega(\varphi-c)
 \leq E_\omega(\psi)\leq E_\omega(\varphi+c)=E_\omega(\varphi)+cV .
\]
Finiteness follows by comparing \(\varphi\) with the constants
\(\inf\varphi\) and \(\sup\varphi\).
\end{proof}

\begin{warning}{Which product each identity concerns}
\Cref{def:model-energy} and \Cref{prop:energy-cocycle} are statements about
the \emph{model} product \eqref{eq:model-mixed-measure}, and their proofs are
intersection theory on one model.  \Cref{prop:bounded-energy-identities} and
\Cref{prop:ch26-energy-lipschitz} are statements about the \emph{bounded}
product of \Cref{thm:bounded-ma-extension}.  Nothing has yet been asserted
about unbounded potentials: for those, both \(E_\omega\) and
\(\MA_\omega\) still have to be defined, in
\Cref{sec:ch26-extension,sec:ch26-nonpluripolar}, and the identities have to
be reproved.
\end{warning}

\section{Mixed energies, cocycles, and integration by parts}
\label{sec:ch26-mixed}

\subsection{The functionals \texorpdfstring{\(I\)}{I} and
\texorpdfstring{\(J\)}{J}}
\label{sec:ch26-ij-def}

The cocycle formula \eqref{eq:energy-cocycle} compares two potentials through
\(n+1\) mixed integrals.  Two combinations of those integrals occur so often
that they deserve names.  They are the non-Archimedean analogues of Aubin's
functionals, and all the inequalities we shall need between them are exact
identities, not estimates.

\begin{definition}[Mixed energies]
\label{def:ch26-ij}
\index{energy functionals!I and J@\(I\) and \(J\)}
For bounded \(\omega\)-psh functions \(\varphi,\psi\) put
\begin{align}
\label{eq:ch26-I-definition}
 I(\varphi,\psi)
 &:=\int(\varphi-\psi)\bigl(\MA_\omega(\psi)-\MA_\omega(\varphi)\bigr),\\
\label{eq:ch26-J-definition}
 J_\psi(\varphi)
 &:=\int(\varphi-\psi)\,\MA_\omega(\psi)
    -\bigl(E_\omega(\varphi)-E_\omega(\psi)\bigr).
\end{align}
\end{definition}

Both are finite on the bounded locus by
\Cref{prop:ch26-energy-lipschitz}.  \(I\) is symmetric in its two arguments
by inspection; \(J_\psi(\varphi)\) is not, and the asymmetry is exactly what
\Cref{thm:ch26-mixed-energy-identities} computes.  Both are invariant under
adding a common constant to \(\varphi\) and \(\psi\), and \(I\) is invariant
under adding \emph{separate} constants, because the two Monge--Amp\`ere
measures have equal mass \(V\); \(J_\psi(\varphi+c)=J_\psi(\varphi)\) as
well, since both terms increase by \(cV\).

\subsection{The integration-by-parts chain}
\label{sec:ch26-ibp-chain}

\begin{lemma}
\label{lem:ch26-dirichlet-coefficients}
Let \(\varphi,\psi\) be bounded \(\omega\)-psh functions and set
\(u:=\varphi-\psi\).  For \(0\leq k\leq n-1\) put
\begin{equation}
\label{eq:ch26-ak}
 a_k(\varphi,\psi)
 :=-\int u\,\ddc u\wedge
     \omega_\varphi^{\,k}\wedge\omega_\psi^{\,n-1-k}.
\end{equation}
Then \(a_k(\varphi,\psi)\geq0\), the family is symmetric in the sense that
\(a_k(\varphi,\psi)=a_{n-1-k}(\psi,\varphi)\), and
\begin{equation}
\label{eq:ch26-ak-telescope}
 \int u\,\MA_\omega(\psi)-\int u\,
   \omega_\varphi^{\,j}\wedge\omega_\psi^{\,n-j}
 =\sum_{k=0}^{j-1}a_k(\varphi,\psi)
 \qquad(0\leq j\leq n).
\end{equation}
\end{lemma}

\begin{proof}
Nonnegativity is \eqref{eq:model-dirichlet-negativity} for model functions
and its bounded extension \Cref{cor:bounded-ma-multilinearity}; the classes
\(\omega_\varphi,\omega_\psi\) are semipositive, so
\eqref{eq:ch20-mixed-hodge} applies.  The symmetry is immediate from
\(u\mapsto-u\) and \(\ddc(-u)=-\ddc u\), which leave the product
\(u\,\ddc u\) unchanged, together with exchanging the roles of the two
families of factors.

For \eqref{eq:ch26-ak-telescope}, apply \Cref{lem:ch26-telescope} with the
roles of \(\varphi\) and \(\psi\) as displayed there --- in the present
notation, with \(\psi-\varphi\) replaced by \(u=\varphi-\psi\) --- to obtain,
for \(1\leq k+1\leq j\),
\[
 \int u\,\omega_\varphi^{\,k}\wedge\omega_\psi^{\,n-k}
 -\int u\,\omega_\varphi^{\,k+1}\wedge\omega_\psi^{\,n-k-1}
 =-\int u\,\ddc u\wedge\omega_\varphi^{\,k}\wedge\omega_\psi^{\,n-1-k}
 =a_k(\varphi,\psi).
\]
Summing over \(0\leq k\leq j-1\) telescopes to
\eqref{eq:ch26-ak-telescope}, the \(k=0\) left endpoint being
\(\int u\,\omega_\psi^{\,n}=\int u\,\MA_\omega(\psi)\).  Every integration by
parts used here is \eqref{eq:model-ibp} at model level, extended to bounded
psh differences by \Cref{cor:bounded-ma-multilinearity}; the extension is
legitimate because \(u\) is a difference of two bounded psh functions and all
remaining factors are bounded psh.
\end{proof}

\begin{theorem}
\label{thm:ch26-mixed-energy-identities}
\index{energy functionals!identities}
Let \(\varphi,\psi\) be bounded \(\omega\)-psh functions and let
\(a_k=a_k(\varphi,\psi)\) be as in \eqref{eq:ch26-ak}.  Then, for the
\emph{bounded} Monge--Amp\`ere product:
\begin{enumerate}
\item\label{itm:ch26-ij:I}
 \(\displaystyle I(\varphi,\psi)=\sum_{k=0}^{n-1}a_k\geq0\);
\item\label{itm:ch26-ij:J}
 \(\displaystyle J_\psi(\varphi)=\frac1{n+1}\sum_{k=0}^{n-1}(n-k)\,a_k\)
 and
 \(\displaystyle J_\varphi(\psi)=\frac1{n+1}\sum_{k=0}^{n-1}(k+1)\,a_k\);
\item\label{itm:ch26-ij:sum}
 \(J_\psi(\varphi)+J_\varphi(\psi)=I(\varphi,\psi)\);
\item\label{itm:ch26-ij:comparison}
 \(\displaystyle
 \frac1{n+1}I(\varphi,\psi)\ \leq\ J_\psi(\varphi)\ \leq\
 \frac{n}{n+1}I(\varphi,\psi)\),
 and the same bounds hold for \(J_\varphi(\psi)\);
\item\label{itm:ch26-ij:positive}
 \(J_\psi(\varphi)\geq0\), with equality if and only if \(a_k=0\) for every
 \(k\), if and only if \(I(\varphi,\psi)=0\).
\end{enumerate}
\end{theorem}

\begin{proof}
\emph{\cref{itm:ch26-ij:I}.}  Take \(j=n\) in
\eqref{eq:ch26-ak-telescope}:
\[
 \int u\,\MA_\omega(\psi)-\int u\,\MA_\omega(\varphi)
 =\sum_{k=0}^{n-1}a_k .
\]
The left-hand side is \(I(\varphi,\psi)\) by
\eqref{eq:ch26-I-definition}.

\emph{\cref{itm:ch26-ij:J}.}  By \eqref{eq:energy-cocycle} in the bounded form
supplied by \Cref{prop:bounded-energy-identities},
\[
 E_\omega(\varphi)-E_\omega(\psi)
 =\frac1{n+1}\sum_{j=0}^n\int u\,
   \omega_\varphi^{\,j}\wedge\omega_\psi^{\,n-j},
\]
so that
\[
 J_\psi(\varphi)
 =\frac1{n+1}\sum_{j=0}^n
   \Bigl(\int u\,\MA_\omega(\psi)
        -\int u\,\omega_\varphi^{\,j}\wedge\omega_\psi^{\,n-j}\Bigr)
 =\frac1{n+1}\sum_{j=0}^n\sum_{k=0}^{j-1}a_k ,
\]
using \eqref{eq:ch26-ak-telescope}.  Interchanging the order of summation,
the coefficient of \(a_k\) is the number of indices \(j\) with
\(k+1\leq j\leq n\), namely \(n-k\).  This is the first formula.  The second
follows by exchanging \(\varphi\) and \(\psi\), which replaces \(u\) by
\(-u\) and \(a_k\) by \(a_{n-1-k}\), and then reindexing.

\emph{\cref{itm:ch26-ij:sum}.}  Add the two formulas in
\cref{itm:ch26-ij:J}: the coefficient of \(a_k\) is
\(\bigl((n-k)+(k+1)\bigr)/(n+1)=1\).

\emph{\cref{itm:ch26-ij:comparison}.}  For \(0\leq k\leq n-1\) one has
\(1\leq n-k\leq n\) and \(1\leq k+1\leq n\).  Since all \(a_k\geq0\), the
bounds follow by comparing \cref{itm:ch26-ij:J} with
\cref{itm:ch26-ij:I} term by term.

\emph{\cref{itm:ch26-ij:positive}.}  Nonnegativity and the equality cases are
immediate from \cref{itm:ch26-ij:J}, \cref{itm:ch26-ij:I} and \(a_k\geq0\).
\end{proof}

\begin{corollary}
\label{cor:ch26-cauchy-schwarz}
Let \(\varphi,\psi\) be bounded \(\omega\)-psh, let \(u=\varphi-\psi\), and
let \(h\) be any difference of two bounded \(\omega\)-psh functions.  Then,
for \(0\leq k\leq n-1\),
\begin{equation}
\label{eq:ch26-cauchy-schwarz}
 \Bigl|\int h\,\ddc u\wedge
    \omega_\varphi^{\,k}\wedge\omega_\psi^{\,n-1-k}\Bigr|^2
 \leq
 a_k(\varphi,\psi)\cdot
 \Bigl(-\int h\,\ddc h\wedge
    \omega_\varphi^{\,k}\wedge\omega_\psi^{\,n-1-k}\Bigr).
\end{equation}
\end{corollary}

\begin{proof}
At model level this is \eqref{eq:ch20-mixed-cauchy-schwarz} applied to the
vertical divisors representing \(\ddc u\) and \(\ddc h\) and to the \(n-1\)
relatively nef classes.  Both sides are limits of the corresponding model
expressions along simultaneous decreasing regularizations, by
\Cref{prop:bounded-energy-identities} and
\Cref{cor:bounded-ma-multilinearity}.
\end{proof}

\subsection{Cocycle formulae}
\label{sec:ch26-cocycle-formulae}

The name \emph{cocycle} for \eqref{eq:energy-cocycle} refers to the following
elementary but constantly used consistency property.  Write
\[
 \mathcal E(\varphi,\psi):=E_\omega(\varphi)-E_\omega(\psi)
\]
for bounded \(\omega\)-psh \(\varphi,\psi\).

\begin{proposition}
\label{prop:ch26-cocycle-chain}
For bounded \(\omega\)-psh \(\varphi,\psi,\chi\),
\begin{equation}
\label{eq:ch26-cocycle-chain}
 \mathcal E(\varphi,\psi)+\mathcal E(\psi,\chi)
 =\mathcal E(\varphi,\chi),
 \qquad
 \mathcal E(\varphi,\psi)=-\mathcal E(\psi,\varphi),
\end{equation}
and, spelled out through \eqref{eq:energy-cocycle},
\begin{equation}
\label{eq:ch26-cocycle-explicit}
\begin{aligned}
 \frac1{n+1}\sum_{j=0}^n
  \int(\varphi-\psi)\,\omega_\varphi^{\,j}\wedge\omega_\psi^{\,n-j}
 &+\frac1{n+1}\sum_{j=0}^n
  \int(\psi-\chi)\,\omega_\psi^{\,j}\wedge\omega_\chi^{\,n-j}\\
 &=\frac1{n+1}\sum_{j=0}^n
  \int(\varphi-\chi)\,\omega_\varphi^{\,j}\wedge\omega_\chi^{\,n-j}.
\end{aligned}
\end{equation}
Moreover, for \(0\leq j\leq n\) the individual mixed integrals obey
\begin{equation}
\label{eq:ch26-mixed-symmetry}
 \int(\varphi-\psi)\,\omega_\varphi^{\,j}\wedge\omega_\psi^{\,n-j}
 =\int(\varphi-\psi)\,\omega_\varphi^{\,n-j}\wedge\omega_\psi^{\,j}
 +\Bigl(\sum_{k=j}^{n-j-1}-\sum_{k=n-j}^{j-1}\Bigr)a_k(\varphi,\psi),
\end{equation}
where an empty sum is zero, so that exactly one of the two sums contributes.
\end{proposition}

\begin{proof}
The identities \eqref{eq:ch26-cocycle-chain} are trivial for the difference
of values of a single function; their content is
\eqref{eq:ch26-cocycle-explicit}, which is what one obtains by substituting
\eqref{eq:energy-cocycle} three times.  This substitution is legitimate
because each of the three energy differences is a finite real number, by
\Cref{prop:ch26-energy-lipschitz}, so no cancellation of infinite quantities
occurs; compare entry~\cref{itm:ch26-ledger:diff} of
\Cref{sec:ch26-extended-real}.  Identity \eqref{eq:ch26-mixed-symmetry} is
\eqref{eq:ch26-ak-telescope} applied twice and subtracted, once with \(j\) and
once with \(n-j\).
\end{proof}

\begin{remark}[Why the cocycle property is not automatic later]
\label{rem:ch26-cocycle-caution}
For unbounded potentials, \eqref{eq:ch26-cocycle-chain} is a statement about
extended reals and may fail to have content: if
\(E_\omega(\psi)=-\infty\), then \(\mathcal E(\varphi,\psi)\) and
\(\mathcal E(\psi,\chi)\) are \(+\infty\) and \(-\infty\), and their sum is
undefined by the ledger.  This is the first place where finiteness of energy
is used as a hypothesis rather than as a conclusion, and it is the reason
\Cref{chap:bfj-differentiability} normalizes before subtracting.  The
explicit form \eqref{eq:ch26-cocycle-explicit} remains meaningful on
\(\cE^1\) once the mixed products have been extended in
\Cref{sec:ch26-mixed-products}, because each individual integral is then
finite; this is proved in \Cref{prop:ch26-e1-cocycle}.
\end{remark}

\begin{proposition}
\label{prop:ch26-reference-change}
Let \(\psi_0\) be an \(\omega\)-psh model function and put
\(\omega':=\omega_{\psi_0}\).  Then \(\varphi\mapsto\varphi-\psi_0\) is a
bijection from bounded \(\omega\)-psh functions to bounded
\(\omega'\)-psh functions, and
\begin{equation}
\label{eq:ch26-reference-change}
 E_{\omega'}(\varphi-\psi_0)
 =E_\omega(\varphi)-E_\omega(\psi_0).
\end{equation}
In particular energy differences do not depend on the choice of reference
form within a fixed generic numerical class.
\end{proposition}

\begin{proof}
The bijection is \Cref{cor:ch21-reference-independence}, and
\(\omega'_{\varphi-\psi_0}=\omega_\varphi\).  Both sides of
\eqref{eq:ch26-reference-change} vanish at \(\varphi=\psi_0\).  By
\eqref{eq:ch26-energy-difference-intersection}, for model \(\varphi\) both
sides equal the legitimate vertical sum
\[
 \frac1{n+1}\sum_{j=0}^n
 \bigl(D_{\varphi-\psi_0}\cdot\omega_\varphi^{\,j}
                  \cdot\omega_{\psi_0}^{\,n-j}\bigr).
\]
The general case follows by decreasing
regularization and \Cref{prop:bounded-energy-identities}.
\end{proof}

\subsection{Where integration by parts is legitimate}
\label{sec:ch26-ibp-scope}

Every identity of this section was obtained by moving a \(\ddc\) from one
factor to another.  Because that step is the one most often performed without
justification, we record precisely what has been used.

\begin{proposition}
\label{prop:ch26-ibp-regimes}
Let \(\eta_1,\ldots,\eta_{n-1}\) be semipositive closed model forms and let
\(T=\eta_1\wedge\cdots\wedge\eta_{n-1}\).
\begin{enumerate}
\item\label{itm:ch26-ibp:model}
 \emph{(Model.)}  For model functions \(f,g\),
 \(\int f\,\ddc g\wedge T=\int g\,\ddc f\wedge T\).  This is
 \eqref{eq:model-ibp}, equivalently \eqref{eq:ch20-ibp}; both sides are the
 vertical intersection number \((D\cdot D'\cdot\eta_1\cdots\eta_{n-1})\).
\item\label{itm:ch26-ibp:bounded}
 \emph{(Bounded.)}  The same identity holds when \(f,g\) are differences of
 bounded \(\omega\)-psh functions and each \(\eta_i\) is replaced by
 \(\omega_{\chi_i}\) for bounded \(\omega\)-psh \(\chi_i\).  This is
 \Cref{cor:bounded-ma-multilinearity}, and its proof is simultaneous
 decreasing regularization together with
 \eqref{eq:weighted-monotone-ma}.
\item\label{itm:ch26-ibp:finite-energy}
 \emph{(Finite energy.)}  For \(f,g\) differences of members of
 \(\cE^1(X,\omega)\), no unrestricted integration-by-parts formula is used in
 this book.  What is used instead is the truncated identity: for every
 \(t>0\) the bounded case \cref{itm:ch26-ibp:bounded} applies to the
 truncations, and the passage \(t\to\infty\) is controlled by the mass tail
 \eqref{eq:ch26-mixed-tail} and by \Cref{lem:integration-ledger}.
\end{enumerate}
\end{proposition}

\begin{proof}
Assertions \cref{itm:ch26-ibp:model} and \cref{itm:ch26-ibp:bounded} are the
quoted results.  Assertion \cref{itm:ch26-ibp:finite-energy} is a statement
about the architecture of the proofs below, and is verified by inspection:
each of \Cref{thm:e1-mixed-products},
\Cref{prop:ch26-e1-cocycle} and
\Cref{thm:ch26-e1-maximum-principle} is proved by truncation, and no
integration by parts is performed at an unbounded potential.
\end{proof}

\begin{warning}{Integration by parts is not available at arbitrary
singularity}
The classical failure is not subtle.  If \(f=g\) is an unbounded psh
function, then \(\int f\,\ddc f\wedge T\) may be \(-\infty\) while every
truncated version is finite, and no cancellation between two infinite
quantities is permitted
(\Cref{rem:forbidden-cancellations}).  \Cref{ex:ch26-infinite-energy}
exhibits a potential on \(\mathbf P^1\) for which this happens, with the
divergence rate computed.  This is why finiteness of energy, and not merely
membership in \(\PSH(X,\omega)\), is the hypothesis of every statement in
\Cref{sec:ch26-mixed-products}.
\end{warning}

\section{Extension to arbitrary psh functions}
\label{sec:ch26-extension}

A bounded psh function is approximated from above by model functions, and
\Cref{prop:bounded-energy-identities} shows that the energy of the
approximants converges.  For an unbounded psh function there are still model
majorants --- \Cref{thm:bfj-decreasing-regularization} provides a whole
decreasing net of them --- but their energies need not converge to a real
number.  The following definition takes the only value that is consistent
with monotonicity.

\begin{definition}[Extended energy]
\label{def:extended-energy}
For arbitrary \(\varphi\in\PSH(X,\omega)\), define
\begin{equation}
\label{eq:extended-energy}
 E_\omega(\varphi):=
 \inf\{E_\omega(\psi)\mid
       \psi\text{ is a psh model function and }\psi\geq\varphi\}
 \in[-\infty,+\infty[.
\end{equation}
\end{definition}

The infimum is over a nonempty set, by
\Cref{thm:bfj-decreasing-regularization}, so it is not \(+\infty\) by the
empty-infimum convention of the ledger; and it is bounded above by
\(E_\omega(\psi)\) for any single model majorant, hence lies in
\([-\infty,+\infty)\).  Two consistency checks are needed before the notation
can be reused.

\begin{lemma}
\label{lem:ch26-consistency}
If \(\varphi\) is a bounded \(\omega\)-psh function, then
\eqref{eq:extended-energy} agrees with the value supplied by
\Cref{prop:bounded-energy-identities}.  If \(\varphi\) is a model function,
both agree with \Cref{def:model-energy}.
\end{lemma}

\begin{proof}
Write \(E^{\mathrm{b}}_\omega\) for the bounded functional of
\Cref{prop:bounded-energy-identities} and \(E_\omega\) for
\eqref{eq:extended-energy}.  Every model majorant \(\psi\geq\varphi\)
satisfies \(E^{\mathrm{b}}_\omega(\varphi)\leq E_\omega^{\mathrm b}(\psi)\)
by monotonicity of the bounded functional, so
\(E^{\mathrm b}_\omega(\varphi)\leq E_\omega(\varphi)\).  Conversely, a
decreasing model net \(\varphi_\alpha\searrow\varphi\) consists of model
majorants and has \(E_\omega(\varphi_\alpha)\to
E^{\mathrm b}_\omega(\varphi)\) by Step~1 of the proof of
\Cref{prop:bounded-energy-identities}, so
\(E_\omega(\varphi)\leq E^{\mathrm b}_\omega(\varphi)\).  For a model
function \(\varphi\), \(\psi=\varphi\) is an admissible majorant and the
infimum is attained.
\end{proof}

\begin{theorem}\label{thm:extended-energy-basic}
The functional \(E_\omega\) is increasing, concave, upper semicontinuous in
the BFJ weak topology, and continuous along decreasing nets. It obeys
\eqref{eq:bfj-energy-translation}.
\end{theorem}

\begin{proof}
\emph{Step 1: monotonicity and translation.}  If \(\varphi\leq\varphi'\), then
every model majorant of \(\varphi'\) is one of \(\varphi\), so the infimum
defining \(E_\omega(\varphi)\) is taken over a larger set and is therefore
smaller.  If \(c\in\mathbf R\), then \(\psi\mapsto\psi+c\) is a bijection
between the model majorants of \(\varphi\) and those of \(\varphi+c\), and it
changes energy by \(cV\) by \eqref{eq:bfj-energy-translation} at model level;
\Cref{lem:extended-real-safe}\,(1) then transports the identity through the
infimum, including when the value is \(-\infty\).

\emph{Step 2: concavity.}  Let \(\varphi_0,\varphi_1\in\PSH(X,\omega)\) and
\(s\in[0,1]\).  If either energy is \(-\infty\) the asserted inequality
\(E_\omega((1-s)\varphi_0+s\varphi_1)\geq(1-s)E_\omega(\varphi_0)+sE_\omega(\varphi_1)\)
is vacuous, the right-hand side being \(-\infty\) --- note that no
\((+\infty)+(-\infty)\) occurs, since \(E_\omega<+\infty\) always.  Assume
therefore both are finite.  Given model majorants \(\psi_i\geq\varphi_i\),
the model function \((1-s)\psi_0+s\psi_1\) is a majorant of
\((1-s)\varphi_0+s\varphi_1\) and is \(\omega\)-psh by
\Cref{thm:ch21-stability}\,\cref{itm:ch21-stab:sum}; hence
\[
 E_\omega\bigl((1-s)\varphi_0+s\varphi_1\bigr)
 \leq E_\omega\bigl((1-s)\psi_0+s\psi_1\bigr).
\]
That is the wrong direction, so we argue instead through decreasing nets.
Choose decreasing model nets \(\psi_i^\alpha\searrow\varphi_i\).  Then
\((1-s)\psi_0^\alpha+s\psi_1^\alpha\searrow(1-s)\varphi_0+s\varphi_1\), and
by Step~4 below --- whose proof uses only Steps~1 and~3 --- the energies
converge to \(E_\omega((1-s)\varphi_0+s\varphi_1)\).  Model concavity
(\Cref{cor:energy-monotone-concave}) gives
\[
 E_\omega\bigl((1-s)\psi_0^\alpha+s\psi_1^\alpha\bigr)
 \geq(1-s)E_\omega(\psi_0^\alpha)+sE_\omega(\psi_1^\alpha),
\]
and letting \(\alpha\) run gives the claim, the right-hand side converging by
Step~4 applied to each factor separately.

\emph{Step 3: upper semicontinuity.}  Fix
\(\varphi_0\in\PSH(X,\omega)\) and \(a>E_\omega(\varphi_0)\).  By
\eqref{eq:extended-energy} there are a model majorant
\(\psi_0\geq\varphi_0\) and \(\delta>0\) with
\(E_\omega(\psi_0)<a-\delta\).  Choose \(\eta>0\) with \(\eta V<\delta\) and
put
\[
 U:=\Bigl\{\varphi\in\PSH(X,\omega)\ \Bigm|\
       \sup_{X^{\An}}(\varphi-\psi_0)<\eta\Bigr\}.
\]
The functional \(\varphi\mapsto\sup_{X^{\An}}(\varphi-\psi_0)\) is
\emph{continuous} for the BFJ topology.  Indeed, put
\(\omega':=\omega_{\psi_0}\), a semipositive closed model form with the same
generic class; by \Cref{cor:ch21-reference-independence} the assignment
\(\varphi\mapsto\varphi-\psi_0\) is a bijection
\(\PSH(X,\omega)\to\PSH(X,\omega')\), and by
\Cref{prop:ch21-sup-on-skeleton} applied to \(\omega'\),
\[
 \sup_{X^{\An}}(\varphi-\psi_0)
 =\max_{\Delta_{\mathcal X}}(\varphi-\psi_0)
\]
for any SNC model \(\mathcal X\) determining \(\omega'\).  The BFJ topology
is the topology of uniform convergence on every dual complex
(\Cref{thm:bfj-skeletal-equicontinuity},
\Cref{rem:ch22-topology-comparison}), and the maximum over the compact set
\(\Delta_{\mathcal X}\) is continuous for uniform convergence there.
Therefore \(U\) is open, and it contains \(\varphi_0\) because
\(\varphi_0\leq\psi_0\).  For \(\varphi\in U\) we have
\(\varphi\leq\psi_0+\eta\), whence by Step~1
\[
 E_\omega(\varphi)\leq E_\omega(\psi_0+\eta)
 =E_\omega(\psi_0)+\eta V<a .
\]
Thus \(\{E_\omega<a\}\) is open for every \(a\), which is upper
semicontinuity.

\emph{Step 4: decreasing nets.}  Let \(\varphi_\alpha\searrow\varphi\) be a
decreasing net in \(\PSH(X,\omega)\); the limit is \(\omega\)-psh by
\Cref{thm:ch21-decreasing-stability}, provided it is not identically
\(-\infty\), which we assume.  Monotonicity gives
\(\inf_\alpha E_\omega(\varphi_\alpha)\geq E_\omega(\varphi)\), the net of
energies being decreasing and hence convergent in \(\Rbar\) by
\Cref{lem:monotone-net-converges}.  For the converse, fix a model majorant
\(\psi\geq\varphi\) and \(\eta>0\).  The sets
\[
 F_\alpha:=\{\varphi_\alpha\geq\psi+\eta\}\subset X^{\An}
\]
are closed --- \(\varphi_\alpha\) is upper semicontinuous and \(\psi\) is
continuous --- and decrease, with
\(\bigcap_\alpha F_\alpha\subseteq\{\varphi\geq\psi+\eta\}=\varnothing\).
Compactness of \(X^{\An}\) makes some \(F_{\alpha_0}\) empty, so
\(\varphi_{\alpha_0}\leq\psi+\eta\) and
\[
 \inf_\alpha E_\omega(\varphi_\alpha)
 \leq E_\omega(\varphi_{\alpha_0})
 \leq E_\omega(\psi)+\eta V.
\]
Letting \(\eta\downarrow0\) and taking the infimum over model majorants
\(\psi\) gives \(\inf_\alpha E_\omega(\varphi_\alpha)\leq E_\omega(\varphi)\),
which completes Step~4.

The compactness argument in Step~4 is the one underlying
\cite[Proposition~6.2]{BFJ15}.  It is specific to decreasing nets of upper
semicontinuous functions on the compact Berkovich space, not an unrestricted
order-functional principle: it uses that the majorant \(\psi\) is
\emph{continuous}, which is where the choice of model functions as
comparison objects enters.
\end{proof}

\begin{proposition}
\label{prop:ch26-truncation-limit}
Let \(\varphi\in\PSH(X,\omega)\) with \(\varphi\leq0\).  Then
\begin{equation}
\label{eq:ch26-truncation-limit}
 E_\omega(\varphi)
 =\lim_{t\to\infty}E_\omega(\varphi^{\langle t\rangle})
 =\inf_{t>0}E_\omega(\varphi^{\langle t\rangle})
 \in[-\infty,0].
\end{equation}
Consequently \(E_\omega(\varphi)>-\infty\) if and only if the truncation
energies are bounded below, and in that case the convergence in
\eqref{eq:ch26-truncation-limit} is monotone.
\end{proposition}

\begin{proof}
The truncations \(\varphi^{\langle t\rangle}\) are bounded \(\omega\)-psh
functions decreasing to \(\varphi\) as \(t\to\infty\); the family is
decreasing because \(\max\{\varphi,-t\}\geq\max\{\varphi,-s\}\) for
\(t\leq s\), and its pointwise limit is \(\varphi\) by definition.  Apply
Step~4 of \Cref{thm:extended-energy-basic} to the net indexed by
\(t\in(0,\infty)\).  The two-sided bound \(E_\omega(\varphi)\leq0\) follows
from \(\varphi\leq0\) and \(E_\omega(0)=0\).
\end{proof}

The extension is concave and monotone, and it is therefore natural to ask for
its supergradients.  The next proposition answers that question completely on
the bounded locus, and it is the exact sense in which the Monge--Amp\`ere
operator is the derivative of the energy at a bounded potential.  It is also
the point at which the abstract pattern of \Cref{chap:topological-tools} ---
a concave functional \(E\) satisfying \eqref{eq:energy-translation} together
with its dual \eqref{eq:dual-energy-abstract} --- acquires its
non-Archimedean content; compare \Cref{subsec:legendre-finite} and
\Cref{prop:concave-biconjugation}.  The dual functional itself is constructed
in \Cref{chap:bfj-finite-energy-measures}.

\begin{proposition}
\label{prop:ch26-supergradient}
Let \(\varphi\) be a bounded \(\omega\)-psh function and let
\(\psi\in\PSH(X,\omega)\) be arbitrary.  Then
\(\int(\psi-\varphi)\MA_\omega(\varphi)\) is well defined in
\([-\infty,+\infty)\), and
\begin{equation}
\label{eq:ch26-supergradient}
 E_\omega(\psi)\ \leq\ E_\omega(\varphi)
   +\int(\psi-\varphi)\,\MA_\omega(\varphi).
\end{equation}
Moreover
\begin{equation}
\label{eq:ch26-energy-as-infimum}
 E_\omega(\psi)
 =\inf\Bigl\{E_\omega(\varphi)+\int(\psi-\varphi)\MA_\omega(\varphi)
   \ \Bigm|\ \varphi\ \text{bounded }\omega\text{-psh}\Bigr\},
\end{equation}
so that \(E_\omega\) is the infimum of the affine functionals defined by the
bounded Monge--Amp\`ere measures.
\end{proposition}

\begin{proof}
Since \(\psi\) is bounded above and \(\varphi\) is bounded, the function
\(\psi-\varphi\) is bounded above, so
\Cref{lem:integration-ledger}\,(1) makes the integral well defined in
\([-\infty,+\infty)\); no forbidden difference occurs, \(\varphi\) being
finite everywhere.

For bounded \(\psi\), \eqref{eq:ch26-supergradient} is the right-hand
inequality of \eqref{eq:energy-sandwich} in its unconditional form
(\Cref{rem:ch26-sandwich-scope}).  For arbitrary \(\psi\), normalize
\(\psi\leq0\) and apply the bounded case to
\(\psi^{\langle t\rangle}\):
\[
 E_\omega(\psi^{\langle t\rangle})
 \leq E_\omega(\varphi)
   +\int(\psi^{\langle t\rangle}-\varphi)\MA_\omega(\varphi).
\]
By \Cref{prop:ch26-truncation-limit} the left-hand side decreases to
\(E_\omega(\psi)\), and by
\Cref{lem:integration-ledger}\,(2) the integral decreases to
\(\int(\psi-\varphi)\MA_\omega(\varphi)\), both limits being taken in
\([-\infty,+\infty)\).  This gives \eqref{eq:ch26-supergradient}.

For \eqref{eq:ch26-energy-as-infimum}, the inequality \(\leq\) is
\eqref{eq:ch26-supergradient}.  For the reverse, take
\(\varphi=\psi^{\langle t\rangle}\), which is bounded and satisfies
\(\psi-\psi^{\langle t\rangle}\leq0\); then
\[
 E_\omega(\psi^{\langle t\rangle})
  +\int(\psi-\psi^{\langle t\rangle})
    \MA_\omega(\psi^{\langle t\rangle})
 \leq E_\omega(\psi^{\langle t\rangle})
 \longrightarrow E_\omega(\psi),
\]
again by \Cref{prop:ch26-truncation-limit}.  Hence the infimum is at most
\(E_\omega(\psi)\).
\end{proof}

\begin{remark}[Differentiability: what is available where]
\label{rem:ch26-differentiability-scope}
Three statements should be kept apart.  Along a segment of \emph{model}
potentials, \(E_\omega\) is a polynomial of degree at most \(n+1\) and its
derivative is \eqref{eq:energy-first-variation}; this is
\Cref{cor:ch26-energy-polynomial}, and it is an identity between intersection
numbers.  Along a segment of \emph{bounded} potentials the same holds, by
\Cref{prop:bounded-energy-identities}.  Along a segment inside \(\cE^1\) the
function \(t\mapsto E_\omega((1-t)\varphi+t\psi)\) is still differentiable,
with derivative \(\int(\psi-\varphi)\MA_\omega(\varphi_t)\), but the proof
needs the finite-energy mixed products of
\Cref{sec:ch26-mixed-products} and is deferred to
\Cref{thm:e1-energy-differentiability}.  Finally, differentiability in a
\emph{continuous} direction that leaves the psh cone is a different and much
harder question: it requires composing with the envelope \(P_\omega\), and
its answer is equivalent to the orthogonality property
\eqref{eq:orthogonality}.  That is the subject of
\Cref{chap:bfj-differentiability}, and nothing in the present chapter
anticipates it.
\end{remark}

\begin{example}[Upper semicontinuity is strict, and \(E_\omega\) is not
continuous]
\label{ex:ch26-usc-strict}
Let \(K=\mathbf C((\pi))\), \(X=\mathbf P^1_K\), \(L=\cO(1)\), and let
\(\omega\) be the curvature of the model metric of
\((\mathbf P^1_R,\cO(1))\) as in \Cref{ex:ch21-p1-refinement}, so that
\(\omega=\delta_\eta\) with \(\eta\) the Gauss point and \(V=1\).  Choose
\(a_1,a_2,\ldots\in R\) in pairwise distinct residue classes as in
\eqref{eq:ch21-separated-centres}, write
\(\varphi_{a}:=\log\lVert z_1-az_0\rVert_0=\min\{0,\log|z-a|\}\) as in
\Cref{ex:ch21-polar-noncompact}, and put
\[
 \tau_N:=\max\{\varphi_{a_N},-1\}\in\PSH(\mathbf P^1,\omega).
\]
Set \(t_N:=-\varphi_{a_N}\in[0,+\infty]\), so that \(t_N=0\) exactly on
\(\{|z-a_N|\geq1\}\) and \(\tau_N=-\min\{t_N,1\}\).  On the ray
\(R_N=[\eta,a_N]\), parametrized by \(t=-\log r\), the function \(t_N\)
restricts to \(t\).  Thus \(\sup\tau_N=0\) and \(\tau_N\) is bounded.  The
outgoing slope of \(\tau_N\) at \(\eta\) is \(-1\), cancelling the atom of
\(\omega\) there, and the slope jumps from \(-1\) to \(0\) at \(t=1\); so, by
\eqref{eq:curve-ma-laplacian} and \Cref{prop:ch26-graph-energy},
\[
 \MA_\omega(\tau_N)=\delta_{t=1\text{ on }R_N},
 \qquad
 E_\omega(\tau_N)=0-\tfrac12\int_0^11\,\rd t=-\tfrac12 .
\]
The open discs \(D_N:=\{|z-a_N|<1\}\) are pairwise disjoint, and
\(\tau_N\) vanishes outside \(D_N\).  A dual complex \(\Delta\) is a finite
graph containing \(\eta\), so it has finite valence at \(\eta\) and
therefore meets only finitely many of the \(D_N\); consequently
\(\tau_N|_\Delta\equiv0\) for \(N\) large.  Hence \(\tau_N\to0\) uniformly on
every dual complex, that is, in the BFJ topology
(\Cref{thm:bfj-skeletal-equicontinuity},
\Cref{rem:ch22-topology-comparison}), and
\[
 \varlimsup_N E_\omega(\tau_N)=-\tfrac12<0=E_\omega(0),
\]
so upper semicontinuity holds strictly and \(E_\omega\) is \emph{not}
continuous for the BFJ topology.  The sequence is not monotone, so
Step~4 of \Cref{thm:extended-energy-basic} does not apply; this is the
distinction announced in \Cref{rem:ch22-topology-comparison}.
\end{example}

\begin{definition}[Finite energy]
\label{def:e1}
\[
 \mathcal E^1(X,\omega)
 :=\{\varphi\in\PSH(X,\omega)\mid E_\omega(\varphi)>-\infty\}.
\]
\end{definition}

By \Cref{prop:ch26-truncation-limit}, a potential normalized by
\(\varphi\leq0\) lies in \(\cE^1(X,\omega)\) precisely when its truncation
energies are bounded below.  The elementary structure of the class is
recorded next; the deeper properties occupy
\Cref{sec:ch26-nonpluripolar,sec:ch26-mixed-products,sec:ch26-e1-gauge}.

\begin{theorem}
\label{thm:ch26-e1-structure}
The class \(\cE^1(X,\omega)\) has the following properties.
\begin{enumerate}
\item\label{itm:ch26-e1:bounded}
 It contains every bounded \(\omega\)-psh function, in particular every
 constant and every \(\omega\)-psh model function.
\item\label{itm:ch26-e1:upward}
 It is an upward-closed subset of \(\PSH(X,\omega)\): if
 \(\varphi\in\cE^1\) and \(\varphi\leq\psi\in\PSH(X,\omega)\), then
 \(\psi\in\cE^1\).
\item\label{itm:ch26-e1:constants}
 It is stable under addition of real constants.
\item\label{itm:ch26-e1:convex}
 It is convex, and more precisely
 \(E_\omega((1-s)\varphi+s\psi)\geq(1-s)E_\omega(\varphi)+sE_\omega(\psi)\)
 for \(\varphi,\psi\in\cE^1\) and \(s\in[0,1]\).
\item\label{itm:ch26-e1:max}
 It is stable under finite maxima, and
 \(E_\omega(\max\{\varphi,\psi\})\geq\max\{E_\omega(\varphi),E_\omega(\psi)\}\).
\item\label{itm:ch26-e1:truncation}
 It is stable under truncation: \(\varphi^{\langle t\rangle}\in\cE^1\) for
 every \(\varphi\in\cE^1\) bounded above and every \(t>0\).
\item\label{itm:ch26-e1:not-closed}
 It is \emph{not} closed in the BFJ topology, and it is not a lattice for the
 pointwise minimum.
\end{enumerate}
\end{theorem}

\begin{proof}
\cref{itm:ch26-e1:bounded} is \Cref{lem:ch26-consistency} together with
\Cref{prop:ch26-energy-lipschitz}.  \cref{itm:ch26-e1:upward} and
\cref{itm:ch26-e1:constants} are Step~1 of
\Cref{thm:extended-energy-basic}.  \cref{itm:ch26-e1:convex} is Step~2 of the
same theorem, and it gives convexity of \(\cE^1\) because the right-hand side
is finite when both energies are.  \cref{itm:ch26-e1:max} follows from
\cref{itm:ch26-e1:upward} because \(\max\{\varphi,\psi\}\geq\varphi\), and
\(\max\{\varphi,\psi\}\in\PSH(X,\omega)\) by
\Cref{thm:ch21-stability}\,\cref{itm:ch21-stab:max}.  \cref{itm:ch26-e1:truncation}
is the case \(\psi=-t\) of \cref{itm:ch26-e1:max}.

For \cref{itm:ch26-e1:not-closed}, \Cref{ex:ch26-infinite-energy} produces
\(\varphi\in\PSH(\mathbf P^1,\omega)\setminus\cE^1\) which is the limit in
the BFJ topology of the sequence of its truncations, each bounded and hence
in \(\cE^1\); so \(\cE^1\) is not closed.  The failure to be a lattice is
already visible for model functions: the pointwise minimum of two
\(\omega\)-psh model functions need not be \(\omega\)-psh, by
\Cref{ex:ch20-min-fails}.
\end{proof}

\begin{warning}{Polar sets and singularity types}
Within a fixed class $\PSH(X,\omega)$, finite energy is invariant under the singularity type of \Cref{def:ch21-polar-singularity}. Indeed, if $\varphi-C\leq\psi\leq\varphi+C$ for $C\geq0$, monotonicity and translation of energy give
\[
 E_\omega(\varphi)-CV\leq E_\omega(\psi)
                  \leq E_\omega(\varphi)+CV.
\]
Thus one energy is finite exactly when the other is finite. More generally, if $\varphi\in\cE^1$ and $\varphi\preceq\psi$, then $\psi\in\cE^1$, by upward closure and invariance under constants.

The polar set alone does not determine finite energy. The potentials in \Cref{ex:ch26-unbounded-finite-energy,ex:ch26-infinite-energy} have the same single rational polar point and different energy behavior. Their difference is unbounded along the corresponding ray, so they do not have the same singularity type. Finite energy implies the tail estimate \eqref{eq:e1-mass-tail}.
\end{warning}

\section{The nonpluripolar measure}
\label{sec:ch26-nonpluripolar}

The bounded product of \Cref{thm:bounded-ma-extension} does not apply to an
unbounded potential, and the naive remedy --- take the weak limit of the
measures of the truncations --- produces a measure that may charge the pole
set and is therefore useless for the variational method.  The correct
construction discards exactly that mass.  It uses one input only, the
locality statement contained in the bounded maximum principle
\eqref{eq:bounded-maximum-principle}, and it produces an increasing rather
than a weak limit.

\begin{definition}[Nonpluripolar Monge--Amp\`ere measure]
\label{def:nonpluripolar-ma}
After translating \(\varphi\) so that \(\varphi\leq0\), define
\begin{equation}
\label{eq:nonpluripolar-ma}
 \MA_\omega(\varphi)
 :=\lim_{t\to\infty}
 \one_{\{\varphi>-t\}}
 \MA_\omega(\varphi^{\langle t\rangle}).
\end{equation}
The limit is increasing on every Borel set. It does not depend on the
translation.
\end{definition}

\begin{lemma}
\label{lem:ch26-layers}
Let \(\varphi\in\PSH(X,\omega)\) with \(\varphi\leq0\) and let
\(0<t\leq s\).  Then \(\{\varphi>-t\}\) is Borel and
\begin{equation}
\label{eq:ch26-layer-consistency}
 \one_{\{\varphi>-t\}}\MA_\omega(\varphi^{\langle s\rangle})
 =\one_{\{\varphi>-t\}}\MA_\omega(\varphi^{\langle t\rangle}).
\end{equation}
Consequently, for every Borel set \(B\) the quantity
\(\MA_\omega(\varphi^{\langle t\rangle})(B\cap\{\varphi>-t\})\) is
nondecreasing in \(t\), and
\begin{equation}
\label{eq:ch26-npp-definition}
 \MA_\omega(\varphi)(B)
 :=\lim_{t\to\infty}
  \MA_\omega(\varphi^{\langle t\rangle})
   \bigl(B\cap\{\varphi>-t\}\bigr)
\end{equation}
defines a positive Borel measure of total mass at most \(V\), carried by
\(\{\varphi>-\infty\}\), which is independent of the normalization
\(\varphi\leq0\) in the sense that replacing \(\varphi\) by \(\varphi-c\)
with \(c\geq0\) yields the same measure.
\end{lemma}

\begin{proof}
Upper semicontinuity makes \(\varphi\) Borel, hence
\(\{\varphi>-t\}\) Borel; this strict superlevel set need not be open in
the ordinary Berkovich topology.  The point is instead that
\[
 \varphi^{\langle t\rangle}
   =\max\{\varphi^{\langle s\rangle},-t\},
 \qquad
 \{\varphi^{\langle s\rangle}>-t\}=\{\varphi>-t\}.
\]
Apply the measure identity \eqref{eq:bounded-maximum-principle} to the
bounded \(\omega\)-psh functions \(\varphi^{\langle s\rangle}\) and the
constant \(-t\).  It gives exactly
\eqref{eq:ch26-layer-consistency}, without any topological openness claim.
Monotonicity in \(t\) follows: for
\(t\leq s\), put \(C_r:=B\cap\{\varphi>-r\}\).  Then
\[
 \begin{aligned}
 \MA_\omega(\varphi^{\langle s\rangle})(C_s)
 &\geq\MA_\omega(\varphi^{\langle s\rangle})(C_t)\\
 &=\MA_\omega(\varphi^{\langle t\rangle})(C_t).
 \end{aligned}
\]
Thus \eqref{eq:ch26-npp-definition} is a well-defined limit in
\([0,V]\), each measure \(\MA_\omega(\varphi^{\langle t\rangle})\) having
mass \(V\).  Restricting to a countable cofinal set of levels
\(t\in\mathbf N\) --- which is legitimate because the limit is monotone ---
countable additivity of the limit follows from countable additivity of each
term together with the monotone convergence theorem for series of
nonnegative numbers.  The resulting set function \(\nu\) is therefore a
finite positive Borel measure.  It is inner regular: if \(B\) is Borel and
\(\varepsilon>0\), choose an integer \(m\) so that
\(\nu(B)-\nu_m(B)<\varepsilon\), where
\[
 \nu_m:=\one_{\{\varphi>-m\}}
          \MA_\omega(\varphi^{\langle m\rangle}).
\]
The finite Radon measure \(\nu_m\) admits a compact set \(K\subseteq B\)
with \(\nu_m(B)-\nu_m(K)<\varepsilon\).  Since \(\nu\geq\nu_m\),
\[
 0\leq\nu(B)-\nu(K)
 \leq\bigl(\nu(B)-\nu_m(B)\bigr)
      +\bigl(\nu_m(B)-\nu_m(K)\bigr)<2\varepsilon.
\]
Thus \(\nu\) is Radon.

That it is carried by \(\{\varphi>-\infty\}\) is immediate from
\eqref{eq:ch26-npp-definition}, since
\(\{\varphi>-\infty\}=\bigcup_{t}\{\varphi>-t\}\).  Finally, replacing
\(\varphi\) by \(\varphi-c\) replaces the layer \(\{\varphi>-t\}\) by
\(\{\varphi-c>-t-c\}=\{\varphi>-t\}\) and the truncation
\((\varphi-c)^{\langle t+c\rangle}\) by
\(\varphi^{\langle t\rangle}-c\), whose Monge--Amp\`ere measure is the same;
so the two increasing limits run over the same cofinal family of measures.
\end{proof}

\begin{proposition}
\label{prop:ch26-npp-basic}
Let \(\varphi\in\PSH(X,\omega)\).  Then the measure
\(\MA_\omega(\varphi)\) of \Cref{def:nonpluripolar-ma}
\begin{enumerate}
\item\label{itm:ch26-npp:bounded}
 coincides with the bounded product of \Cref{thm:bounded-ma-extension} when
 \(\varphi\) is bounded;
\item\label{itm:ch26-npp:mass}
 has total mass at most \(V\);
\item\label{itm:ch26-npp:polar}
 gives no mass to any Borel pluripolar set;
\item\label{itm:ch26-npp:local}
 is local: if \(\varphi=\psi\) on an open set \(U\), then
 \(\MA_\omega(\varphi)|_U=\MA_\omega(\psi)|_U\).
\end{enumerate}
\end{proposition}

\begin{proof}
\cref{itm:ch26-npp:bounded}: if \(\varphi\geq-t_0\), then
\(\varphi^{\langle t\rangle}=\varphi\) and \(\{\varphi>-t\}=X^{\An}\) for
\(t>t_0\).  \cref{itm:ch26-npp:mass} is \Cref{lem:ch26-layers}.
\cref{itm:ch26-npp:polar}: on each layer the measure is the restriction of a
bounded product, which gives no mass to Borel pluripolar sets by
\Cref{prop:bounded-ma-no-polar-mass}; an increasing limit of measures each
vanishing on a fixed Borel set vanishes on it.
\cref{itm:ch26-npp:local}: on \(U\cap\{\varphi>-t\}=U\cap\{\psi>-t\}\) the
truncations agree, so \eqref{eq:ch26-layer-consistency} and the locality
clause of \Cref{thm:bounded-maximum-principle} give equality of the
corresponding layers; pass to the limit.
\end{proof}

For arbitrary psh \(\varphi\), this measure has mass at most \(V\), is
supported on \(\{\varphi>-\infty\}\), and does not charge pluripolar sets.
It may lose mass at the pole set;
\Cref{ex:ch26-infinite-energy} exhibits a potential on \(\mathbf P^1\) whose
nonpluripolar measure is \emph{zero}, all the mass having escaped to a single
rational point.

\begin{definition}
A psh function has \emph{full Monge--Amp\`ere mass} if
\(\MA_\omega(\varphi)(X^{\An})=V\).
\end{definition}

\begin{warning}{Which product \(\MA_\omega\) denotes}
From here to the end of the chapter, \(\MA_\omega(\varphi)\) denotes the
nonpluripolar measure \eqref{eq:nonpluripolar-ma} whenever \(\varphi\) is
unbounded, and the bounded product of \Cref{thm:bounded-ma-extension}
otherwise; the two agree on the overlap by
\Cref{prop:ch26-npp-basic}\,\cref{itm:ch26-npp:bounded}, which is what makes
the single symbol admissible.  Neither of them is the local
Bedford--Taylor product of \Cref{chap:cld-bedford-taylor}, and neither is
the Chambert--Loir measure of a metric that is merely continuous unless the
comparison \Cref{thm:bfj-cld-ma-common-locus} applies.
\end{warning}

\begin{proposition}\label{prop:e1-integrability-tail}
Let \(\varphi\in\mathcal E^1(X,\omega)\), normalized by \(\varphi\leq0\).
Then
\begin{equation}
\label{eq:e1-mass-tail}
 \MA_\omega(\varphi^{\langle t\rangle})
       (\{\varphi\leq-t\})=o(t^{-1}),
\end{equation}
so \(\varphi\) has full Monge--Amp\`ere mass. Moreover
\begin{equation}
\label{eq:e1-integrability}
 -E_\omega(\varphi)
 \leq\int(-\varphi)\MA_\omega(\varphi)
 \leq(n+1)(-E_\omega(\varphi))<+\infty.
\end{equation}
\end{proposition}

\begin{proof}
\emph{Step 1: the tail bound.}  Put
\(\mu_t=\MA_\omega(\varphi^{\langle t\rangle})\), a positive measure of mass
\(V\) supplied by the bounded product.  By
\eqref{eq:energy-cocycle}, in the bounded form of
\Cref{prop:bounded-energy-identities}, applied to the pair
\(\varphi^{\langle t\rangle}\leq\varphi^{\langle t/2\rangle}\), and by
positivity of all \(n+1\) mixed measures,
\[
 E_\omega(\varphi^{\langle t/2\rangle})
 -E_\omega(\varphi^{\langle t\rangle})
 \geq\frac1{n+1}
 \int\bigl(\varphi^{\langle t/2\rangle}
          -\varphi^{\langle t\rangle}\bigr)\,\rd \mu_t ,
\]
where we have kept only the summand with all \(n\) entries equal to
\(\varphi^{\langle t\rangle}\).  On the set \(\{\varphi\leq-t\}\) one has
\(\varphi^{\langle t/2\rangle}=-t/2\) and
\(\varphi^{\langle t\rangle}=-t\), so the integrand equals \(t/2\) there and
is nonnegative everywhere.  Hence
\begin{equation}
\label{eq:ch26-tail-chebyshev}
 E_\omega(\varphi^{\langle t/2\rangle})
 -E_\omega(\varphi^{\langle t\rangle})
 \geq\frac{t}{2(n+1)}\,\mu_t\bigl(\{\varphi\leq-t\}\bigr).
\end{equation}
By \Cref{prop:ch26-truncation-limit} both energies converge to the finite
number \(E_\omega(\varphi)\) as \(t\to\infty\), so the left-hand side of
\eqref{eq:ch26-tail-chebyshev} tends to \(0\).  This is
\eqref{eq:e1-mass-tail}.

\emph{Step 2: full mass.}  By \eqref{eq:ch26-npp-definition},
\[
 \MA_\omega(\varphi)(X^{\An})
 =\lim_{t\to\infty}\mu_t(\{\varphi>-t\})
 =\lim_{t\to\infty}\bigl(V-\mu_t(\{\varphi\leq-t\})\bigr)=V,
\]
by Step~1.

\emph{Step 3: the left inequality of \eqref{eq:e1-integrability}.}  Fix
\(t>0\) and apply \eqref{eq:energy-sandwich} to
\(\varphi^{\langle t\rangle}\leq0\):
\[
 E_\omega(0)-E_\omega(\varphi^{\langle t\rangle})
 \leq\int\bigl(0-\varphi^{\langle t\rangle}\bigr)
      \MA_\omega(\varphi^{\langle t\rangle})
 =\int(-\varphi^{\langle t\rangle})\,\rd\mu_t .
\]
Since \(E_\omega(0)=0\), this reads
\(-E_\omega(\varphi^{\langle t\rangle})\leq\int(-\varphi^{\langle
t\rangle})\rd\mu_t\).

\emph{Step 4: the right inequality.}  In \eqref{eq:model-energy} for
\(\varphi^{\langle t\rangle}\), every one of the \(n+1\) summands
\(\int\varphi^{\langle t\rangle}
\omega_{\varphi^{\langle t\rangle}}^{\,j}\wedge\omega^{\,n-j}\)
is nonpositive, because \(\varphi^{\langle t\rangle}\leq0\) and the measures
are positive; the summand with \(j=n\) is
\(\int\varphi^{\langle t\rangle}\rd\mu_t\).  Hence
\[
 \int(-\varphi^{\langle t\rangle})\,\rd\mu_t
 \leq(n+1)\bigl(-E_\omega(\varphi^{\langle t\rangle})\bigr).
\]

\emph{Step 5: passage to the limit.}  Combining Steps~3 and~4 with
\(E_\omega(\varphi^{\langle t\rangle})\geq E_\omega(\varphi)\),
\[
 -E_\omega(\varphi^{\langle t\rangle})
 \leq\int(-\varphi^{\langle t\rangle})\,\rd\mu_t
 \leq(n+1)\bigl(-E_\omega(\varphi)\bigr)
\]
for every \(t\).  It remains to identify the limit of the middle term with
\(\int(-\varphi)\MA_\omega(\varphi)\).  Since
\(\varphi^{\langle t\rangle}=\varphi\) on \(\{\varphi>-t\}\) and equals
\(-t\) on its complement,
\begin{equation}
\label{eq:ch26-tail-split}
 \int(-\varphi^{\langle t\rangle})\,\rd\mu_t
 =\int_{\{\varphi>-t\}}(-\varphi)\,\rd\mu_t
  +t\,\mu_t\bigl(\{\varphi\leq-t\}\bigr).
\end{equation}
The second term tends to \(0\) by \eqref{eq:e1-mass-tail}.  For the first,
\eqref{eq:ch26-layer-consistency} identifies \(\one_{\{\varphi>-t\}}\mu_t\)
with the restriction of \(\MA_\omega(\varphi)\) to the layer
\(\{\varphi>-t\}\), so
\[
 \int_{\{\varphi>-t\}}(-\varphi)\,\rd\mu_t
 =\int_{\{\varphi>-t\}}(-\varphi)\,\MA_\omega(\varphi)
 \nearrow\int(-\varphi)\MA_\omega(\varphi)
\]
by monotone convergence over the increasing family of layers, whose union
carries \(\MA_\omega(\varphi)\) and on which the integrand \(-\varphi\) is
nonnegative.  Hence
\(\int(-\varphi^{\langle t\rangle})\rd\mu_t\to
\int(-\varphi)\MA_\omega(\varphi)\), and the displayed two-sided bound passes
to the limit, using \(E_\omega(\varphi^{\langle t\rangle})\to
E_\omega(\varphi)\) from \Cref{prop:ch26-truncation-limit}.  All integrals
here are of nonnegative functions, so no undefined difference is formed, in
accordance with entry~\cref{itm:ch26-ledger:integral} of
\Cref{sec:ch26-extended-real}.
\end{proof}

\begin{corollary}
\label{cor:ch26-quantitative-tail}
Let \(\varphi\in\cE^1(X,\omega)\) with \(\varphi\leq0\).  Then for every
\(t>0\),
\begin{equation}
\label{eq:ch26-quantitative-tail}
 \begin{aligned}
 \MA_\omega(\varphi^{\langle t\rangle})
   \bigl(\{\varphi\leq-t\}\bigr)
 &\leq\frac{n+1}{t}\bigl(-E_\omega(\varphi)\bigr),\\
 \MA_\omega(\varphi)\bigl(\{\varphi\leq-t\}\bigr)
 &\leq\frac{n+1}{t}\bigl(-E_\omega(\varphi)\bigr).
 \end{aligned}
\end{equation}
\end{corollary}

\begin{proof}
Steps~4 and~5 of the previous proof give
\(\int(-\varphi^{\langle t\rangle})\rd\mu_t\leq
(n+1)\bigl(-E_\omega(\varphi^{\langle t\rangle})\bigr)
\leq(n+1)\bigl(-E_\omega(\varphi)\bigr)\).
Chebyshev's inequality applied to the nonnegative function
\(-\varphi^{\langle t\rangle}\), which equals \(t\) on
\(\{\varphi\leq-t\}\), gives the first bound.  The second is Chebyshev
applied to \(-\varphi\) and \(\MA_\omega(\varphi)\), using
\eqref{eq:e1-integrability}.
\end{proof}

The estimate \eqref{eq:ch26-quantitative-tail} is the workhorse of the
chapter: it converts a single scalar bound \(E_\omega(\varphi)\geq-C\) into a
uniform rate for every truncation argument that follows, and its uniformity
over a family with a common energy bound is what makes
\Cref{prop:e1-compact-superlevels} useful.

\begin{corollary}
\label{cor:e1-capacity-tail}
For \(\varphi\in\mathcal E^1(X,\omega)\),
\begin{equation}
\label{eq:e1-capacity-tail}
                       \Capa_\omega\{\varphi<-t\}\longrightarrow0.
\end{equation}
Thus \Cref{thm:e1-quasicontinuity} applies.
\end{corollary}

\begin{proof}
Normalize \(\varphi\leq0\). Apply
\eqref{eq:capacity-comparison} to a truncation and the constant obstacle
\(-t\), with the interpolation parameter \(1/2\), and then let the
truncation level tend to infinity. The right side is bounded by a constant
times
\[
 \MA_\omega(\varphi)(\{\varphi<-t/2+1\})
 \leq\frac{2}{t-2}\int(-\varphi)\MA_\omega(\varphi),
\]
which tends to zero by \eqref{eq:e1-integrability}.
\end{proof}

\begin{warning}{Extended-real arithmetic}
For \(\varphi,\psi\in\mathcal E^1\), the expression
\(\psi-\varphi\) is not defined where both functions equal \(-\infty\).
Every integral below involving that difference is obtained from bounded
truncations or from a monotone cocycle identity. It is never interpreted as
pointwise subtraction on the common pole set.
\end{warning}

\section{Mixed products in finite energy}
\label{sec:ch26-mixed-products}

We now prove the extension announced in
\Cref{sec:ch26-four-products}: the bounded mixed product of
\Cref{thm:bounded-ma-extension} extends to \(n\)-tuples of finite-energy
potentials, without loss of mass.  The proof is a monotone-limit argument
with three separate ingredients, and we isolate them before assembling the
theorem, because two of them are of independent use.

\subsection{Locality of the bounded mixed product}
\label{sec:ch26-mixed-locality}

\Cref{thm:bounded-maximum-principle} is a statement about the \emph{diagonal}
bounded product.  What the truncation argument needs is the mixed version.
It is not an additional theorem: it follows from the diagonal case by
polarization, and we give the argument because the polarization is performed
inside the convex set of \(\omega\)-psh functions and therefore has to be
arranged with some care.

\begin{lemma}
\label{lem:ch26-mixed-locality}
Let \(u,v,w_2,\ldots,w_n\) be bounded \(\omega\)-psh functions.  Then
\begin{equation}
\label{eq:ch26-mixed-maximum-principle}
 \one_{\{u>v\}}\MA_\omega(\max\{u,v\},w_2,\ldots,w_n)
 =\one_{\{u>v\}}\MA_\omega(u,w_2,\ldots,w_n).
\end{equation}
Moreover, if \(u_1,\ldots,u_n\) and \(v_1,\ldots,v_n\) are bounded
\(\omega\)-psh functions and \(U\subseteq X^{\An}\) is open with
\(u_i=v_i\) on \(U\) for every \(i\), then
\begin{equation}
\label{eq:ch26-mixed-locality}
 \one_U\,\MA_\omega(u_1,\ldots,u_n)
 =\one_U\,\MA_\omega(v_1,\ldots,v_n).
\end{equation}
\end{lemma}

\begin{proof}
We first prove \eqref{eq:ch26-mixed-maximum-principle}.  Put
\(A:=\{u>v\}\).  For \(s=(s_2,\ldots,s_n)\in\mathbf R_{\geq0}^{n-1}\)
and \(S:=\sum_{j=2}^ns_j\), set
\[
 u_s:=\frac{u+\sum_{j=2}^ns_jw_j}{1+S},
 \qquad
 v_s:=\frac{v+\sum_{j=2}^ns_jw_j}{1+S}.
\]
These are bounded \(\omega\)-psh functions,
\(\{u_s>v_s\}=A\), and
\[
 \max\{u_s,v_s\}
 =\frac{\max\{u,v\}+\sum_{j=2}^ns_jw_j}{1+S}.
\]
Applying \eqref{eq:bounded-maximum-principle} to \(u_s,v_s\), then
multiplying by \((1+S)^n\), gives the following identity of measures on
\(A\):
\[
 \left(\omega_{\max\{u,v\}}+\sum_{j=2}^ns_j\omega_{w_j}\right)^{\!n}
 =
 \left(\omega_u+\sum_{j=2}^ns_j\omega_{w_j}\right)^{\!n}.
\]
Evaluate on an arbitrary Borel subset of \(A\).  Both sides are
polynomials in \(s_2,\ldots,s_n\); equality for every nonnegative \(s\)
implies equality of their coefficients.  The coefficient of
\(s_2\cdots s_n\) is \(n!\) times the mixed measure in
\eqref{eq:ch26-mixed-maximum-principle}, proving that identity.  For
\(n=1\) the same argument is simply
\eqref{eq:bounded-maximum-principle} itself.

We now prove the open-set assertion.
Let \(\Sigma:=\{s\in\mathbf R_{\geq0}^n:\sum_is_i=1\}\) be the standard
simplex.  For \(s\in\Sigma\) the functions
\(u_s:=\sum_is_iu_i\) and \(v_s:=\sum_is_iv_i\) are bounded
\(\omega\)-psh, by
\Cref{thm:ch21-stability}\,\cref{itm:ch21-stab:sum} together with
\(\sum_is_i=1\), and they agree on \(U\).  The locality clause of
\Cref{thm:bounded-maximum-principle} gives
\begin{equation}
\label{eq:ch26-polarization-identity}
 \one_U\,\MA_\omega(u_s)=\one_U\,\MA_\omega(v_s)
 \qquad(s\in\Sigma).
\end{equation}
By multilinearity of the bounded mixed product
(\Cref{cor:bounded-ma-multilinearity}),
\[
 \MA_\omega(u_s)
 =\sum_{|\alpha|=n}\frac{n!}{\alpha!}\,s^\alpha\,
   \MA_\omega(u_1^{(\alpha_1)},\ldots,u_n^{(\alpha_n)}),
\]
the sum being over multi-indices \(\alpha\in\mathbf N^n\) of total degree
\(n\), and likewise for \(v_s\).  Fix a Borel set \(B\subseteq U\) and
evaluate: \eqref{eq:ch26-polarization-identity} says that the two
homogeneous polynomials of degree \(n\) in \(s\),
\[
 s\longmapsto\sum_{|\alpha|=n}\frac{n!}{\alpha!}s^\alpha
  \MA_\omega(u^{(\alpha)})(B)
 \quad\text{and}\quad
 s\longmapsto\sum_{|\alpha|=n}\frac{n!}{\alpha!}s^\alpha
  \MA_\omega(v^{(\alpha)})(B),
\]
agree on \(\Sigma\).  A homogeneous polynomial of degree \(n\) on
\(\mathbf R^n\) vanishing on \(\Sigma\) vanishes identically, because
\(\Sigma\) spans \(\mathbf R^n\) affinely and the polynomial is determined by
its restriction to the ray through each point of \(\Sigma\).  Hence the
coefficients agree, and taking \(\alpha=(1,\ldots,1)\) gives
\(\MA_\omega(u_1,\ldots,u_n)(B)=\MA_\omega(v_1,\ldots,v_n)(B)\) for every
Borel \(B\subseteq U\).
\end{proof}

\subsection{The mixed tail estimate}
\label{sec:ch26-mixed-tail}

The second ingredient is the quantitative statement that replaces
\eqref{eq:e1-mass-tail} for mixed products.  Its proof rests on one
elementary inequality between a mixed product and the diagonal product of the
average, which we state separately.

\begin{lemma}
\label{lem:ch26-average-domination}
Let \(u_1,\ldots,u_n\) be bounded \(\omega\)-psh functions and put
\(\bar u:=\frac1n\sum_iu_i\).  Then
\begin{equation}
\label{eq:ch26-average-domination}
 \MA_\omega(u_1,\ldots,u_n)
 \leq\frac{n^n}{n!}\,\MA_\omega(\bar u)
\end{equation}
as measures.
\end{lemma}

\begin{proof}
By \Cref{cor:bounded-ma-multilinearity},
\[
 \MA_\omega(\bar u)
 =\frac1{n^n}\sum_{|\alpha|=n}\frac{n!}{\alpha!}
   \MA_\omega(u^{(\alpha)}),
\]
a sum of positive measures.  Retaining only the term with
\(\alpha=(1,\ldots,1)\), whose coefficient is \(n!/1=n!\), gives
\(\MA_\omega(\bar u)\geq\frac{n!}{n^n}\MA_\omega(u_1,\ldots,u_n)\).
\end{proof}

\begin{lemma}
\label{lem:ch26-mixed-tail}
Let \(\varphi_1,\ldots,\varphi_n\in\cE^1(X,\omega)\) with
\(\varphi_i\leq0\), and put
\(\bar\varphi:=\frac1n\sum_i\varphi_i\in\cE^1(X,\omega)\).  Write
\[
 U_t:=\bigcap_{i=1}^n\{\varphi_i>-t\},
 \qquad
 B_t:=X^{\An}\setminus U_t=\bigcup_{i=1}^n\{\varphi_i\leq-t\}.
\]
Then, for every \(t>0\),
\begin{equation}
\label{eq:ch26-mixed-tail}
 \MA_\omega\bigl(\varphi_1^{\langle t\rangle},\ldots,
   \varphi_n^{\langle t\rangle}\bigr)(B_t)
 \ \leq\ \frac{n^{n+1}(n+1)}{n!}\cdot
   \frac{-E_\omega(\bar\varphi)}{t}.
\end{equation}
In particular the left-hand side tends to \(0\) as \(t\to\infty\), uniformly
over any family of \(n\)-tuples for which \(-E_\omega(\bar\varphi)\) is
bounded above.
\end{lemma}

\begin{proof}
That \(\bar\varphi\in\cE^1\) is
\Cref{thm:ch26-e1-structure}\,\cref{itm:ch26-e1:convex}, and
\(\bar\varphi\leq0\).  Put
\[
 \bar\varphi_t:=\frac1n\sum_{i=1}^n\varphi_i^{\langle t\rangle},
\]
a bounded \(\omega\)-psh function with \(-t\leq\bar\varphi_t\leq0\).

\emph{Step 1: the bad set is a sublevel of \(\bar\varphi_t\).}  Let
\(x\in B_t\), say \(\varphi_i(x)\leq-t\).  Then
\(\varphi_i^{\langle t\rangle}(x)=-t\), while
\(\varphi_j^{\langle t\rangle}(x)\leq0\) for every \(j\), because
\(\varphi_j\leq0\).  Hence
\[
 \bar\varphi_t(x)\leq\frac1n(-t)=-\frac tn ,
\]
so \(B_t\subseteq\{\bar\varphi_t\leq-t/n\}\).

\emph{Step 2: Chebyshev.}  By \Cref{lem:ch26-average-domination} applied to
the \(n\)-tuple of truncations,
\[
 \begin{aligned}
 &\MA_\omega\bigl(\varphi_1^{\langle t\rangle},\ldots,
    \varphi_n^{\langle t\rangle}\bigr)(B_t)\\
 &\quad\leq\frac{n^n}{n!}\,
   \MA_\omega(\bar\varphi_t)
      \bigl(\{\bar\varphi_t\leq-t/n\}\bigr)\\
 &\quad\leq\frac{n^n}{n!}\cdot\frac{n}{t}
   \int(-\bar\varphi_t)\,\MA_\omega(\bar\varphi_t),
 \end{aligned}
\]
the last step being Chebyshev's inequality for the nonnegative function
\(-\bar\varphi_t\).

\emph{Step 3: the energy bound.}  The function \(\bar\varphi_t\) is bounded,
\(\omega\)-psh and \(\leq0\), so Step~4 of the proof of
\Cref{prop:e1-integrability-tail} gives
\[
 \int(-\bar\varphi_t)\MA_\omega(\bar\varphi_t)
 \leq(n+1)\bigl(-E_\omega(\bar\varphi_t)\bigr).
\]
Moreover \(\bar\varphi_t\geq\bar\varphi\), so
\(E_\omega(\bar\varphi_t)\geq E_\omega(\bar\varphi)\) by
\Cref{thm:extended-energy-basic}, whence
\(-E_\omega(\bar\varphi_t)\leq-E_\omega(\bar\varphi)\).  Combining the three
steps gives \eqref{eq:ch26-mixed-tail}.
\end{proof}

\begin{remark}[Why the average, and not each factor separately]
\label{rem:ch26-why-average}
One might try to bound the mixed measure of \(B_t\) by summing the diagonal
tail bounds \eqref{eq:e1-mass-tail} for the individual \(\varphi_i\).  That
does not work, because \eqref{eq:e1-mass-tail} controls
\(\MA_\omega(\varphi_i^{\langle t\rangle})\) and not the mixed measure, and
no inequality bounds a mixed Monge--Amp\`ere measure of a set by the
diagonal measures of the same set.  \Cref{lem:ch26-average-domination} is the
substitute: it replaces the \(n\) potentials by a single one, at the cost of
the explicit constant \(n^n/n!\), which is harmless because only the rate
\(O(t^{-1})\) matters.  The device also explains why \(\cE^1\) has to be
convex before mixed products can be constructed, and hence why
\Cref{thm:ch26-e1-structure}\,\cref{itm:ch26-e1:convex} precedes this section.
\end{remark}

\subsection{The extension theorem}
\label{sec:ch26-extension-theorem}

\begin{theorem}\label{thm:e1-mixed-products}
For \(\varphi_1,\ldots,\varphi_n\in\mathcal E^1(X,\omega)\), the mixed
measure
\[
       \omega_{\varphi_1}\wedge\cdots\wedge\omega_{\varphi_n}
\]
is defined by common truncation and has mass \(V\). It is symmetric,
multilinear under psh convex combinations, does not charge pluripolar sets,
and is continuous when each entry decreases to an \(\mathcal E^1\) limit.
The maximum and comparison principles extend from bounded potentials to
\(\mathcal E^1\).
\end{theorem}

\begin{proof}
Choose one common constant \(C\) such that \(\varphi_i-C\leq0\) for every
\(i\), and replace the whole tuple by this common translate.  Keep the
notation \(U_t,B_t,\bar\varphi\) of
\Cref{lem:ch26-mixed-tail} and write
\(\nu_t:=\MA_\omega(\varphi_1^{\langle t\rangle},\ldots,
\varphi_n^{\langle t\rangle})\), a positive measure of mass \(V\) supplied by
\Cref{thm:bounded-ma-extension}.

\emph{Step 1: the truncated measures are consistent on the layers.}  For
\(t\leq s\) the bounded \(\omega\)-psh functions
\(\varphi_i^{\langle t\rangle}\) and
\(\varphi_i^{\langle s\rangle}\) agree on the Borel layer \(U_t\), for
every \(i\), but ordinary open-set locality cannot be invoked.  Instead,
replace the entries one at a time.  Since
\[
 \varphi_i^{\langle t\rangle}
 =\max\{\varphi_i^{\langle s\rangle},-t\},
 \qquad
 \{\varphi_i^{\langle s\rangle}>-t\}=\{\varphi_i>-t\},
\]
the mixed maximum identity
\eqref{eq:ch26-mixed-maximum-principle}, restricted further to
\(U_t\subseteq\{\varphi_i>-t\}\), shows that replacing the \(i\)-th entry
does not change the measure on \(U_t\).  Iterating over \(i=1,\ldots,n\)
gives
\begin{equation}
\label{eq:ch26-mixed-layer-consistency}
 \one_{U_t}\nu_s=\one_{U_t}\nu_t .
\end{equation}

\emph{Step 2: the increasing limit.}  Exactly as in
\Cref{lem:ch26-layers}, \eqref{eq:ch26-mixed-layer-consistency} makes
\(t\mapsto\nu_t(B\cap U_t)\) nondecreasing for every Borel \(B\), so
\begin{equation}
\label{eq:ch26-mixed-definition}
 \bigl(\omega_{\varphi_1}\wedge\cdots\wedge\omega_{\varphi_n}\bigr)(B)
 :=\lim_{t\to\infty}\nu_t(B\cap U_t)
\end{equation}
defines a positive Borel measure, and restricting the limit to
\(t\in\mathbf N\) --- legitimate by monotonicity --- gives countable
additivity as before.  The compact-approximation argument in the proof of
\Cref{lem:ch26-layers}, applied to the increasing sequence
\(\one_{U_m}\nu_m\), shows that this finite Borel measure is Radon.  We write
\(\MA_\omega(\varphi_1,\ldots,\varphi_n)\) for it.  The choice of the common
normalizing constant is harmless: for \(c\geq0\),
\[
 (\varphi_i-c)^{\langle t+c\rangle}
   =\varphi_i^{\langle t\rangle}-c,
 \qquad
 \{\varphi_i-c>-(t+c)\}=\{\varphi_i>-t\},
\]
simultaneously for every \(i\).  Bounded mixed products are unchanged by
adding constants to their entries, so the two definitions run through the
same cofinal family of restricted measures.  Invariance under
\emph{separate} constants will follow after monotone continuity, in Step~6.

\emph{Step 3: total mass \(V\).}  By \eqref{eq:ch26-mixed-definition} and
\(\nu_t(X^{\An})=V\),
\[
 \MA_\omega(\varphi_1,\ldots,\varphi_n)(X^{\An})
 =\lim_{t\to\infty}\bigl(V-\nu_t(B_t)\bigr)=V
\]
by \Cref{lem:ch26-mixed-tail}. Finiteness of energy gives this vanishing tail and also the uniform tail bound used in Step~5 below. For arbitrary psh entries the limit
in \eqref{eq:ch26-mixed-definition} still exists, but its mass may be
strictly smaller than \(V\), and may be zero
(\Cref{ex:ch26-infinite-energy}).

\emph{Step 4: symmetry, polar sets, and locality.}  Symmetry is inherited
from \(\nu_t\).  Pluripolar sets carry no mass because on each layer the
measure is the restriction of a bounded mixed product, to which
\Cref{prop:bounded-ma-no-polar-mass} applies.  Locality on open sets follows
from \Cref{lem:ch26-mixed-locality} layer by layer.

\emph{Step 5: monotone continuity.}  Let \(\varphi_i^\alpha\searrow\varphi_i\)
be decreasing nets in \(\cE^1\), all \(\leq0\), with
\(\varphi_i\in\cE^1\).  Put
\(\bar\varphi^\alpha=\frac1n\sum_i\varphi_i^\alpha\geq\bar\varphi\), so that
\(-E_\omega(\bar\varphi^\alpha)\leq-E_\omega(\bar\varphi)=:C_0<\infty\) and
the bound \eqref{eq:ch26-mixed-tail} is uniform in \(\alpha\).  Let
\(f\in C^0(X^{\An})\) with \(\lVert f\rVert_\infty\leq1\).  For every
\(t>0\),
\[
\begin{split}
 \Bigl|\int f\,\MA_\omega(\varphi_1^\alpha,\ldots,\varphi_n^\alpha)
      -\int f\,\MA_\omega(\varphi_1,\ldots,\varphi_n)\Bigr|
 \leq{}&\Bigl|\int f\,\nu_t^\alpha-\int f\,\nu_t\Bigr|\\
 &+2\nu_t^\alpha(B_t^\alpha)+2\nu_t(B_t),
\end{split}
\]
where \(\nu^\alpha_t\) and \(B_t^\alpha\) are formed from the
\(\varphi_i^\alpha\).  Here we used that, by
\eqref{eq:ch26-mixed-definition}, the measures
\(\MA_\omega(\varphi_1,\ldots,\varphi_n)\) and \(\nu_t\) agree on \(U_t\) and
have the same total mass \(V\), so that for a bounded Borel \(f\)
\[
 \begin{aligned}
 &\Bigl|\int f\MA_\omega(\varphi_1,\ldots,\varphi_n)-\int f\,\nu_t\Bigr|\\
 &\qquad\leq\lVert f\rVert_\infty
   \bigl(\MA_\omega(\varphi_1,\ldots,\varphi_n)(B_t)+\nu_t(B_t)\bigr)\\
 &\qquad\leq2\lVert f\rVert_\infty\,\nu_t(B_t),
 \end{aligned}
\]
and similarly for the net.  The last two terms are therefore at most
\(4\kappa_nC_0/t\) with
\(\kappa_n=n^{n+1}(n+1)/n!\), uniformly in \(\alpha\).  For fixed \(t\), the
truncations \((\varphi_i^\alpha)^{\langle t\rangle}\) form decreasing nets of
bounded \(\omega\)-psh functions with limits
\(\varphi_i^{\langle t\rangle}\), so
\eqref{eq:weighted-monotone-ma} gives \(\int f\nu^\alpha_t\to\int f\nu_t\)
after approximating \(f\) uniformly by model functions, which is legitimate
because all measures involved have mass \(V\).  Choosing \(t\) large and then
\(\alpha\) large makes the total error arbitrarily small.  Hence
\(\MA_\omega(\varphi_1^\alpha,\ldots,\varphi_n^\alpha)\to
\MA_\omega(\varphi_1,\ldots,\varphi_n)\) weakly.

\emph{Step 6: separate constants and multilinearity.}  First let
\(c_1,\ldots,c_n\in\mathbf R\).  The bounded potentials
\(\varphi_i^{\langle t\rangle}+c_i\) decrease to \(\varphi_i+c_i\), and
at every bounded level
\[
\MA_\omega(\varphi_1^{\langle t\rangle}+c_1,\ldots,
                  \varphi_n^{\langle t\rangle}+c_n)
 =\MA_\omega(\varphi_1^{\langle t\rangle},\ldots,
                  \varphi_n^{\langle t\rangle}).
\]
After one further common translation, all entries of both tuples are
nonpositive, so Step~5 applies to both sides and shows that the
finite-energy mixed product is unchanged when arbitrary constants are added
separately to its entries.

Now let
\(\varphi_1=(1-\lambda)\chi+\lambda\chi'\), where
\(\chi,\chi'\in\cE^1(X,\omega)\) and \(0\leq\lambda\leq1\), and fix
\(\varphi_2,\ldots,\varphi_n\in\cE^1(X,\omega)\).  After translating the
potentials, assume that they are all nonpositive.  For \(t>0\) put
\[
 q_t:=(1-\lambda)\chi^{\langle t\rangle}
             +\lambda(\chi')^{\langle t\rangle},
 \qquad
 q_{j,t}:=\varphi_j^{\langle t\rangle}\quad(2\leq j\leq n).
\]
Then \(q_t\searrow\varphi_1\) and
\(q_{j,t}\searrow\varphi_j\).  At bounded level,
\Cref{cor:bounded-ma-multilinearity} gives
\[
 \begin{aligned}
 \MA_\omega(q_t,q_{2,t},\ldots,q_{n,t})
  ={}&(1-\lambda)
     \MA_\omega(\chi^{\langle t\rangle},q_{2,t},\ldots,q_{n,t})\\
   &+\lambda
     \MA_\omega((\chi')^{\langle t\rangle},q_{2,t},\ldots,q_{n,t}).
 \end{aligned}
\]
Apply the monotone continuity established in Step~5 to all three terms and
let \(t\to\infty\).  The resulting identity is affineness in the first
entry.  Symmetry gives the same conclusion in every entry, which is the
asserted multilinearity under psh convex combinations.  This argument does
not require a strict superlevel set of an upper semicontinuous function to
be open.

\emph{Step 7: maximum and comparison principles.}  These are
\Cref{thm:ch26-e1-maximum-principle} and
\Cref{cor:ch26-e1-comparison} below, proved by truncation from the bounded
statements \eqref{eq:bounded-maximum-principle} and
\eqref{eq:bounded-comparison-principle}.
\end{proof}

\begin{theorem}
\label{thm:ch26-e1-maximum-principle}
Let \(\varphi,\psi\in\cE^1(X,\omega)\).  Then
\(\max\{\varphi,\psi\}\in\cE^1(X,\omega)\) and
\begin{equation}
\label{eq:ch26-e1-maximum-principle}
 \one_{\{\varphi>\psi\}}\MA_\omega(\max\{\varphi,\psi\})
 =\one_{\{\varphi>\psi\}}\MA_\omega(\varphi),
\end{equation}
for the finite-energy product of \Cref{thm:e1-mixed-products}.  The same
identity holds for mixed products in which the remaining \(n-1\) entries are
arbitrary members of \(\cE^1(X,\omega)\).
\end{theorem}

\begin{proof}
Membership is \Cref{thm:ch26-e1-structure}\,\cref{itm:ch26-e1:max}.
Normalize \(\varphi,\psi\leq0\) and let \(\chi:=\max\{\varphi,\psi\}\), so
that \(\chi^{\langle t\rangle}=\max\{\varphi^{\langle
t\rangle},\psi^{\langle t\rangle}\}\).  Fix \(t>0\) and put
\(W_t:=\{\varphi>-t\}\cap\{\psi>-t\}\), a Borel layer.  On \(W_t\) the
inequality \(\varphi>\psi\) is equivalent to
\(\varphi^{\langle t\rangle}>\psi^{\langle t\rangle}\), so
\eqref{eq:bounded-maximum-principle} gives
\[
 \one_{W_t\cap\{\varphi>\psi\}}\MA_\omega(\chi^{\langle t\rangle})
 =\one_{W_t\cap\{\varphi>\psi\}}\MA_\omega(\varphi^{\langle t\rangle}).
\]
Since \(\chi\geq\varphi\), one has \(\{\chi\leq-t\}\subseteq\{\varphi\leq
-t\}\), so \(W_t\) is contained in the layer of \(\chi\) as well.  Passing to
the increasing limit in \(t\) through \eqref{eq:ch26-npp-definition} on both
sides yields \eqref{eq:ch26-e1-maximum-principle}: the left-hand sides
increase to \(\one_{\{\varphi>\psi\}}\MA_\omega(\chi)\) and the right-hand
sides to \(\one_{\{\varphi>\psi\}}\MA_\omega(\varphi)\), because
\(\bigcup_tW_t=\{\varphi>-\infty\}\cap\{\psi>-\infty\}\) carries both
measures --- for \(\MA_\omega(\varphi)\) on the set \(\{\varphi>\psi\}\) this
uses that \(\{\psi=-\infty\}\) is pluripolar and hence
\(\MA_\omega(\varphi)\)-null, by
\Cref{prop:ch26-npp-basic}\,\cref{itm:ch26-npp:polar}.  The mixed statement is
identical, using \eqref{eq:ch26-mixed-maximum-principle} in place of
\eqref{eq:bounded-maximum-principle}.
\end{proof}

\begin{corollary}
\label{cor:ch26-e1-comparison}
For \(\varphi,\psi\in\cE^1(X,\omega)\),
\begin{equation}
\label{eq:ch26-e1-comparison}
 \int_{\{\varphi<\psi\}}\MA_\omega(\psi)
 \leq\int_{\{\varphi<\psi\}}\MA_\omega(\varphi).
\end{equation}
\end{corollary}

\begin{proof}
Fix \(\varepsilon>0\) and put \(\chi:=\max\{\varphi,\psi-\varepsilon\}\in\cE^1\).
By \Cref{thm:ch26-e1-maximum-principle} applied twice,
\[
 \begin{aligned}
 \one_{\{\varphi>\psi-\varepsilon\}}\MA_\omega(\chi)
 &=\one_{\{\varphi>\psi-\varepsilon\}}\MA_\omega(\varphi),\\
 \one_{\{\varphi<\psi-\varepsilon\}}\MA_\omega(\chi)
 &=\one_{\{\varphi<\psi-\varepsilon\}}\MA_\omega(\psi).
 \end{aligned}
\]
All three measures have total mass \(V\) by
\Cref{thm:e1-mixed-products}.  Hence
\[
\begin{split}
 V&=\MA_\omega(\chi)(\{\varphi<\psi-\varepsilon\})
   +\MA_\omega(\chi)(\{\varphi\geq\psi-\varepsilon\})\\
  &\geq\MA_\omega(\psi)(\{\varphi<\psi-\varepsilon\})
   +\MA_\omega(\varphi)(\{\varphi>\psi-\varepsilon\})\\
  &=\MA_\omega(\psi)(\{\varphi<\psi-\varepsilon\})
   +V-\MA_\omega(\varphi)(\{\varphi\leq\psi-\varepsilon\}),
\end{split}
\]
which rearranges to
\(\MA_\omega(\psi)(\{\varphi<\psi-\varepsilon\})\leq
\MA_\omega(\varphi)(\{\varphi\leq\psi-\varepsilon\})\).  Letting
\(\varepsilon\downarrow0\), the left-hand sets increase to
\(\{\varphi<\psi\}\) and the right-hand sets also increase, to
\(\bigcup_{\varepsilon>0}\{\varphi\leq\psi-\varepsilon\}
=\{\varphi<\psi\}\cup\{\varphi=\psi=-\infty\}\); both families are
increasing as \(\varepsilon\downarrow0\), since the constraint weakens.
The extra set is
pluripolar, hence \(\MA_\omega(\varphi)\)-null by
\Cref{prop:ch26-npp-basic}\,\cref{itm:ch26-npp:polar}.  Continuity from below
on both sides, both measures being finite, gives
\eqref{eq:ch26-e1-comparison}.
\end{proof}

\begin{proposition}
\label{prop:ch26-e1-cocycle}
Let \(\varphi,\psi\in\cE^1(X,\omega)\) with \(\varphi\leq\psi\), and for
\(0\leq j\leq n\) put
\[
 I_j:=\int(\psi-\varphi)\,
   \omega_\varphi^{\,j}\wedge\omega_\psi^{\,n-j}\in[0,+\infty],
\]
the products being the finite-energy products of
\Cref{thm:e1-mixed-products}.  Then every \(I_j\) is well defined, the
sandwich
\begin{equation}
\label{eq:ch26-e1-sandwich}
 I_0=\int(\psi-\varphi)\MA_\omega(\psi)
 \ \leq\ E_\omega(\psi)-E_\omega(\varphi)\ \leq\
 \int(\psi-\varphi)\MA_\omega(\varphi)=I_n
\end{equation}
holds, and
\begin{equation}
\label{eq:ch26-e1-cocycle}
 \frac1{n+1}\sum_{j=0}^nI_j
 \ \leq\ E_\omega(\psi)-E_\omega(\varphi).
\end{equation}
In particular every \(I_j\) is finite, and
\(\int(-\varphi)\MA_\omega(\psi)<+\infty\).
\end{proposition}

\begin{proof}
Normalize \(\psi\leq0\), so \(\varphi\leq\psi\leq0\).  The difference
\(\psi-\varphi\) is defined and nonnegative on
\(\{\varphi>-\infty\}\), which carries all the measures in question by
\Cref{prop:ch26-npp-basic}\,\cref{itm:ch26-npp:polar}; on the complement it is
of the forbidden form \((-\infty)-(-\infty)\) and is left undefined, which is
harmless by \Cref{lem:integration-ledger}\,(3).  Each \(I_j\) is therefore
defined in \([0,+\infty]\).  Throughout, \(U_t:=\{\varphi>-t\}\); since
\(\varphi\leq\psi\), this is exactly the layer
\(\{\varphi>-t\}\cap\{\psi>-t\}\) associated with any \(n\)-tuple whose
entries lie in \(\{\varphi,\psi\}\), so
\eqref{eq:ch26-mixed-layer-consistency} identifies the restriction of the
truncated bounded mixed measure to \(U_t\) with that of the finite-energy
mixed measure.

\emph{Step 1: the upper bound in \eqref{eq:ch26-e1-sandwich}.}  Apply the
supergradient inequality \eqref{eq:ch26-supergradient} with the bounded
potential \(\varphi^{\langle t\rangle}\) and the arbitrary potential
\(\psi\):
\[
 E_\omega(\psi)
 \leq E_\omega(\varphi^{\langle t\rangle})
   +\int\bigl(\psi-\varphi^{\langle t\rangle}\bigr)
     \MA_\omega(\varphi^{\langle t\rangle}).
\]
Write \(\mu_t=\MA_\omega(\varphi^{\langle t\rangle})\) and split the integral
over \(U_t\) and its complement.  On \(U_t\) the integrand is
\(\psi-\varphi\geq0\) and \(\mu_t\) restricts to \(\MA_\omega(\varphi)\), so
that part increases to \(I_n\) by monotone convergence.  On the complement
\(\{\varphi\leq-t\}\) the integrand is \(\psi+t\leq t\), so that part is at
most \(t\,\mu_t(\{\varphi\leq-t\})\), which tends to \(0\) by
\eqref{eq:e1-mass-tail}.  Letting \(t\to\infty\) and using
\(E_\omega(\varphi^{\langle t\rangle})\to E_\omega(\varphi)\)
(\Cref{prop:ch26-truncation-limit}) gives
\(E_\omega(\psi)-E_\omega(\varphi)\leq I_n\).

\emph{Step 2: the lower bound in \eqref{eq:ch26-e1-sandwich}.}  Apply
\eqref{eq:ch26-supergradient} with the bounded potential
\(\psi^{\langle t\rangle}\) and the arbitrary potential \(\varphi\):
\[
 E_\omega(\varphi)
 \leq E_\omega(\psi^{\langle t\rangle})
   +\int\bigl(\varphi-\psi^{\langle t\rangle}\bigr)
     \MA_\omega(\psi^{\langle t\rangle}).
\]
Write \(\nu_t=\MA_\omega(\psi^{\langle t\rangle})\) and split over
\(W_t:=\{\psi>-t\}\) and its complement.  On \(W_t\) the integrand is
\(\varphi-\psi\leq0\) and \(\nu_t\) restricts to \(\MA_\omega(\psi)\), so that
part decreases to \(-I_0\) by monotone convergence.  On the complement
\(\{\psi\leq-t\}\) one has \(\varphi\leq\psi\leq-t\), so the integrand
\(\varphi+t\) is \(\leq0\) and that part is nonpositive.  Hence
\[
 E_\omega(\varphi)\leq E_\omega(\psi^{\langle t\rangle})
   +\int_{W_t}(\varphi-\psi)\MA_\omega(\psi),
\]
and letting \(t\to\infty\) gives
\(E_\omega(\varphi)\leq E_\omega(\psi)-I_0\), which is the lower bound.  In
particular \(I_0<+\infty\), and since
\(\int(-\psi)\MA_\omega(\psi)<+\infty\) by
\eqref{eq:e1-integrability}, also
\(\int(-\varphi)\MA_\omega(\psi)=I_0+\int(-\psi)\MA_\omega(\psi)<+\infty\).

\emph{Step 3: the cocycle inequality.}  Apply \eqref{eq:energy-cocycle} in
the bounded form of \Cref{prop:bounded-energy-identities} to the pair
\(\varphi^{\langle t\rangle}\leq\psi^{\langle t\rangle}\):
\begin{equation}
\label{eq:ch26-truncated-cocycle}
 E_\omega(\psi^{\langle t\rangle})-E_\omega(\varphi^{\langle t\rangle})
 =\frac1{n+1}\sum_{j=0}^n
   \int\bigl(\psi^{\langle t\rangle}-\varphi^{\langle t\rangle}\bigr)\,
   \omega_{\varphi^{\langle t\rangle}}^{\,j}\wedge
   \omega_{\psi^{\langle t\rangle}}^{\,n-j}.
\end{equation}
The left-hand side converges to \(E_\omega(\psi)-E_\omega(\varphi)\), a
difference of two finite numbers.  Every summand on the right is nonnegative,
and restricting the \(j\)-th one to \(U_t\) --- where the integrand equals
\(\psi-\varphi\) and the measure equals the finite-energy mixed measure ---
gives the lower bound
\[
 \int\bigl(\psi^{\langle t\rangle}-\varphi^{\langle t\rangle}\bigr)
  \omega_{\varphi^{\langle t\rangle}}^{\,j}\wedge
  \omega_{\psi^{\langle t\rangle}}^{\,n-j}
 \ \geq\
 \int_{U_t}(\psi-\varphi)\,
  \omega_\varphi^{\,j}\wedge\omega_\psi^{\,n-j},
\]
whose right-hand side increases to \(I_j\).  Taking the lower limit as
\(t\to\infty\) in \eqref{eq:ch26-truncated-cocycle} and using that a finite
sum of lower limits is at most the lower limit of the sum yields
\eqref{eq:ch26-e1-cocycle}.  Finiteness of each \(I_j\) follows.
\end{proof}

\begin{remark}[Why only an inequality]
\label{rem:ch26-cocycle-inequality}
At bounded level \eqref{eq:ch26-truncated-cocycle} is an identity, and one
would like to pass to the limit in each summand separately.  The obstruction
is the contribution of the complement \(\{\varphi\leq-t\}\) to the \(j\)-th
mixed integral: the integrand there is at most \(t\), and
\eqref{eq:ch26-mixed-tail} bounds the mixed mass of that set only by
\(O(t^{-1})\), so the product is \(O(1)\) and need not tend to zero.  For the
extreme summands \(j=0\) and \(j=n\) the difficulty disappears --- in Step~1
the tail is controlled by the sharper diagonal estimate
\eqref{eq:e1-mass-tail}, which is \(o(t^{-1})\), and in Step~2 the tail
contribution has a favourable sign --- and this is why the sandwich
\eqref{eq:ch26-e1-sandwich} is an equality-free two-sided bound rather than a
consequence of a term-by-term limit.  Everything used later in this book
follows from \eqref{eq:ch26-e1-sandwich} and
\eqref{eq:ch26-e1-cocycle}, so we do not pursue the sharper mixed tail
estimate that would upgrade \eqref{eq:ch26-e1-cocycle} to an identity.
\end{remark}

\begin{warning}{The cocycle statements for unordered pairs}
\Cref{prop:ch26-e1-cocycle} is stated for \(\varphi\leq\psi\), and the
hypothesis is not decorative: for an unordered pair the pointwise difference
\(\psi-\varphi\) is of the forbidden shape \((-\infty)-(-\infty)\) on the
common pole set, and the integrals \(I_j\) need not be defined.  The
permitted reduction is to interpose \(\chi:=\max\{\varphi,\psi\}\in\cE^1\)
and apply the ordered case to the two pairs \(\varphi\leq\chi\) and
\(\psi\leq\chi\), combining them by the chain rule
\eqref{eq:ch26-cocycle-chain}:
\[
 E_\omega(\psi)-E_\omega(\varphi)
 =\bigl(E_\omega(\chi)-E_\omega(\varphi)\bigr)
  -\bigl(E_\omega(\chi)-E_\omega(\psi)\bigr),
\]
which is legitimate because all three energies are finite real numbers.  This
is the reduction used in \Cref{sec:ch26-e1-gauge} and in
\Cref{thm:e1-energy-differentiability}; it is not an additional theorem, but
it is the only reading permitted outside the ordered case.
\end{warning}

\section{The class \texorpdfstring{\(\cE^1\)}{E1} and its gauges}
\label{sec:ch26-e1-gauge}

\Cref{thm:ch26-e1-structure} recorded the order structure of \(\cE^1\).  We
now ask what quantitative structure the energy puts on it.  The answer is
more delicate than the complex-geometric literature might suggest, and the
delicacy is not an artefact of the non-Archimedean setting: it already
appears on a metric graph with two vertices.

\subsection{The energy gauge}
\label{sec:ch26-gauge}

\begin{definition}[The energy gauge]
\label{def:ch26-gauge}
\index{energy gauge}
For \(\varphi,\psi\in\cE^1(X,\omega)\) set
\begin{equation}
\label{eq:ch26-gauge}
 d_\omega(\varphi,\psi)
 :=2E_\omega\bigl(\max\{\varphi,\psi\}\bigr)
   -E_\omega(\varphi)-E_\omega(\psi).
\end{equation}
\end{definition}

All three terms are finite, by
\Cref{thm:ch26-e1-structure}\,\cref{itm:ch26-e1:max}, so
\eqref{eq:ch26-gauge} is a well-defined real number; no forbidden
cancellation occurs.

\begin{proposition}
\label{prop:ch26-gauge-elementary}
For \(\varphi,\psi,\chi\in\cE^1(X,\omega)\) and \(c\in\mathbf R\):
\begin{enumerate}
\item\label{itm:ch26-gauge:pos}
 \(d_\omega(\varphi,\psi)\geq0\), with \(d_\omega(\varphi,\varphi)=0\);
\item\label{itm:ch26-gauge:sym}
 \(d_\omega(\varphi,\psi)=d_\omega(\psi,\varphi)\);
\item\label{itm:ch26-gauge:transl}
 \(d_\omega(\varphi+c,\psi+c)=d_\omega(\varphi,\psi)\);
\item\label{itm:ch26-gauge:ordered}
 if \(\varphi\leq\psi\), then
 \(d_\omega(\varphi,\psi)=E_\omega(\psi)-E_\omega(\varphi)\);
\item\label{itm:ch26-gauge:dominated}
 if \(\varphi\leq\chi\) and \(\psi\leq\chi\), then
 \(d_\omega(\varphi,\psi)\leq d_\omega(\varphi,\chi)+d_\omega(\chi,\psi)\);
\item\label{itm:ch26-gauge:chain}
 if \(\varphi\leq\chi\leq\psi\), then
 \(d_\omega(\varphi,\psi)=d_\omega(\varphi,\chi)+d_\omega(\chi,\psi)\).
\end{enumerate}
\end{proposition}

\begin{proof}
\cref{itm:ch26-gauge:pos} and \cref{itm:ch26-gauge:sym} are immediate from
monotonicity of \(E_\omega\) and \(\max\{\varphi,\psi\}\geq\varphi,\psi\).
\cref{itm:ch26-gauge:transl} follows from
\(\max\{\varphi+c,\psi+c\}=\max\{\varphi,\psi\}+c\) and
\eqref{eq:bfj-energy-translation}: the three terms change by \(2cV\),
\(-cV\), \(-cV\).  \cref{itm:ch26-gauge:ordered} is the case
\(\max\{\varphi,\psi\}=\psi\).  For \cref{itm:ch26-gauge:dominated}, use
\cref{itm:ch26-gauge:ordered} twice and monotonicity:
\[
 \begin{aligned}
 d_\omega(\varphi,\chi)+d_\omega(\chi,\psi)
 &=2E_\omega(\chi)-E_\omega(\varphi)-E_\omega(\psi)\\
 &\geq 2E_\omega(\max\{\varphi,\psi\})\\
 &\quad-E_\omega(\varphi)-E_\omega(\psi),
 \end{aligned}
\]
since \(\max\{\varphi,\psi\}\leq\chi\).  \cref{itm:ch26-gauge:chain} is
\cref{itm:ch26-gauge:ordered} applied three times, the three energies
telescoping.
\end{proof}

\subsection{Nondegeneracy}
\label{sec:ch26-gauge-nondegenerate}

\begin{proposition}
\label{prop:ch26-strict-monotone}
Let \(\varphi\leq\psi\) in \(\cE^1(X,\omega)\) with
\(E_\omega(\varphi)=E_\omega(\psi)\).  Then:
\begin{enumerate}
\item\label{itm:ch26-strict:ae}
 \(\varphi=\psi\) almost everywhere with respect to each of the \(n+1\)
 measures \(\omega_\varphi^{\,j}\wedge\omega_\psi^{\,n-j}\),
 \(0\leq j\leq n\);
\item\label{itm:ch26-strict:bounded}
 if moreover \(\varphi\) and \(\psi\) are bounded, then all \(n+1\) of those
 measures coincide, and in particular
 \(\MA_\omega(\varphi)=\MA_\omega(\psi)\);
\item\label{itm:ch26-strict:cont}
 if in addition \(\psi\) is continuous, then \(\varphi=\psi\).
\end{enumerate}
\end{proposition}

\begin{proof}
\cref{itm:ch26-strict:ae}.  By \eqref{eq:ch26-e1-cocycle} the \(n+1\)
nonnegative numbers \(I_j\) satisfy
\(\sum_jI_j\leq(n+1)\bigl(E_\omega(\psi)-E_\omega(\varphi)\bigr)=0\), so each
vanishes.  Since the integrand \(\psi-\varphi\) is nonnegative, this says
exactly that \(\psi=\varphi\) almost everywhere for each of the \(n+1\)
measures.

\cref{itm:ch26-strict:bounded}.  Assume \(\varphi,\psi\) bounded and put
\(u=\varphi-\psi\).  By \Cref{thm:ch26-mixed-energy-identities}\,%
\cref{itm:ch26-ij:J} and \eqref{eq:ch26-ak-telescope}, vanishing of all the
integrals forces \(J_\psi(\varphi)=0\), hence \(a_k(\varphi,\psi)=0\) for
every \(k\) by \Cref{thm:ch26-mixed-energy-identities}\,%
\cref{itm:ch26-ij:positive}.  Let \(h\) be any model function.  By
\Cref{rem:ch26-omega-theta} we may assume that \(\omega\) is
\(\mathcal X\)-positive on a model determining \(h\); then
\(\omega+C^{-1}\ddc h\) is relatively ample for \(C\) large, so \(C^{-1}h\)
is an \(\omega\)-psh model function and \(h=C\cdot(C^{-1}h-0)\) is a positive
multiple of a difference of two bounded \(\omega\)-psh functions.  The
Cauchy--Schwarz inequality \eqref{eq:ch26-cauchy-schwarz} with \(k=n-1\),
applied to \(C^{-1}h\) and rescaled, gives
\[
 \Bigl|\int h\,\ddc u\wedge\omega_\varphi^{\,n-1}\Bigr|^2
 \leq C^2\,a_{n-1}(\varphi,\psi)\cdot
  \Bigl(-\int C^{-1}h\,\ddc(C^{-1}h)\wedge\omega_\varphi^{\,n-1}\Bigr)=0 .
\]
Now \(\ddc u\wedge\omega_\varphi^{\,n-1}
=\MA_\omega(\varphi)-\MA_\omega(\varphi^{(n-1)},\psi)\) is a signed Radon
measure integrating to zero against every model function, and model functions
are uniformly dense in \(C^0(X^{\An})\); hence
\(\MA_\omega(\varphi)=\MA_\omega(\varphi^{(n-1)},\psi)\).  Repeating with
\(k=n-2,\ldots,0\) identifies all \(n+1\) mixed measures with one another.

\cref{itm:ch26-strict:cont}.  Suppose \(\psi\) is continuous, hence bounded.
Assume first that \(\varphi\) is bounded as well.  By
\cref{itm:ch26-strict:ae} with \(j=n\) we have \(\psi=\varphi\) almost
everywhere for \(\MA_\omega(\varphi)\); in particular
\(\psi\leq\varphi\) almost everywhere for \(\MA_\omega(\varphi)\).  After
subtracting a constant we may assume \(\varphi\leq0\).  Then
\Cref{prop:bounded-capacity-domination}, applied with the pair
\((\psi,\varphi)\) in place of the pair \((\varphi,\psi)\) occurring there,
gives \(\psi\leq\varphi\) everywhere, hence \(\varphi=\psi\).

For unbounded \(\varphi\), choose \(t\geq-\inf_{X^{\An}}\psi\).  Then
\[
 \varphi^{\langle t\rangle}=\max\{\varphi,-t\}
 \leq\max\{\psi,-t\}=\psi .
\]
Monotonicity gives
\[
 E_\omega(\varphi)\leq E_\omega(\varphi^{\langle t\rangle})
 \leq E_\omega(\psi)=E_\omega(\varphi),
\]
so \(E_\omega(\varphi^{\langle t\rangle})=E_\omega(\psi)\).  The bounded case
applies to the pair \(\varphi^{\langle t\rangle}\leq\psi\) and yields
\(\varphi^{\langle t\rangle}=\psi\) for every such \(t\).  Letting
\(t\to\infty\) gives \(\varphi=\psi\); in particular \(\varphi\) was bounded
after all.
\end{proof}

\begin{corollary}
\label{cor:ch26-gauge-nondegenerate}
Let \(\varphi,\psi\in\cE^1(X,\omega)\) with
\(d_\omega(\varphi,\psi)=0\), and put \(\chi:=\max\{\varphi,\psi\}\).  Then:
\begin{enumerate}
\item\label{itm:ch26-nondeg:ae}
 \(\chi=\varphi\) almost everywhere for \(\MA_\omega(\varphi)\), and
 \(\chi=\psi\) almost everywhere for \(\MA_\omega(\psi)\); in particular
 \(\psi\leq\varphi\) almost everywhere for \(\MA_\omega(\varphi)\) and
 \(\varphi\leq\psi\) almost everywhere for \(\MA_\omega(\psi)\);
\item\label{itm:ch26-nondeg:one}
 if \(\varphi\) and \(\psi\) are bounded and \(\psi\) is continuous, then
 \(\psi\leq\varphi\) everywhere;
\item\label{itm:ch26-nondeg:both}
 if \(\varphi\) and \(\psi\) are both continuous, then \(\varphi=\psi\).
\end{enumerate}
\end{corollary}

\begin{proof}
Vanishing of \(d_\omega\) forces
\(E_\omega(\chi)=E_\omega(\varphi)=E_\omega(\psi)\): indeed
\(d_\omega(\varphi,\psi)=(E_\omega(\chi)-E_\omega(\varphi))
+(E_\omega(\chi)-E_\omega(\psi))\) is a sum of two nonnegative terms.

\cref{itm:ch26-nondeg:ae}.  Apply
\Cref{prop:ch26-strict-monotone}\,\cref{itm:ch26-strict:ae} to the ordered
pair \(\varphi\leq\chi\) with \(j=n\): the integral
\(\int(\chi-\varphi)\MA_\omega(\varphi)\) vanishes, so \(\chi=\varphi\)
almost everywhere for \(\MA_\omega(\varphi)\); since \(\chi\geq\psi\), this
gives \(\psi\leq\varphi\) almost everywhere for that measure.  Exchange
\(\varphi\) and \(\psi\).

\cref{itm:ch26-nondeg:one}.  By \cref{itm:ch26-nondeg:ae},
\(\psi\leq\varphi\) almost everywhere for \(\MA_\omega(\varphi)\).  After
subtracting a constant, \(\varphi\leq0\); \(\psi\) is continuous, so
\Cref{prop:bounded-capacity-domination}, applied with the pair
\((\psi,\varphi)\) in place of \((\varphi,\psi)\) there, gives
\(\psi\leq\varphi\) everywhere.

\cref{itm:ch26-nondeg:both}.  Apply \cref{itm:ch26-nondeg:one} twice.
\end{proof}

\begin{warning}{What has and has not been proved about nondegeneracy}
\Cref{cor:ch26-gauge-nondegenerate} proves \(d_\omega(\varphi,\psi)=0
\Rightarrow\varphi=\psi\) only when both potentials are continuous, and the
one-sided conclusion \cref{itm:ch26-nondeg:one} only when they are bounded.
For a single continuous competitor the argument stops after one inequality,
because the domination principle proved in this book,
\Cref{prop:bounded-capacity-domination}, requires the \emph{smaller}
function to be continuous, and continuity of \(\psi\) is available on only
one side.  What would close the gap is uniqueness of solutions of the
Monge--Amp\`ere equation modulo constants; that is the subject of
\Cref{chap:bfj-continuity}, and it is not available at this point of the
development.  The trivially valued analogue, where the envelope property is
available in full and a genuine metric exists, is
\Cref{chap:bj-strong-topologies}; see also \cite[\S12]{BJ22}.
\end{warning}

\subsection{The triangle inequality fails}
\label{sec:ch26-triangle}

It is tempting to regard \eqref{eq:ch26-gauge} as a distance.  The following
computation, carried out entirely on the two-vertex metric graph of
\Cref{ex:two-vertex-ma}, shows that it is not.

\begin{example}[The energy gauge is not a metric]
\label{ex:ch26-triangle-fails}
Let \(X\) be a smooth projective curve with a regular semistable model whose
dual complex is the segment \(\Gamma=[x_0,x_1]\) of length one, and let
\(\omega=\delta_{x_0}+\delta_{x_1}\), so \(n=1\) and \(V=2\).  Parametrize
\(\Gamma\) by \(t\in[0,1]\).  By
\Cref{prop:ch21-curve-criterion}, a piecewise affine function \(u\) on
\(\Gamma\) is (the restriction of) an \(\omega\)-psh model function precisely
when it is convex on the open edge and its outgoing slopes at the two
endpoints are at least \(-1\); after a subdivision of the model, breakpoints
at rational interior parameters are permitted.  By
\Cref{prop:ch26-graph-energy} below,
\begin{equation}
\label{eq:ch26-triangle-energy-formula}
 E_\omega(u)=u(0)+u(1)-\tfrac12\int_0^1|u'|^2\,\rd t .
\end{equation}
Take
\[
 \varphi(t)=1-t,\qquad
 \chi(t)=t,\qquad
 \psi(t)=0 .
\]
Each is affine with slope in \([-1,1]\), hence \(\omega\)-psh, and
\(\max\{\varphi,\chi\}(t)=\max\{1-t,t\}\) is convex with slopes \(\mp1\) and
a breakpoint at \(t=\tfrac12\), hence \(\omega\)-psh on the model obtained by
one further blowup.  Formula \eqref{eq:ch26-triangle-energy-formula} gives
\[
 E_\omega(\varphi)=1+0-\tfrac12=\tfrac12,
 \qquad
 E_\omega(\chi)=0+1-\tfrac12=\tfrac12,
 \qquad
 E_\omega(\psi)=0,
\]
and, since \(\max\{\varphi,\chi\}(0)=\max\{\varphi,\chi\}(1)=1\) with
\(|u'|\equiv1\),
\[
 E_\omega(\max\{\varphi,\chi\})=1+1-\tfrac12=\tfrac32 .
\]
Because \(\varphi\geq\psi\) and \(\chi\geq\psi\), one has
\(\max\{\varphi,\psi\}=\varphi\) and \(\max\{\psi,\chi\}=\chi\).  Therefore
\[
 d_\omega(\varphi,\psi)=2\cdot\tfrac12-\tfrac12-0=\tfrac12,
 \qquad
 d_\omega(\psi,\chi)=2\cdot\tfrac12-0-\tfrac12=\tfrac12,
\]
while
\[
 d_\omega(\varphi,\chi)=2\cdot\tfrac32-\tfrac12-\tfrac12=2 .
\]
Thus
\[
 d_\omega(\varphi,\chi)=2
 >1=d_\omega(\varphi,\psi)+d_\omega(\psi,\chi),
\]
and the triangle inequality fails by a factor of two.  All three potentials
are model functions, all energies are finite, and no singular behaviour is
involved: the failure is a feature of the maximum operation, not of the
finite-energy class.

The mechanism is visible in the numbers.  The gauge measures how much energy
must be added to reach a common upper bound, and \(\max\{\varphi,\chi\}\)
has a corner that neither \(\varphi\) nor \(\chi\) has; the corner
contributes Dirichlet energy that is invisible when each of \(\varphi,\chi\)
is compared with the flat \(\psi\).  Replacing the maximum by the
\emph{rooftop} envelope repairs the example --- see
\Cref{ex:ch26-rooftop-repair} --- but the corresponding triangle inequality
is a theorem of a different order and is not proved here.
\end{example}

\begin{definition}[The rooftop gauge]
\label{def:ch26-rooftop-gauge}
Let \(\varphi,\psi\in\cE^1(X,\omega)\) and suppose that the competitor set
\(\{u\in\PSH(X,\omega):u\leq\min\{\varphi,\psi\}\}\) is nonempty, which is
automatic when \(\varphi-\psi\) is bounded.  Let
\(P_\omega(\min\{\varphi,\psi\})\) denote the \(\omega\)-psh envelope of the
upper semicontinuous obstacle \(\min\{\varphi,\psi\}\), which is
\(\omega\)-psh by \Cref{def:ch23-envelope} and
\Cref{cor:ch23-usc-obstacle}, and set
\begin{equation}
\label{eq:ch26-rooftop-gauge}
 d_\omega^{\,\vartriangle}(\varphi,\psi)
 :=E_\omega(\varphi)+E_\omega(\psi)
   -2E_\omega\bigl(P_\omega(\min\{\varphi,\psi\})\bigr).
\end{equation}
\end{definition}

\begin{proposition}
\label{prop:ch26-rooftop-elementary}
Suppose \(P_\omega(\min\{\varphi,\psi\})\in\cE^1(X,\omega)\).  Then
\(d^{\,\vartriangle}_\omega(\varphi,\psi)\) is a nonnegative real number,
symmetric in its arguments and invariant under adding a common constant; it
vanishes when \(\varphi=\psi\); and if \(\varphi\leq\psi\) then
\(d^{\,\vartriangle}_\omega(\varphi,\psi)=E_\omega(\psi)-E_\omega(\varphi)
=d_\omega(\varphi,\psi)\).
\end{proposition}

\begin{proof}
The envelope satisfies \(P_\omega(\min\{\varphi,\psi\})\leq\varphi,\psi\), so
both differences in \eqref{eq:ch26-rooftop-gauge} are nonnegative by
monotonicity of \(E_\omega\).  Symmetry is clear.  Invariance under a common
constant follows from
\(P_\omega(\min\{\varphi+c,\psi+c\})=P_\omega(\min\{\varphi,\psi\})+c\), which
is \Cref{prop:ch23-formal-envelope-calculus}, together with
\eqref{eq:bfj-energy-translation}.  If \(\varphi\leq\psi\) then
\(\min\{\varphi,\psi\}=\varphi\) is already \(\omega\)-psh, so the envelope is
\(\varphi\) itself.
\end{proof}

\begin{warning}{No metric is claimed on \(\cE^1\)}
Neither \eqref{eq:ch26-gauge} nor \eqref{eq:ch26-rooftop-gauge} is asserted
here to satisfy the triangle inequality; for the first,
\Cref{ex:ch26-triangle-fails} shows that it does not.  Consequently the
strong topology of \Cref{sec:ch26-strong} is defined as a topology, and the
words \emph{complete}, \emph{metric}, and \emph{geodesic} are not used of
\(\cE^1(X,\omega)\) anywhere in this book without a proof.  The strong
\emph{metric} structure is developed, in the trivially valued setting where
the envelope property is available in full, in
\Cref{chap:bj-strong-topologies}.
\end{warning}

\section{Strong convergence and compactness of normalized sublevels}
\label{sec:ch26-strong}

\begin{definition}[Strong convergence]
A net \(\varphi_\alpha\in\mathcal E^1(X,\omega)\) converges
\emph{strongly} to \(\varphi\) if it converges in the BFJ weak topology and
\[
                   E_\omega(\varphi_\alpha)\to E_\omega(\varphi).
\]
\end{definition}

Equivalently, the \emph{strong topology} on \(\cE^1(X,\omega)\) is the
coarsest topology refining the BFJ topology for which \(E_\omega\) is
continuous; since \(E_\omega\) is already upper semicontinuous
(\Cref{thm:extended-energy-basic}), it is generated by the weakly open sets
together with the sets \(\{E_\omega>c\}\), \(c\in\mathbf R\).  A net
converges strongly if and only if it converges in that topology, by the net
criterion \Cref{prop:net-criterion}.  We record the elementary consequences.

\begin{proposition}
\label{prop:ch26-strong-basic}
Let \(\varphi_\alpha,\varphi\in\cE^1(X,\omega)\).
\begin{enumerate}
\item\label{itm:ch26-strong:weak}
 Strong convergence implies weak convergence, and the two notions differ:
 \Cref{ex:ch26-usc-strict} converges weakly but not strongly.
\item\label{itm:ch26-strong:monotone}
 A decreasing net with limit in \(\cE^1\) converges strongly
 (\Cref{prop:e1-monotone-continuity}).  For an increasing net
 \(\varphi_\alpha\nearrow\) with regularized supremum
 \(\varphi=(\sup_\alpha\varphi_\alpha)^*\in\cE^1\) one has weak convergence
 and \(\sup_\alpha E_\omega(\varphi_\alpha)\leq E_\omega(\varphi)\), but no
 strong convergence is asserted here.
\item\label{itm:ch26-strong:uniform}
 Uniform convergence of bounded potentials implies strong convergence, by
 \eqref{eq:ch26-energy-lipschitz}.
\item\label{itm:ch26-strong:gauge}
 If \(\varphi_\alpha\to\varphi\) strongly and either
 \(\varphi_\alpha\geq\varphi\) for all \(\alpha\) or
 \(\varphi_\alpha\leq\varphi\) for all \(\alpha\), then
 \(d_\omega(\varphi_\alpha,\varphi)\to0\).
\end{enumerate}
\end{proposition}

\begin{proof}
\cref{itm:ch26-strong:weak} is the definition together with
\Cref{ex:ch26-usc-strict}, where \(E_\omega(\tau_N)=-\tfrac12\) does not
converge to \(E_\omega(0)=0\).  \cref{itm:ch26-strong:monotone} for decreasing
nets is \Cref{prop:e1-monotone-continuity}; for an increasing net
\(\varphi_\alpha\nearrow\), the regularized supremum
\(\varphi:=(\sup_\alpha\varphi_\alpha)^*\) is \(\omega\)-psh by
\Cref{cor:ch23-increasing-net}, which also gives
\(\varphi=\sup_\alpha\varphi_\alpha\) on every dual complex, so that the net
converges to \(\varphi\) in the BFJ topology; and
\(\sup_\alpha E_\omega(\varphi_\alpha)\leq E_\omega(\varphi)\) is
monotonicity of \(E_\omega\) (\Cref{thm:extended-energy-basic}).
\cref{itm:ch26-strong:uniform} is \Cref{prop:ch26-energy-lipschitz}.
\cref{itm:ch26-strong:gauge} is
\Cref{prop:ch26-gauge-elementary}\,\cref{itm:ch26-gauge:ordered}, which for an
ordered pair gives \(d_\omega(\varphi_\alpha,\varphi)
=|E_\omega(\varphi_\alpha)-E_\omega(\varphi)|\).
\end{proof}

\begin{warning}{The two monotone directions are not symmetric}
Continuity of \(E_\omega\) along \emph{decreasing} nets is a theorem
(Step~4 of \Cref{thm:extended-energy-basic}), and its proof uses that model
majorants are continuous and that \(X^{\An}\) is compact.  The corresponding
statement for \emph{increasing} nets is not proved in this book: what
\eqref{eq:extended-energy} gives for free is only the inequality
\(\sup_\alpha E_\omega(\varphi_\alpha)\leq E_\omega(\varphi)\), and the
reverse inequality would require producing, for each \(\varepsilon>0\), a
single member of the net whose energy is within \(\varepsilon\) of
\(E_\omega(\varphi)\).  Upper semicontinuity of \(E_\omega\) is of no help,
since it points in the same direction as monotonicity.  We therefore do not
use increasing nets as a source of strong convergence anywhere in
\Cref{part:bfj}; every strong-convergence argument below is either monotone
decreasing, uniform, or a direct verification of energy convergence.
\end{warning}

\begin{proposition}\label{prop:e1-monotone-continuity}
If \(\varphi_\alpha\searrow\varphi\in\mathcal E^1(X,\omega)\), then
\[
 E_\omega(\varphi_\alpha)\longrightarrow E_\omega(\varphi),
\qquad
 \MA_\omega(\varphi_\alpha)\longrightarrow\MA_\omega(\varphi)
\]
weakly. In particular, the convergence is strong.
\end{proposition}

\begin{proof}
Energy continuity is \Cref{thm:extended-energy-basic}.  Weak convergence of
the measures is the diagonal case of Step~5 of the proof of
\Cref{thm:e1-mixed-products}, whose hypotheses are met: the net is decreasing
with limit in \(\cE^1\), and after subtracting a constant all members are
\(\leq0\).  Explicitly, the two error terms in that estimate are the tails
\eqref{eq:e1-mass-tail} for \(\varphi_\alpha\) and \(\varphi\), bounded
uniformly by \((n+1)(-E_\omega(\varphi))/t\) by
\eqref{eq:ch26-quantitative-tail} and \(E_\omega(\varphi_\alpha)\geq
E_\omega(\varphi)\); and the fixed-level term, which converges by
\eqref{eq:weighted-monotone-ma} after uniform approximation of the continuous
test function by model functions, all measures having the common mass \(V\).
\end{proof}

\begin{proposition}\label{prop:e1-compact-superlevels}
For \(C\geq0\), the set
\begin{equation}
\label{eq:normalized-energy-superlevel}
 \mathcal K_C:=
 \{\varphi\in\mathcal E^1(X,\omega)\mid
   \sup_{X^{\An}}\varphi=0,\ E_\omega(\varphi)\geq-C\}
\end{equation}
is compact in the BFJ weak topology.
\end{proposition}

\begin{proof}
The normalized psh space is compact by
\Cref{thm:bfj-psh-compactness}. The supremum functional is continuous
in the BFJ topology, and \(E_\omega\) is upper semicontinuous. Hence the
conditions \(\sup\varphi=0\) and \(E_\omega(\varphi)\geq-C\) define a closed
subset. The energy bound also ensures that every point of the subset lies in
\(\mathcal E^1\).
\end{proof}

\begin{corollary}
\label{cor:ch26-net-superlevel}
Let \((\varphi_\alpha)\) be a net in \(\cE^1(X,\omega)\) with
\(\sup_{X^{\An}}\varphi_\alpha\) bounded and
\(\inf_\alpha E_\omega(\varphi_\alpha)>-\infty\).  Then some subnet converges
in the BFJ topology, uniformly on every dual complex, to a limit
\(\varphi\in\cE^1(X,\omega)\), and
\begin{equation}
\label{eq:ch26-net-energy-usc}
 E_\omega(\varphi)\ \geq\
 \varlimsup_\alpha E_\omega(\varphi_\alpha).
\end{equation}
Moreover, for every \(t>0\),
\begin{equation}
\label{eq:ch26-uniform-tail-family}
 \begin{aligned}
 \sup_\alpha\ \MA_\omega(\varphi_\alpha)
   \bigl(\{\varphi_\alpha<\sup\varphi_\alpha-t\}\bigr)
 &\leq \frac{(n+1)C'}{t},\\
 C'&:=\sup_\alpha\bigl(-E_\omega(\varphi_\alpha)
   +V\sup_{X^{\An}}\varphi_\alpha\bigr).
 \end{aligned}
\end{equation}
\end{corollary}

\begin{proof}
Normalizing by \(\sup\varphi_\alpha\), which changes energies by a bounded
amount through \eqref{eq:bfj-energy-translation}, we may assume
\(\sup_{X^{\An}}\varphi_\alpha=0\) and
\(E_\omega(\varphi_\alpha)\geq-C\) for a single \(C\).  The net then lies in
the compact set \(\mathcal K_C\) of \eqref{eq:normalized-energy-superlevel},
so \Cref{prop:net-criterion}\,(3) --- equivalently
\Cref{cor:ch22-net-compactness} --- supplies a convergent subnet, and the
mode of convergence is uniform convergence on every dual complex by
\Cref{thm:bfj-skeletal-equicontinuity}.  The limit lies in
\(\mathcal K_C\subseteq\cE^1\) because the set is closed, and
\eqref{eq:ch26-net-energy-usc} is upper semicontinuity of \(E_\omega\) along
the subnet, extended to the whole net by
\Cref{prop:net-semicontinuity}.  The tail bound
\eqref{eq:ch26-uniform-tail-family} is
\eqref{eq:ch26-quantitative-tail} applied to each normalized member, the
normalizing constant being absorbed into \(C'\).

We stress that no sequence is extracted.  The space \(\PSH_0(X,\omega)\) is
compact but is not known to be metrizable, and
\Cref{lem:ch21-no-cofinal-sequence} shows that the underlying system of
models admits no cofinal sequence; the sequential reductions available in
this theory are those of
\Cref{cor:psh-sequential-regularization} and
\Cref{prop:capacity-subsequence}, and neither applies to an arbitrary net in
\(\mathcal K_C\).
\end{proof}

\begin{remark}
The set in \eqref{eq:normalized-energy-superlevel} is a \emph{superlevel} of
the negative energy, although it is often informally called a bounded-energy
sublevel. Its compactness is weak compactness. No claim of strong compactness
is hidden in the terminology.
\end{remark}

\begin{proposition}
\label{prop:ch26-no-strong-compactness}
The set \(\mathcal K_C\) of \eqref{eq:normalized-energy-superlevel} is not
compact for the strong topology when \(C>0\).
\end{proposition}

\begin{proof}
Take \(K=\mathbf C((\pi))\), \(X=\mathbf P^1\), and \(\omega\) as in
\Cref{ex:ch26-usc-strict}, so \(V=1\), and let \(a_1,a_2,\ldots\) and
\(t_N=-\varphi_{a_N}\) be as constructed there.  The sequence
\((\tau_N)\) of that example has \(E_\omega(\tau_N)=-\tfrac12\) and so lies
in \(\mathcal K_C\) only for \(C\geq\tfrac12\); for small \(C\) it must be
rescaled.  Fix \(C>0\), put \(\varepsilon:=\min\{1,2C\}\in(0,1]\) and
\(\tau_N^\varepsilon:=\max\{\varphi_{a_N},-\varepsilon\}
 =-\min\{t_N,\varepsilon\}\).  The outgoing slope at \(\eta\) along \(R_N\)
is \(-1\), cancelling the atom of \(\omega\) there, and the slope jumps from
\(-1\) to \(0\) at \(t=\varepsilon\); so
\(\MA_\omega(\tau_N^\varepsilon)=\delta_{t=\varepsilon\text{ on }R_N}\) and,
by \Cref{prop:ch26-graph-energy},
\(E_\omega(\tau_N^\varepsilon)=-\tfrac\varepsilon2\geq-C\).  Since
\(\sup\tau_N^\varepsilon=0\), the sequence lies in \(\mathcal K_C\), and the
disjointness argument of \Cref{ex:ch26-usc-strict} gives
\(\tau_N^\varepsilon\to0\) weakly.  If a subnet converged
strongly, its weak limit would be \(0\) --- the weak topology being Hausdorff
and the whole sequence converging weakly to \(0\) --- and its energies would
converge to \(E_\omega(0)=0\), contradicting
\(E_\omega(\tau_N^\varepsilon)=-\tfrac\varepsilon2<0\).  Hence no subnet converges strongly, and
\(\mathcal K_C\) is not strongly compact for any \(C>0\).
\end{proof}

\section{Dimension one: Thuillier's space
\texorpdfstring{\(\mathrm W^1\)}{W1} and finite-energy measures}
\label{sec:ch26-thuillier}

The Dirichlet form \eqref{eq:ch26-dirichlet-form} was introduced on model
functions, where it is an intersection number.  On a curve it can be
completed, and the completion is a Hilbert space in which the entire
finite-energy theory takes place.  This is Thuillier's construction
\cite[\S3.5]{Thu05}; we record it because it is the one-dimensional shadow
of everything this chapter and
\Cref{chap:bfj-finite-energy-measures} do, and because the shadow is
Hilbertian while the general theory is not
(\Cref{ex:ch26-triangle-fails}).

Throughout this section \(C\) is a connected, proper, smooth, strictly
\(K\)-analytic curve; \(\mathrm A^0(C)\) is the space of smooth functions in
the sense of \cite[\S3.2]{Thu05}, that is, the functions that factor as
\(\Phi\circ\tau_{\mathcal Y}\) through a piecewise affine \(\Phi\) on the
skeleton of some SNC formal model, and \(\mathrm A^0_c(\Omega)\) is
the subspace of those compactly supported in an open \(\Omega\);
\(\mathrm D^0(C)\) is the space of currents of degree \(0\), canonically the
space of all real functions on the set \(I(C)\) of quasimonomial points with
the topology of pointwise convergence; and \(\ddc\) is the operator of
\Cref{sec:ch08-harmonic}, which on model functions is the one used
throughout this chapter (\Cref{rem:ch26-thuillier-dirichlet}).

\refstepcounter{theorem}\label{thm:ch26-thuillier-w1}
\begin{quote}
\textbf{Foundational input \thetheorem\ (Thuillier's Hilbert space
\(\mathrm W^1\) and finite-energy measures; source boundary).}
Let \(C\) be a connected proper smooth strictly \(K\)-analytic curve and put
\(\langle\varphi,\psi\rangle_{\mathrm D}=-\int_C\varphi\,\ddc\psi\) for
\(\varphi,\psi\in\mathrm A^0(C)\).
\begin{enumerate}
\item\label{itm:ch26-thu:hilbert}
 The form \(\langle\cdot,\cdot\rangle_{\mathrm D}\) is symmetric and
 positive, and its kernel is exactly the constants.  For each
 \(x\in I(C)\) the expression
 \(\langle\varphi,\psi\rangle_x=\varphi(x)\psi(x)
   +\langle\varphi,\psi\rangle_{\mathrm D}\)
 is an inner product on \(\mathrm A^0(C)\); the canonical injection
 \(\mathrm A^0(C)\hookrightarrow\mathrm D^0(C)\) extends to a continuous
 injection of the completion, whose image is a subspace
 \(\mathrm W^1(C)\subseteq\mathrm D^0(C)\) independent of \(x\), and the
 topologies induced by the various \(x\) coincide.
\item\label{itm:ch26-thu:criterion}
 A family \((u_\lambda)\) in \(\mathrm W^1(C)\) converging to \(u\) in
 \(\mathrm D^0(C)\) converges weakly to \(u\) in \(\mathrm W^1(C)\) if and
 only if \(\langle u_\lambda,u_\lambda\rangle_{\mathrm D}\) is bounded.
\item\label{itm:ch26-thu:energy}
 Let \(\Omega\subseteq C^{\An}\) be open with \(\partial\Omega\) a nonempty
 subset of \(I(C)\), and let \(\mathrm W^1(\Omega)\) be the closure of
 \(\mathrm A^0_c(\Omega)\) in \(\mathrm W^1(C)\).  For a positive Radon
 measure \(\mu\) on \(\Omega\) with potential \(g_\mu\) the following are
 equivalent: \(\mu\in\ddc\mathrm W^1(\Omega)\); \(g_\mu\) is subharmonic and
 lies in \(\mathrm W^1(\Omega)\); \(g_\mu\) is subharmonic and lies in
 \(L^1(\rd\mu)\).  When they hold, \(\mu\) is said to have \emph{finite
 energy}, and
 \[
  E_\Omega(\mu)=-\int_\Omega g_\mu\,\rd\mu
   =\langle g_\mu,g_\mu\rangle_{\mathrm D}
 \]
 is the squared \(\mathrm W^1\)-norm of its potential.  The measures of
 finite energy form a convex cone; consequently, for every closed
 \(F\subseteq\Omega\) carrying at least one probability measure of finite
 energy there is a unique energy-minimizing probability measure on \(F\), by
 the projection theorem in \(\mathrm W^1(\Omega)\).
\item\label{itm:ch26-thu:dirichlet-principle}
 \emph{(Dirichlet principle.)}  For a signed measure \(\nu\) of total mass
 zero on \(C^{\An}\),
 \begin{equation}
 \label{eq:ch26-thuillier-dirichlet-principle}
  \sup_{\varphi\in\mathrm A^0(C)}
   \Bigl(2\int_C\varphi\,\rd\nu
         -\langle\varphi,\varphi\rangle_{\mathrm D}\Bigr)
  =E(\nu),
 \end{equation}
 both sides being \(+\infty\) simultaneously; when finite, the supremum is
 attained at the potential \(g\) with \(\ddc g=-\nu\).
\item\label{itm:ch26-thu:polar}
 Every measure of finite energy gives measure zero to every polar set.
\end{enumerate}
Items \cref{itm:ch26-thu:hilbert}--\cref{itm:ch26-thu:polar} are imported in
full from \cite[Propositions~3.5.2, 3.5.7, 3.5.9, 3.5.10 and
Th\'eor\`eme~3.5.11]{Thu05}; no proof is claimed here.
\end{quote}

\begin{remark}[How this matches the general theory]
\label{rem:ch26-thuillier-comparison}
Four points of contact, and one factor.

\emph{Potentials versus measures.}  This chapter's class \(\cE^1\)
(\Cref{def:e1}) consists of \emph{potentials} of finite energy; Thuillier's
\(\mathrm E(\Omega)^+\) consists of \emph{measures} of finite energy.  The
two are exchanged by the Legendre transform of
\Cref{chap:bfj-finite-energy-measures}, and
\cref{itm:ch26-thu:dirichlet-principle} is the curve case of the variational
principle used there.

\emph{The factor.}  Take \(n=1\), let \(\mu\in\mathcal M_V(C)\) be a positive
measure of total mass \(V\), and put \(\nu=\mu_0-\mu\), a signed measure of
total mass zero.  By
\eqref{eq:ch26-energy-dirichlet}, for every model function \(\varphi\),
\[
 E_\omega(\varphi)-\int\varphi\,\rd\mu
 =\int\varphi\,\rd\nu-\tfrac12\langle\varphi,\varphi\rangle_{\mathrm D}
 =\tfrac12\Bigl(2\int\varphi\,\rd\nu
        -\langle\varphi,\varphi\rangle_{\mathrm D}\Bigr).
\]
Taking suprema and comparing with \eqref{eq:dual-energy-measure} and
\eqref{eq:ch26-thuillier-dirichlet-principle} gives
\begin{equation}
\label{eq:ch26-dual-vs-thuillier}
 E^\vee_\omega(\mu)\leq\tfrac12\,E(\mu_0-\mu),
\end{equation}
the factor \(\tfrac12\) coming from \eqref{eq:model-energy} at \(n=1\).  The
inequality is not upgraded to an equality here: the supremum in
\eqref{eq:dual-energy-measure} is over \(\omega\)-psh \emph{model} functions,
whereas that in \eqref{eq:ch26-thuillier-dirichlet-principle} is over all
smooth functions.  What is true, and worth noticing, is that the maximizer
supplied by \cref{itm:ch26-thu:dirichlet-principle} is the potential \(g\)
with \(\ddc g=\mu-\mu_0\), which is \(\omega\)-psh precisely because
\(\mu\geq0\); so the two problems have the same solution whenever that
solution is a model function.  \Cref{sec:ch28-curves} carries out the
corresponding computation with numbers.

\emph{Polar sets.}  Item \cref{itm:ch26-thu:polar} is the curve analogue of
\Cref{prop:ch26-npp-basic}\,\cref{itm:ch26-npp:polar}, but it is proved by a
different route: there the measure is a nonpluripolar product, built layer
by layer so as to discard polar mass, and the statement is an assembly of
\Cref{prop:bounded-ma-no-polar-mass}; here the measure is arbitrary of finite
energy, and the statement follows from \(g_\mu\in L^1(\rd\mu)\), which forces
\(\mu(\{g_\mu=-\infty\})=0\), together with the potential-theoretic fact that
\(\mu(E\cap\{g_\mu>-\infty\})=0\) for every compact polar \(E\) and every
positive Radon measure \(\mu\) \cite[Th\'eor\`eme~3.4.22]{Thu05}.
Thuillier's version applies
to measures that are not Monge--Amp\`ere measures of anything, which is
exactly the generality \Cref{chap:bfj-finite-energy-measures} needs.

\emph{Why the Hilbert structure does not survive.}  Item
\cref{itm:ch26-thu:hilbert} makes \(\mathrm W^1\) a Hilbert space because the
Dirichlet form is bilinear, which by
\Cref{prop:ch26-dirichlet-vertical}\,\cref{itm:ch26-dir:energy} is the same
as saying that \(E_\omega\) is affine-plus-quadratic.  For \(n\geq2\) the
energy is a polynomial of degree \(n+1\)
(\eqref{eq:ch26-energy-intersection}), the gauge
\eqref{eq:ch26-gauge} is not induced by any inner product, and
\Cref{ex:ch26-triangle-fails} shows it is not even a metric.  No completion
theorem of the kind stated in \cref{itm:ch26-thu:hilbert} is available or
claimed in this book outside dimension one.
\end{remark}

\section{Worked calculations}
\label{sec:ch26-examples}

The six computations below are carried to explicit numbers.  Together they
test every normalization of the chapter: the constant \(V\) in
\eqref{eq:bfj-energy-translation}, the sign of the Dirichlet term, the
factor \(n!\) in the toric formula, the tail rate
\eqref{eq:ch26-quantitative-tail}, and the two failures announced in
\Cref{sec:ch26-payoff}.

\subsection{Energy on a metric graph}
\label{sec:ch26-ex-graph}

\begin{proposition}
\label{prop:ch26-graph-energy}
Let \(X\) be a smooth projective curve, let \(\mathcal X\) be a regular
semistable model with skeleton \(\Gamma\) equipped with its modulus lengths,
and let \(\omega\) be a closed model form determined on \(\mathcal X\) with
curvature measure \(\mu_0=\MA_\omega(0)\) of mass \(V=\deg L\).  Then, for
every piecewise affine \(\omega\)-psh function \(u\) on \(\Gamma\),
\begin{equation}
\label{eq:ch26-graph-energy}
 E_\omega(u)=\int_\Gamma u\,\rd\mu_0
             -\frac12\int_\Gamma|u'|^2\,\rd s .
\end{equation}
More generally, for two such functions \(u,v\),
\begin{equation}
\label{eq:ch26-graph-energy-difference}
 E_\omega(v)-E_\omega(u)
 =\int_\Gamma(v-u)\,\rd\mu_0
  -\frac12\int_\Gamma\bigl(|v'|^2-|u'|^2\bigr)\,\rd s .
\end{equation}
\end{proposition}

\begin{proof}
Here \(n=1\), so \eqref{eq:model-energy} reads
\[
 E_\omega(u)=\tfrac12\Bigl(\int u\,\omega+\int u\,\omega_u\Bigr)
 =\int u\,\rd\mu_0+\tfrac12\int u\,\ddc u ,
\]
using \(\omega_u=\omega+\ddc u\) and \(\int u\,\omega=\int u\,\rd\mu_0\),
which is \eqref{eq:curve-ma-laplacian} at \(u=0\).  The graph Dirichlet
identity \eqref{eq:graph-dirichlet-sign} gives
\(\int_\Gamma u\,\Delta_\Gamma u=-\int_\Gamma|u'|^2\rd s\), and
\(\ddc u=\Delta_\Gamma u\) by \eqref{eq:ch31-graph-laplacian}.  This is
\eqref{eq:ch26-graph-energy}, and
\eqref{eq:ch26-graph-energy-difference} follows by subtraction, both
quantities being finite.
\end{proof}

\begin{example}[Energy on the two-vertex graph]
\label{ex:ch26-two-vertex-energy}
Use \Cref{ex:two-vertex-ma}, with edge length one and reference
\(\omega=m_0\delta_{x_0}+m_1\delta_{x_1}\). If
\(\varphi(x_0)=a\), \(\varphi(x_1)=b\), then
\[
 \int\varphi\,\ddc\varphi=-(a-b)^2.
\]
Since \(n=1\), formula~\eqref{eq:model-energy} becomes
\[
 E_\omega(\varphi)
 =\int\varphi\,\omega+\frac12\int\varphi\,\ddc\varphi
 =m_0a+m_1b-\frac12(a-b)^2.
\]
The affine part records the reference mass; the negative quadratic term is
the graph Dirichlet energy. Concavity is now visible without invoking the
Hodge index theorem.

By \Cref{ex:two-vertex-ma} the two Monge--Amp\`ere masses are
\(m_0+b-a\) and \(m_1+a-b\), so the admissible region is the window
\[
 -m_0\ \leq\ b-a\ \leq\ m_1 ,
\]
which is \eqref{eq:ch31-two-component-window}; it collapses to
\(|a-b|\leq\min\{m_0,m_1\}\) only when \(m_0=m_1\), and is otherwise
asymmetric about \(a=b\).  Two further checks are available.  First, the
first variation \eqref{eq:energy-first-variation} along the segment from
\((0,0)\) to \((a,b)\) is
\[
 \frac{\rd}{\rd t}E_\omega(ta,tb)
 =a\bigl(m_0+t(b-a)\bigr)+b\bigl(m_1+t(a-b)\bigr)
 =m_0a+m_1b-t(a-b)^2,
\]
which is indeed \(\int(\varphi-0)\MA_\omega(t\varphi)\).  Second, the second
variation is the constant \(-(a-b)^2\leq0\), which is
\eqref{eq:energy-second-variation} with \(n=1\); it vanishes exactly when
\(a=b\), that is, when \(\varphi\) is constant, in accordance with the
equality clause of \Cref{prop:ch20-relative-hodge}.
\end{example}

\begin{example}[A chain of four components, with numbers]
\label{ex:ch26-chain-energy}
Continue \Cref{ex:ch20-chain-three}: \(\Gamma=[0,3]\) subdivided at the
integers, all edges of length one, \(\mu_0=\delta_0+\delta_3\) and \(V=2\).
Model functions are the piecewise affine \(u\) with vertex values
\((u_0,u_1,u_2,u_3)\), and \(\omega\)-psh means
\eqref{eq:ch20-chain-constraints}.  By \eqref{eq:ch26-graph-energy},
\begin{equation}
\label{eq:ch26-chain-energy}
 E_\omega(u)=u_0+u_3
   -\tfrac12\bigl((u_1-u_0)^2+(u_2-u_1)^2+(u_3-u_2)^2\bigr).
\end{equation}
Three admissible vectors, all normalized by \(\max_iu_i=0\):
\begin{center}
\small
\begin{tabular}{c|c|c|c}
\toprule
\((u_0,u_1,u_2,u_3)\)&\(\MA_\omega(u)\)&\(E_\omega(u)\)&comment\\
\midrule
\((0,0,0,0)\)&\(\delta_{x_0}+\delta_{x_3}\)&\(0\)&the reference\\
\((0,-1,-1,0)\)&\(\delta_{x_1}+\delta_{x_2}\)&\(-1\)&the tent\\
\((0,-1,-2,-2)\)&\(\delta_{x_2}+\delta_{x_3}\)&\(-3\)&one-sided\\
\bottomrule
\end{tabular}
\end{center}
The four left-hand sides of \eqref{eq:ch20-chain-constraints} are precisely
the four masses of \(\MA_\omega(u)\) at the vertices, so the middle column is
read off from them: for the tent they are \((0,1,1,0)\), and for the third
vector \((0,0,1,1)\), all nonnegative, so both vectors are admissible.  The
energies come from \eqref{eq:ch26-chain-energy}: for the tent,
\(0+0-\tfrac12(1+0+1)=-1\); for the third vector,
\(0+(-2)-\tfrac12(1+1+0)=-3\).  In each case the total mass is \(2=V\), the
perturbation \(\ddc u\) having total mass zero.

The energy of the tent may also be obtained from the sandwich
\eqref{eq:energy-sandwich}.  With \(\varphi=(0,-1,-1,0)\leq\psi=(0,0,0,0)\),
\[
 \int(\psi-\varphi)\MA_\omega(\psi)=0+0=0,
 \qquad
 \int(\psi-\varphi)\MA_\omega(\varphi)=1+1=2 ,
\]
so \(0\leq E_\omega(\psi)-E_\omega(\varphi)\leq2\); the true value is \(1\),
the midpoint, as it must be for a quadratic with vanishing linear part along
this segment.  Finally
\(I(\varphi,\psi)=\int(\varphi-\psi)(\MA_\omega(\psi)-\MA_\omega(\varphi))
=0-(-2)=2\) and
\(J_\psi(\varphi)=0-(-1)=1\), \(J_\varphi(\psi)=2-1=1\), confirming
\Cref{thm:ch26-mixed-energy-identities}\,\cref{itm:ch26-ij:sum} and the
bounds \cref{itm:ch26-ij:comparison} with \(n=1\), which read
\(\tfrac12I\leq J\leq\tfrac12I\), that is, equality.
\end{example}

\subsection{An unbounded potential of finite energy}
\label{sec:ch26-ex-unbounded}

\begin{example}[An unbounded \(\cE^1\) potential on \(\mathbf P^1\), with its
energy computed]
\label{ex:ch26-unbounded-finite-energy}
Keep the setting of \Cref{ex:ch21-p1-refinement}: \(K\) complete discretely
valued, \(X=\mathbf P^1_K\), \(L=\cO(1)\), and \(\omega\) the curvature of
the model metric of \((\mathbf P^1_R,\cO(1))\), so that
\(\mu_0=\MA_\omega(0)=\delta_\eta\) at the Gauss point and \(V=1\).  Let
\(R_0=[\eta,0]\) be the ray obtained from the tower of blowups at the closure
of \(\{z=0\}\), parametrized by \(t=-\log r\in[0,\infty)\), the point of parameter $t<\infty$ being the analytic point $\zeta_{0,|\pi|^t}$, which is divisorial when $t$ is rational.

\emph{The potential.}  Let \(s_k:=-1/k\) for \(k\geq1\) and let
\(u:[0,\infty)\to(-\infty,0]\) be the piecewise affine function with
\(u(0)=0\) and slope \(s_k\) on \([k-1,k]\).  Thus
\[
 \begin{aligned}
 u(k)&=-H_k:=-\sum_{j=1}^k\frac1j,\\
 u(k+)-u(k-)\ \text{slope jump}
   &=\Delta_k:=\frac1k-\frac1{k+1}\\
   &=\frac1{k(k+1)} .
 \end{aligned}
\]
Write \(\varphi_0:=\log\lVert z_1\rVert_0=\min\{0,\log|z|\}\), the potential
of \Cref{ex:ch21-p1-refinement}, and put
\(\mathbf t:=-\varphi_0:X^{\An}\to[0,+\infty]\); on \(R_0\) the function
\(\mathbf t\) restricts to the parameter \(t\), and
\(\mathbf t=0\) exactly on \(\{|z|\geq1\}\).  By
\Cref{ex:ch21-p1-refinement} the retraction to the skeleton \([0,m]\) of the
\(m\)-fold blowup satisfies \(\varphi_0\circ p_{\mathcal X_m}
=\max\{\varphi_0,-m\}\), that is,
\(\mathbf t\circ p_{\mathcal X_m}=\min\{\mathbf t,m\}\).  Define
\[
 \Phi_m:=u\bigl(\min\{\mathbf t,m\}\bigr)
 =u_m\circ p_{\mathcal X_m},
 \qquad u_m(t):=u(\min\{t,m\}).
\]
Since \(u_m\) is piecewise affine on \([0,m]\) with rational breakpoints and
slopes, \(\Phi_m\) is a model function determined on \(\mathcal X_m\), and it
is \(\omega\)-psh by the vertex criterion
\eqref{eq:ch21-curve-vertex-inequalities}: at \(t=0\) the only nonzero
outgoing slope is \(s_1=-1\) and \(\mu_0(\{t=0\})=1\), giving
\(1+(-1)=0\geq0\); at an interior integer \(t=k<m\) the two outgoing slopes
sum to \(\Delta_k>0\); at \(t=m\) the single outgoing slope is
\(-s_m=1/m>0\); and elsewhere both terms vanish.  The sequence
\((\Phi_m)_m\) is decreasing, because \(u\) is, so
\(\varphi:=\inf_m\Phi_m=u(\mathbf t)\) is \(\omega\)-psh by
\Cref{thm:ch21-decreasing-stability}, with
\(\varphi(\eta)=u(0)=0\) so that \(\varphi\not\equiv-\infty\).  Its polar set
is \(\{\mathbf t=+\infty\}=\{z=0\}\), a single rational point.

\emph{The Monge--Amp\`ere measures.}  By
\eqref{eq:curve-ma-laplacian},
\[
 \MA_\omega(\Phi_m)
 =\sum_{k=1}^{m-1}\Delta_k\,\delta_{t=k}+\frac1m\,\delta_{t=m},
\]
of total mass \(\sum_{k=1}^{m-1}(\tfrac1k-\tfrac1{k+1})+\tfrac1m=1=V\), as
\eqref{eq:mixed-total-mass} demands.  The truncation
\(\varphi^{\langle H_m\rangle}\) equals \(\Phi_m\), because \(u\) is
decreasing with \(u(m)=-H_m\).  Passing to the increasing limit
\eqref{eq:ch26-npp-definition},
\begin{equation}
\label{eq:ch26-unbounded-measure}
 \MA_\omega(\varphi)=\sum_{k\geq1}\frac{1}{k(k+1)}\,\delta_{t=k},
 \qquad
 \MA_\omega(\varphi)(X^{\An})=\sum_{k\geq1}
   \Bigl(\frac1k-\frac1{k+1}\Bigr)=1=V .
\end{equation}
So \(\varphi\) has full mass, as \Cref{prop:e1-integrability-tail} requires
once finiteness of the energy is known.

\emph{The energy.}  By \eqref{eq:ch26-graph-energy} applied on the skeleton
\([0,m]\), with \(\mu_0=\delta_{t=0}\) and \(\Phi_m(0)=0\),
\[
 E_\omega(\Phi_m)
 =0-\frac12\int_0^m|u'|^2\,\rd t
 =-\frac12\sum_{k=1}^m\frac1{k^2}.
\]
Hence, by \Cref{prop:ch26-truncation-limit},
\begin{equation}
\label{eq:ch26-unbounded-energy}
 E_\omega(\varphi)=\lim_{m\to\infty}E_\omega(\Phi_m)
 =-\frac12\sum_{k\geq1}\frac1{k^2}
 =-\frac{\pi^2}{12}\approx-0.8225 .
\end{equation}
In particular \(\varphi\in\cE^1(\mathbf P^1,\omega)\) and \(\varphi\) is
unbounded.

\emph{Consistency checks.}  First, the integral in
\eqref{eq:e1-integrability}:
\[
 \int(-\varphi)\MA_\omega(\varphi)
 =\sum_{k\geq1}\frac{H_k}{k(k+1)}
 =\sum_{k\geq1}H_k\Bigl(\frac1k-\frac1{k+1}\Bigr)
 =\frac{\pi^2}{6},
\]
the last equality by Abel summation:
\(\sum_{k=1}^NH_k(\tfrac1k-\tfrac1{k+1})
=1-\tfrac{H_{N+1}}{N+1}+\sum_{j=2}^{N+1}j^{-2}\to1+(\tfrac{\pi^2}{6}-1)\).
Since \(n=1\) and \(\int\varphi\,\rd\mu_0=\varphi(\eta)=0\), the general
identity \eqref{eq:model-energy} gives
\(\int(-\varphi)\MA_\omega(\varphi)=(n+1)(-E_\omega(\varphi))\) exactly; and
indeed \(2\cdot\tfrac{\pi^2}{12}=\tfrac{\pi^2}{6}\).  Thus
\eqref{eq:e1-integrability} holds with equality on the right, which is the
sharpness of that inequality.

Second, the tail.  From \eqref{eq:ch26-unbounded-measure},
\[
 \MA_\omega(\varphi)\bigl(\{\varphi\leq-H_m\}\bigr)
 =\sum_{k\geq m}\frac1{k(k+1)}=\frac1m,
\]
while \(H_m\sim\log m\), so the tail decays like \(\rme^{-H_m}\), far faster
than the rate \(O(H_m^{-1})\) guaranteed by
\eqref{eq:ch26-quantitative-tail}.  The general bound is therefore not sharp for this potential. The next example shows why the finite-energy hypothesis is needed.
\end{example}

\subsection{A potential of infinite energy}
\label{sec:ch26-ex-infinite}

\begin{example}[Infinite energy and vanishing nonpluripolar mass]
\label{ex:ch26-infinite-energy}
Keep \(X=\mathbf P^1_K\), \(L=\cO(1)\) and \(\omega\) as above, and let
\[
 \varphi_0:=\log\lVert z_1\rVert_0\in\PSH(\mathbf P^1,\omega),
\]
the potential of \Cref{ex:ch21-p1-refinement}, that is
\(\varphi_0=-\mathbf t\) in the notation of
\Cref{ex:ch26-unbounded-finite-energy}.  On the ray \(R_0\) it is
\(\varphi_0(t)=-t\), it vanishes on \(\{|z|\geq1\}\), and its polar set is
the single rational point \(z=0\).  Its truncations are
\(\varphi_0^{\langle m\rangle}=\varphi_0\circ p_{\mathcal X_m}\), with
\[
 \MA_\omega(\varphi_0^{\langle m\rangle})=\delta_{t=m},
 \qquad
 E_\omega(\varphi_0^{\langle m\rangle})
 =0-\tfrac12\int_0^m1\,\rd t=-\frac m2 .
\]
Hence, by \Cref{prop:ch26-truncation-limit},
\begin{equation}
\label{eq:ch26-infinite-energy}
 E_\omega(\varphi_0)=\lim_{m\to\infty}\Bigl(-\frac m2\Bigr)=-\infty,
\end{equation}
so \(\varphi_0\notin\cE^1(\mathbf P^1,\omega)\).  Moreover the layers
\(\{\varphi_0>-m\}\) carry no mass at all:
\(\MA_\omega(\varphi_0^{\langle s\rangle})(\{\varphi_0>-m\})=0\) for
\(s\geq m\), the whole mass sitting at \(t=s\).  Therefore
\begin{equation}
\label{eq:ch26-zero-npp-mass}
 \MA_\omega(\varphi_0)=0,
\end{equation}
the zero measure: all the mass has escaped to the polar point \(z=0\), where
the nonpluripolar product discards it.  The weak limit of the truncated
measures, by contrast, is \(\delta_{\{z=0\}}\), of mass \(1=V\).  The two
limits differ, and only the first is the object constructed in
\Cref{def:nonpluripolar-ma}.

Three morals are worth extracting.

\emph{(i)}  Full Monge--Amp\`ere mass is not automatic; the failure is total
here, not partial.  This is the precise sense in which
\Cref{prop:e1-integrability-tail} has content.

\emph{(ii)}  The truncations \(\varphi_0^{\langle m\rangle}\) are bounded,
hence in \(\cE^1\), and they converge to \(\varphi_0\) in the BFJ topology
--- they decrease to it pointwise and converge uniformly on every dual
complex, each dual complex meeting \(R_0\) in a bounded segment.  Thus
\(\cE^1\) is not closed in the BFJ topology, proving
\Cref{thm:ch26-e1-structure}\,\cref{itm:ch26-e1:not-closed}.  Their energies
\(-m/2\) do diverge, consistently with monotone continuity
\Cref{prop:e1-monotone-continuity}, which asserts convergence to
\(E_\omega(\varphi_0)=-\infty\); this is why that proposition is stated for
limits in \(\cE^1\) when a \emph{weak convergence of measures} conclusion is
wanted.

\emph{(iii)}  Scaling does not help.  For \(0<\lambda\leq1\) the function
\(\lambda\varphi_0=\lambda\varphi_0+(1-\lambda)\cdot0\) is
\(\omega\)-psh by
\Cref{thm:ch21-stability}\,\cref{itm:ch21-stab:sum}, equals \(-\lambda t\) on
\(R_0\), and the same computation gives
\(E_\omega(\lambda\varphi_0^{\langle m\rangle})=-\lambda^2m/2\to-\infty\) and
\(\MA_\omega(\lambda\varphi_0)=(1-\lambda)\delta_\eta\), of mass
\(1-\lambda<V\).  Thus for every \(\lambda\in(0,1]\) the potential
\(\lambda\varphi_0\) has infinite energy and loses exactly mass $\lambda$. The loss is partial for $0<\lambda<1$ and total for $\lambda=1$.  Comparing with
\Cref{ex:ch26-unbounded-finite-energy}, whose potential also has a single
rational polar point: the polar set does not decide finiteness of energy, as
announced in the warning after \Cref{thm:ch26-e1-structure}.
\end{example}

\begin{example}[The rooftop gauge repairs
\Cref{ex:ch26-triangle-fails}]
\label{ex:ch26-rooftop-repair}
In the notation of \Cref{ex:ch26-triangle-fails}, the obstacle
\(\min\{\varphi,\chi\}(t)=\min\{1-t,t\}\) is concave with value \(0\) at both
endpoints.  An \(\omega\)-psh competitor \(v\) is convex on \([0,1]\), so
\(v(t)\leq(1-t)v(0)+tv(1)\leq0\); hence
\(P_\omega(\min\{\varphi,\chi\})=0=\psi\).  Since
\(\min\{\varphi,\psi\}=\psi=\min\{\psi,\chi\}\) is already \(\omega\)-psh,
\[
 d^{\,\vartriangle}_\omega(\varphi,\chi)
 =\tfrac12+\tfrac12-0=1,
 \qquad
 d^{\,\vartriangle}_\omega(\varphi,\psi)
 =\tfrac12+0-0=\tfrac12
 =d^{\,\vartriangle}_\omega(\psi,\chi),
\]
so the triangle inequality holds with equality for this triple.  This is
consistent with, but of course does not prove, the general triangle
inequality for \eqref{eq:ch26-rooftop-gauge}; no such statement is claimed
here.
\end{example}

\subsection{Toric energy}
\label{sec:ch26-ex-toric}

We now prove the toric energy formula and then compute it on a surface.  The
conventions are those of \Cref{chap:bfj-examples}: for a polarized toric
variety \((X_\Sigma,L_P)\) with polytope \(P\subset M_{\mathbf R}\),
invariant semipositive metrics correspond to convex functions \(F\) on
\(N_{\mathbf R}\) whose subgradient image is contained in \(-P\), with
\(\trop(x)(m)=-\log|\chi^m(x)|\); the Monge--Amp\`ere measure pushes forward
to \(n!\,\MA_{\mathbf R}(F)\) by \eqref{eq:toric-bfj-alexandrov}, of total
mass \(n!\vol_M(P)=(L^n)=V\).

\begin{definition}[Roof function]
\label{def:ch26-roof}
For a convex \(F:N_{\mathbf R}\to\mathbf R\) with subgradient image contained
in \(-P\), the \emph{roof function} of \(F\) is
\begin{equation}
\label{eq:ch26-roof}
 g_F(y):=\inf_{x\in N_{\mathbf R}}
   \bigl(\langle y,x\rangle+F(x)\bigr),
 \qquad y\in P .
\end{equation}
It is concave on \(P\), and \(g_{F+c}=g_F+c\).
\end{definition}

Comparing with \eqref{eq:appD-conjugate}, \(g_F(y)=-F^*(-y)\), where \(F^*\)
is the Legendre transform of \(F\); concavity and finiteness on \(P\) are
\Cref{prop:appD-legendre-tests}.

\begin{theorem}
\label{thm:ch26-toric-energy}
\index{toric variety!Monge--Ampere energy}
Let \(F,F_0\) be convex functions on \(N_{\mathbf R}\) with subgradient
images contained in \(-P\), defining invariant semipositive metrics on
\((X_\Sigma,L_P)\), and suppose \(F-F_0\) is bounded.  Then
\begin{equation}
\label{eq:ch26-toric-energy}
 E(F)-E(F_0)=n!\int_P(g_F-g_{F_0})\,\rd\lambda,
\end{equation}
where \(\lambda\) is the lattice Lebesgue measure on \(M_{\mathbf R}\) and
\(E\) is the Monge--Amp\`ere energy relative to the model form of \(F_0\).
\end{theorem}

\begin{proof}
Put \(F_t:=(1-t)F_0+tF\) for \(t\in[0,1]\); each \(F_t\) is convex with
subgradient image in \(-P\), because the subgradient of a convex combination
is contained in the corresponding combination of the subgradients and \(-P\)
is convex.  By \eqref{eq:energy-first-variation}, in the form supplied by
\Cref{prop:bounded-energy-identities} --- legitimate because \(F-F_0\) is
bounded, so the corresponding potentials are bounded ---
\[
 \frac{\rd}{\rd t}E(F_t)
 =\int(F-F_0)\,\MA_{L}(F_t)
 =n!\int_{N_{\mathbf R}}(F-F_0)\,\MA_{\mathbf R}(F_t),
\]
the second equality being \eqref{eq:toric-bfj-alexandrov}, the integrand
being invariant and hence a function of \(\trop\).

By \Cref{prop:appD-legendre-pushforward}, \eqref{eq:appD-legendre-pushforward},
the Alexandrov measure of \(F_t\) is the pushforward of Lebesgue measure on
its subgradient image \(-P\) under \(\nabla F_t^*\).  Substituting
\(m=-y\), which preserves \(\lambda\), and using
\(g_{F_t}(y)=-F_t^*(-y)\), whence
\(\nabla g_{F_t}(y)=(\nabla F_t^*)(-y)\), this reads
\begin{equation}
\label{eq:ch26-toric-pushforward}
 \int_{N_{\mathbf R}}h\,\MA_{\mathbf R}(F_t)
 =\int_{P}h\bigl(\nabla g_{F_t}(y)\bigr)\,\rd\lambda(y)
\end{equation}
for every bounded Borel \(h\), the gradient existing \(\lambda\)-almost
everywhere on \(P\) because a finite concave function is differentiable
almost everywhere on the relative interior of its domain.

Now fix \(y\in P\) and consider
\[
 h_y(t):=g_{F_t}(y)
 =\inf_{x\in N_{\mathbf R}}
   \bigl(\langle y,x\rangle+(1-t)F_0(x)+tF(x)\bigr),
 \qquad t\in[0,1].
\]
For each fixed \(x\) the expression in brackets is affine in \(t\), so
\(h_y\) is concave; and since the two infimands differ by at most
\(\lVert F-F_0\rVert_\infty\), the function \(h_y\) is Lipschitz with
constant \(\lVert F-F_0\rVert_\infty\).  Being concave and Lipschitz, it is
absolutely continuous, so \(h_y(1)-h_y(0)=\int_0^1h_y'(t)\,\rd t\).  At a
value of \(t\) at which the infimum defining \(h_y(t)\) is attained at a
unique point \(x_t(y)\) --- equivalently at which \(g_{F_t}\) is
differentiable at \(y\), in which case \(x_t(y)=\nabla g_{F_t}(y)\) by the
subgradient description \eqref{eq:appD-subgradient} --- the concave function
\(h_y\) has derivative
\[
 h_y'(t)=(F-F_0)\bigl(\nabla g_{F_t}(y)\bigr),
\]
because the affine function \(t\mapsto\langle y,x_t(y)\rangle+F_t(x_t(y))\)
touches \(h_y\) at \(t\) from above and has slope
\((F-F_0)(x_t(y))\), so that slope is both a supergradient and, by the same
argument applied on the other side, the derivative.

Integrating in \(y\) and applying Fubini --- legitimate because the
integrand is bounded by \(\lVert F-F_0\rVert_\infty\) on the finite measure
space \([0,1]\times P\) --- gives
\[
\begin{split}
 n!\int_P(g_F-g_{F_0})\,\rd\lambda
 &=n!\int_P\int_0^1(F-F_0)
     \bigl(\nabla g_{F_t}(y)\bigr)\,\rd t\,\rd\lambda(y)\\
 &=\int_0^1n!\int_{N_{\mathbf R}}(F-F_0)\,
     \MA_{\mathbf R}(F_t)\,\rd t\\
 &=\int_0^1\frac{\rd}{\rd t}E(F_t)\,\rd t
  =E(F)-E(F_0),
\end{split}
\]
where the second equality is \eqref{eq:ch26-toric-pushforward}.  This is
\eqref{eq:ch26-toric-energy}.
\end{proof}

\begin{example}[Toric roof functions]
For an invariant semipositive metric on a polarized toric variety with
polytope \(P\), choose the roof-function convention in which increasing
metric weights increase the roof \(g\). Then, with lattice measure,
\[
              E_\omega(\varphi_g)-E_\omega(\varphi_{g_0})
              =n!\int_P(g-g_0)\,\rd \lambda.
\]
In dimension one, \(P=[0,d]\) and \(g(s)=c+as\) give
\[
              E(g)-E(0)=cd+\frac a2d^2.
\]
The constant term is \(cV\), exactly as in
\eqref{eq:bfj-energy-translation}. The Alexandrov and mixed-volume viewpoint is
developed for superforms and supercurrents in \cite{Ber26Convex}; here it is
a real/toric model for, not a substitute for, the global BFJ construction.
\end{example}

\begin{example}[A toric surface, carried to numbers]
\label{ex:ch26-toric-surface}
Let \(X=\mathbf P^1\times\mathbf P^1\) with \(L=\cO(1,1)\), so \(n=2\),
\(P=[0,1]^2\), \(\vol_M(P)=1\) and \(V=2!\cdot1=2=(L^2)\).  Let
\(F_0(x_1,x_2)=\max\{0,-x_1\}+\max\{0,-x_2\}\), the support function of
\(-P\), which is the model potential of the standard toric metric.  For
\(s>0\) put
\[
 F_s(x_1,x_2):=\max\{-x_1,-s\}+\max\{-x_2,-s\} .
\]
Each \(F_s\) is convex with subgradient image \(-P\), and \(F_s-F_0\) is
bounded by \(2s\); so \(F_s\) defines a bounded invariant semipositive metric.

\emph{Energy by the roof formula.}  The one-variable roof of
\(f_s(x)=\max\{-x,-s\}\) is
\(g_{f_s}(y)=\inf_x(xy+f_s(x))=s(y-1)\) for \(y\in[0,1]\), the infimum being
attained at \(x=s\); and \(g_{F_0}=0\).  Since \(F_s\) is separable, so is
its roof: \(g_{F_s}(y_1,y_2)=s(y_1-1)+s(y_2-1)\).  Hence, by
\eqref{eq:ch26-toric-energy},
\begin{equation}
\label{eq:ch26-toric-surface-energy}
 E(F_s)-E(F_0)
 =2!\int_{[0,1]^2}\bigl(s(y_1-1)+s(y_2-1)\bigr)\rd y
 =2\Bigl(-\frac s2-\frac s2\Bigr)=-2s .
\end{equation}

\emph{Energy by the first variation, independently.}  Put
\(F_t^{\,\sharp}:=(1-t)F_0+tF_s\); a direct computation gives
\(F_t^{\,\sharp}(x_1,x_2)=f_t(x_1)+f_t(x_2)\) where \(f_t\) has slopes
\(-1\) on \((-\infty,0)\), \(-t\) on \((0,s)\), and \(0\) on \((s,\infty)\).
The subgradient of \(F_t^{\,\sharp}\) is a product of intervals, so by
\eqref{eq:appD-polyhedral-atom} and \eqref{eq:polyhedral-alexandrov-atom},
\[
 \MA_L(F_t^{\,\sharp})
 =2\bigl[(1-t)^2\delta_{(0,0)}+t(1-t)\delta_{(0,s)}
        +t(1-t)\delta_{(s,0)}+t^2\delta_{(s,s)}\bigr],
\]
of total mass \(2\bigl((1-t)+t\bigr)^2=2=V\).  Since
\(F_s-F_0=w(x_1)+w(x_2)\) with \(w(x)=-\min\{\max\{x,0\},s\}\), one has
\((F_s-F_0)(0,0)=0\), \((F_s-F_0)(0,s)=(F_s-F_0)(s,0)=-s\) and
\((F_s-F_0)(s,s)=-2s\).  Therefore
\[
 \frac{\rd}{\rd t}E(F_t^{\,\sharp})
 =\int(F_s-F_0)\MA_L(F_t^{\,\sharp})
 =2\bigl[-2st(1-t)-2st^2\bigr]=-4st ,
\]
and \(\int_0^1(-4st)\rd t=-2s\), confirming
\eqref{eq:ch26-toric-surface-energy}.

\emph{The mixed measures and the functionals \(I,J\).}  Reading off the
coefficients of the binomial expansion
\(\MA_L(F^{\,\sharp}_t)=\sum_j\binom2jt^j(1-t)^{2-j}
\MA_L(F_s^{(j)},F_0^{(2-j)})\) gives
\[
 \MA_L(F_0)=2\delta_{(0,0)},\qquad
 \MA_L(F_s,F_0)=\delta_{(0,s)}+\delta_{(s,0)},\qquad
 \MA_L(F_s)=2\delta_{(s,s)},
\]
each of mass \(2=V\), as \eqref{eq:mixed-total-mass} requires.  Consequently
\[
 I(F_s,F_0)=\int(F_s-F_0)\bigl(\MA_L(F_0)-\MA_L(F_s)\bigr)
 =0-(-4s)=4s,
\]
\[
 J_{F_0}(F_s)=\int(F_s-F_0)\MA_L(F_0)-\bigl(E(F_s)-E(F_0)\bigr)
 =0+2s=2s,
\]
\[
 J_{F_s}(F_0)=\int(F_0-F_s)\MA_L(F_s)-\bigl(E(F_0)-E(F_s)\bigr)
 =4s-2s=2s .
\]
Thus \(J_{F_0}(F_s)+J_{F_s}(F_0)=4s=I(F_s,F_0)\), confirming
\Cref{thm:ch26-mixed-energy-identities}\,\cref{itm:ch26-ij:sum}, and
\[
 \tfrac13I=\tfrac{4s}{3}\ \leq\ 2s\ \leq\ \tfrac23I=\tfrac{8s}{3},
\]
confirming \cref{itm:ch26-ij:comparison} with \(n=2\).  Finally the Dirichlet
coefficients of \eqref{eq:ch26-ak} are
\[
 a_1=-\Bigl(\int(F_s-F_0)\MA_L(F_s)
   -\int(F_s-F_0)\MA_L(F_s,F_0)\Bigr)=-(-4s+2s)=2s,
\]
\[
 a_0=-\Bigl(\int(F_s-F_0)\MA_L(F_s,F_0)
   -\int(F_s-F_0)\MA_L(F_0)\Bigr)=-(-2s-0)=2s,
\]
so \(a_0+a_1=4s=I\) and
\(\tfrac13(2a_0+a_1)=\tfrac13(6s)=2s=J_{F_0}(F_s)\), in agreement with
\Cref{thm:ch26-mixed-energy-identities}\,\cref{itm:ch26-ij:J}.  Every entry of
the theorem has now been verified on an example in which no two of the
measures coincide.
\end{example}

\subsection{A compactness argument carried out}
\label{sec:ch26-ex-compactness}

\begin{example}[Compactness of a normalized energy superlevel, explicitly]
\label{ex:ch26-compactness-worked}
Take \(K=\mathbf C((\pi))\), \(X=\mathbf P^1\), \(L=\cO(1)\), and \(\omega\)
as in \Cref{ex:ch26-unbounded-finite-energy}, so \(V=1\).  For
\(\lambda\in[0,1]\) and \(m\in\mathbf N\) define
\[
 \sigma_{\lambda,m}:=\max\{\lambda\varphi_0,\,-\lambda m\},
\]
where \(\varphi_0\) is the potential of \Cref{ex:ch26-infinite-energy}.  On
the ray \(R_0\) one has \(\sigma_{\lambda,m}(t)=-\lambda\min\{t,m\}\), and
\(\sigma_{\lambda,m}\equiv0\) elsewhere; so
\(\sup\sigma_{\lambda,m}=0\).  By \eqref{eq:ch26-graph-energy},
\[
 E_\omega(\sigma_{\lambda,m})=-\tfrac12\lambda^2m .
\]
Fix \(C>0\).  Then \(\sigma_{\lambda,m}\in\mathcal K_C\) exactly when
\(\lambda^2m\leq2C\).

\emph{A net with an energy bound.}  Consider the sequence
\(\sigma_m:=\sigma_{\lambda_m,m}\) with
\(\lambda_m=m^{-1/2}\), so \(E_\omega(\sigma_m)=-\tfrac12\) for every
\(m\) and \(\sigma_m\in\mathcal K_{1/2}\).  In the notation of
\Cref{ex:ch26-unbounded-finite-energy},
\(\sigma_m=-m^{-1/2}\min\{\mathbf t,m\}\), and on any dual complex
\(\Delta\) the function \(\mathbf t=-\varphi_0\) is finite and continuous,
hence bounded by some \(T_\Delta\) --- this is
\Cref{prop:ch21-first-consequences}\,\cref{itm:ch21-fc:finite}.  Therefore
\(\lVert\sigma_m\rVert_{\infty,\Delta}\leq m^{-1/2}T_\Delta\to0\), so
\(\sigma_m\to0\) uniformly on every dual complex.  By
\Cref{cor:ch26-net-superlevel} the limit is \(0\in\mathcal K_{1/2}\), and
\eqref{eq:ch26-net-energy-usc} reads
\(0=E_\omega(0)\geq-\tfrac12\), strictly.  No subsequence converges strongly.

\emph{A net without an energy bound.}  Now take
\(\lambda_m=1\) for all \(m\), that is \(\sigma_m=\varphi_0^{\langle
m\rangle}\).  Then \(E_\omega(\sigma_m)=-m/2\to-\infty\), the sequence leaves
every \(\mathcal K_C\), and its BFJ limit \(\varphi_0\) is outside
\(\cE^1\).  Compactness of \(\PSH_0(\mathbf P^1,\omega)\) still applies and
produces the limit \(\varphi_0\); what fails is the last clause of
\Cref{prop:e1-compact-superlevels}, and it fails precisely because the energy
bound has been dropped.

\emph{Reading the tail bound.}  For \(\sigma_{\lambda,m}\in\mathcal K_C\) one
has \(\MA_\omega(\sigma_{\lambda,m})
=(1-\lambda)\delta_\eta+\lambda\delta_{t=m}\) and hence
\[
 \MA_\omega(\sigma_{\lambda,m})
   \bigl(\{\sigma_{\lambda,m}<-t\}\bigr)
 =\begin{cases}\lambda,&t<\lambda m,\\ 0,&t\geq\lambda m.\end{cases}
\]
The bound \eqref{eq:ch26-uniform-tail-family} with \(n=1\) asserts that this
is at most \(2C/t\); for \(t<\lambda m\) that reads
\(\lambda\leq2C/t\), which holds because
\(t<\lambda m\) and \(\lambda^2m\leq2C\) give
\(\lambda t<\lambda^2m\leq2C\).  The estimate is therefore sharp along this
family, in the sense that equality is approached as \(t\uparrow\lambda m\).
\end{example}

\section{Scope: hypotheses and later extensions}
\label{sec:ch26-scope}

\subsection{Where each hypothesis enters}
\label{sec:ch26-hypothesis-use}

\begin{center}
\footnotesize
\begin{tabular}{p{.27\textwidth}|p{.21\textwidth}|p{.44\textwidth}}
\toprule
\textbf{Statement}&\textbf{Uses}&\textbf{Comment}\\
\midrule
\Cref{prop:ch26-energy-intersection},
\Cref{prop:energy-cocycle}
 &intersection theory on models; relative Hodge index
 &Discreteness of the valuation is used only through the existence of the
  models of \eqref{eq:ch20-standing}; residue characteristic is irrelevant.
  Regularity of the model is needed for the sign
  \eqref{eq:ch20-mixed-hodge}.\\
\Cref{prop:bounded-energy-identities}
 &\Cref{thm:bounded-ma-extension};
  \Cref{thm:bfj-decreasing-regularization}
 &Decreasing regularization is where
  \(\operatorname{char}k=0\) enters this chapter, through multiplier ideals
  in \Cref{chap:bfj-compactness}.  In positive residue characteristic the
  substitute is the test-ideal technique of \cite{GJKM19}.\\
\Cref{thm:extended-energy-basic}
 &compactness of \(X^{\An}\); continuity of model functions
 &Step~4 uses only that model majorants are continuous and that
  \(X^{\An}\) is compact; it is valid over any complete nontrivially valued
  field once psh functions and decreasing regularization are available.\\
\Cref{prop:e1-integrability-tail},
\Cref{cor:ch26-quantitative-tail}
 &\eqref{eq:energy-cocycle}; positivity of mixed measures
 &Purely formal given the bounded theory; no capacity input.\\
\Cref{lem:ch26-mixed-locality}
 &\eqref{eq:bounded-maximum-principle}
 &Inherits the capacity-theoretic hypotheses of
  \Cref{chap:bfj-capacity}, hence \(\operatorname{char}k=0\) through
  regularization.\\
\Cref{thm:e1-mixed-products}
 &\Cref{lem:ch26-mixed-locality};
  \Cref{lem:ch26-mixed-tail}
 &Ampleness of \(\{\omega\}\) is used only through \(V>0\) and through the
  bounded theory.  Nothing in the proof is special to
  \(\operatorname{char}k=0\) beyond the inputs listed above.\\
\Cref{prop:ch26-strict-monotone}\,\cref{itm:ch26-strict:cont}
 &\Cref{prop:bounded-capacity-domination}
 &Continuity of one competitor is genuinely used; see the warning after
  \Cref{cor:ch26-gauge-nondegenerate}.\\
\Cref{prop:e1-compact-superlevels},
\Cref{cor:ch26-net-superlevel}
 &\Cref{thm:bfj-psh-compactness}
 &Compactness of the normalized psh space carries the full standing
  signature, including smoothness of \(X\), through
  \Cref{chap:bfj-compactness}.\\
\Cref{thm:ch26-toric-energy}
 &\Cref{thm:toric-bfj-alexandrov};
  \eqref{eq:appD-legendre-pushforward}
 &Toric only.  The boundedness hypothesis \(F-F_0\) bounded is used to place
  the potentials in the bounded theory; the unbounded toric case is not
  treated here.\\
\bottomrule
\end{tabular}
\end{center}

\subsection{What later work removes, and what it does not}
\label{sec:ch26-later}

\emph{(E1) Arbitrary complete nontrivially valued fields.}  The definitions
\eqref{eq:model-energy}, \eqref{eq:extended-energy} and
\eqref{eq:nonpluripolar-ma} make sense whenever the model theory of
\Cref{chap:bfj-model-forms} does, and the formal parts of
\Cref{thm:extended-energy-basic} survive.  What does not survive as proved
here is everything resting on decreasing regularization
\Cref{thm:bfj-decreasing-regularization}, and hence on
\eqref{eq:weighted-monotone-ma}: without it, neither
\Cref{prop:bounded-energy-identities} nor
\Cref{thm:e1-mixed-products} has a proof in this book.  Boucksom--Gubler--Martin
work with non-Archimedean volumes over more general fields
\cite{BGM21}, and Boucksom--Eriksson develop energy-type functionals in the
semiample setting \cite[Theorems~7.14 and~8.4]{BE21}; neither yields an
unrestricted extension of the present chapter, and we do not claim one.

\emph{(E2) Positive residue characteristic.}  Only the regularization input
is affected, and it is affected decisively: multiplier ideals are replaced by
asymptotic test ideals in \cite{GJKM19}, whose scope is recorded in
\Cref{chap:bfj-envelopes}.  All statements of the present chapter that do not
invoke \Cref{thm:bfj-decreasing-regularization},
\Cref{thm:bounded-ma-extension} or
\Cref{thm:bounded-maximum-principle} survive verbatim; in particular
\Cref{prop:ch26-energy-intersection}, \Cref{prop:energy-cocycle},
\Cref{lem:ch26-telescope}, \Cref{thm:ch26-mixed-energy-identities} and
\Cref{cor:ch26-cauchy-schwarz}, all of which are intersection theory on a
fixed regular model.

\emph{(E3) The trivially valued theory.}  The trivially valued analogue of
this chapter is \Cref{chap:bj-energy,chap:bj-finite-energy}, and the strong
topology there is genuinely metrizable because the envelope property is
available in full; see \Cref{chap:bj-strong-topologies}.  The two theories are
parallel but not deducible from one another: the present chapter uses a fixed
ground valuation and a special fibre, and the homogeneity that carries the BJ
arguments is absent here.

\begin{warning}{The scope table is about this chapter's proofs}
The entries above describe what the arguments \emph{given here} require.
They are not claims that the corresponding statements are false under weaker
hypotheses.  Where a published extension exists it is named; where none is
named, the honest reading is ``not proved in this book''.
\end{warning}

\section{Dictionary}
\label{sec:ch26-dictionary}

\subsection{Objects and their two descriptions}
\label{sec:ch26-dict-objects}

\begin{center}
\small
\begin{tabular}{p{.29\textwidth}|p{.30\textwidth}|p{.33\textwidth}}
\toprule
\textbf{Algebraic or measure datum}&\textbf{Energy-theoretic object}
 &\textbf{Established at}\\
\midrule
mixed vertical sum
\(\sum_{j=0}^n(D_\varphi\cdot\omega_\varphi^j\cdot\omega^{n-j})\)
 &\((n+1)E_\omega(\varphi)\)
 &\Cref{prop:ch26-energy-intersection},
  \eqref{eq:ch26-energy-intersection}\\
relative Hodge index sign \eqref{eq:ch20-mixed-hodge}
 &concavity of \(E_\omega\) along psh segments
 &\eqref{eq:energy-second-variation},
  \Cref{cor:energy-monotone-concave}\\
Cauchy--Schwarz \eqref{eq:ch20-mixed-cauchy-schwarz}
 &comparison of \(I\) and \(J\); rigidity
 &\Cref{cor:ch26-cauchy-schwarz},
  \Cref{prop:ch26-strict-monotone}\\
positivity of all mixed measures
 &monotonicity of \(E_\omega\) and the sandwich
 &\eqref{eq:energy-sandwich}\\
graph Dirichlet energy \(\int_\Gamma|u'|^2\)
 &\(-2\bigl(E_\omega(u)-\int u\,\rd\mu_0\bigr)\) on a curve
 &\Cref{prop:ch26-graph-energy}\\
integral of the roof over the polytope
 &toric energy, with the factor \(n!\)
 &\Cref{thm:ch26-toric-energy}\\
truncation \(\varphi^{\langle t\rangle}\) and its layers
 &nonpluripolar measure \(\MA_\omega(\varphi)\)
 &\Cref{def:nonpluripolar-ma}, \Cref{lem:ch26-layers}\\
mass escaping to the polar set
 &failure of \(\varphi\in\cE^1\)
 &\Cref{prop:e1-integrability-tail},
  \Cref{ex:ch26-infinite-energy}\\
uniform tail bound over a family
 &weak compactness of \(\mathcal K_C\)
 &\eqref{eq:ch26-quantitative-tail},
  \Cref{prop:e1-compact-superlevels}\\
\bottomrule
\end{tabular}
\end{center}

\subsection{Reading the dictionary backwards}
\label{sec:ch26-dict-backward}

Four entries are one-way, and it is worth saying which.

\emph{From energy to the potential} there is no arrow: two distinct
potentials may have the same energy, and \(E_\omega\) is not injective even
modulo constants --- \Cref{ex:ch26-two-vertex-energy} shows that the level
sets of \eqref{eq:ch26-two-vertex-energy} are ellipses in the
\((a,b)\)-plane.  What is true is the rigidity statement
\Cref{prop:ch26-strict-monotone}, and only along an ordered pair.

\emph{From full Monge--Amp\`ere mass to finite energy} there is no arrow
either.  Full mass is a consequence of finite energy
(\Cref{prop:e1-integrability-tail}) and not equivalent to it; the class of
full-mass potentials is strictly larger, as it already is in the complex
theory \cite{GZ07}.

\emph{From weak convergence to strong convergence} there is no arrow, by
\Cref{ex:ch26-usc-strict}; the missing ingredient is convergence of energy,
which is a genuinely additional hypothesis and is what
\Cref{prop:ch26-no-strong-compactness} exploits.

\emph{From the energy gauge to a metric} there is no arrow, by
\Cref{ex:ch26-triangle-fails}.

\subsection{Numbers to remember}
\label{sec:ch26-dict-numbers}

\begin{center}
\small
\begin{tabular}{p{.40\textwidth}|p{.22\textwidth}|p{.30\textwidth}}
\toprule
\textbf{Quantity}&\textbf{Value}&\textbf{Where}\\
\midrule
energy on the two-vertex graph
 &\(m_0a+m_1b-\tfrac12(a-b)^2\)
 &\eqref{eq:ch26-two-vertex-energy}\\
energy of the tent on a four-chain
 &\(-1\)
 &\Cref{ex:ch26-chain-energy}\\
energy of the harmonic-slope tower potential
 &\(-\pi^2/12\)
 &\eqref{eq:ch26-unbounded-energy}\\
its Monge--Amp\`ere integral
 \(\int(-\varphi)\MA_\omega(\varphi)\)
 &\(\pi^2/6\)
 &\Cref{ex:ch26-unbounded-finite-energy}\\
energy of \(\log\lVert z_1\rVert_0\) on \(\mathbf P^1\)
 &\(-\infty\)
 &\eqref{eq:ch26-infinite-energy}\\
its nonpluripolar mass
 &\(0\)
 &\eqref{eq:ch26-zero-npp-mass}\\
toric energy of \(F_s\) on \(\mathbf P^1\times\mathbf P^1\)
 &\(-2s\)
 &\eqref{eq:ch26-toric-surface-energy}\\
\((I,J_{F_0}(F_s),J_{F_s}(F_0))\) there
 &\((4s,2s,2s)\)
 &\Cref{ex:ch26-toric-surface}\\
gauge defect in \Cref{ex:ch26-triangle-fails}
 &\(2\) versus \(1\)
 &\Cref{ex:ch26-triangle-fails}\\
\bottomrule
\end{tabular}
\end{center}

\section*{Source notes}

The Monge--Amp\`ere energy, the class \(\mathcal E^1\), the nonpluripolar
mass, monotone continuity, locality, comparison, and differentiability along
psh segments are \cite[\S6]{BFJ15}, in the setting of the original BFJ
signature \eqref{eq:ch20-standing}; the model and bounded inputs quoted
throughout are \cite[Theorem~3.1]{BFJ15} and
\cite[Theorem~5.1 and Corollaries~5.2--5.3]{BFJ15}, and the upper
semicontinuity and monotone continuity of \(E_\omega\) are
\cite[Proposition~6.2]{BFJ15}.  Differentiability along \(\cE^1\) segments,
which is proved in \Cref{chap:bfj-differentiability} rather than here, is
\cite[Proposition~6.12]{BFJ15}.  The underlying model formalism --- closed
model forms, the operator \(\ddc\), semipositivity, and the relative Hodge
index inequality --- is \cite[\S\S2 and 4--5]{BFJ16}.

\paragraph{Provenance, statement by statement.}
\Cref{def:model-energy} is the non-Archimedean instance of the
Aubin--Mabuchi energy; its normalization is fixed once and for all in
\eqref{eq:appJ-energy}, and \Cref{prop:ch26-energy-intersection} is the
identification with a mixed vertical \((n+1)\)-fold intersection sum that
\Cref{chap:bfj-model-forms} announced after
\Cref{cor:ch20-concavity} and left unproved there.  The cocycle
\eqref{eq:energy-cocycle} is \eqref{eq:appJ-energy-cocycle}; the first
variation \eqref{eq:energy-first-variation} is
\eqref{eq:appJ-energy-variation}.  The functionals \(I\) and \(J\) of
\Cref{def:ch26-ij} are the non-Archimedean analogues of Aubin's, and the
identities of \Cref{thm:ch26-mixed-energy-identities}, including the
comparison \(\frac1{n+1}I\leq J\leq\frac n{n+1}I\), reproduce the standard
complex statements found in \cite[\S2]{BEGZ10} and, in the form used for
finite-energy classes, in \cite{GZ07}.  We have proved them here from the
telescoping identity \eqref{eq:ch26-ak-telescope} rather than quoting them,
because the sign of every term is the relative Hodge index inequality
\eqref{eq:ch20-mixed-hodge} and it is worth seeing that dependence
explicitly.

\Cref{lem:ch26-mixed-locality} is the polarization step that turns the
diagonal maximum principle \eqref{eq:bounded-maximum-principle} into a mixed
one; the argument is elementary but is often left implicit.
\Cref{lem:ch26-mixed-tail} and its constant \(n^{n+1}(n+1)/n!\) are the
quantitative form of the full-mass statement; the complex prototype is
\cite[Proposition~1.6]{BEGZ10}, whose proof is organized differently.
\Cref{thm:e1-mixed-products} is \cite[\S6]{BFJ15}; the proof given here is
self-contained and isolates the three ingredients so that each can be
audited against the hypothesis table of \Cref{sec:ch26-hypothesis-use}.

\paragraph{Dimension one, and Thuillier.}
\Cref{sec:ch26-dirichlet-vertical} and \Cref{sec:ch26-thuillier} import
Thuillier's potential theory on Berkovich curves \cite{Thu05}, which
predates the BFJ theory and is its one-dimensional ancestor.  The Dirichlet
form \eqref{eq:ch26-dirichlet-form} is \cite[\S3.5]{Thu05}, where it is
introduced on compactly supported smooth functions and where it is stated
that, for a discrete valuation, it generalizes the intersection form on
vertical divisors of integral models;
\Cref{prop:ch26-dirichlet-vertical} is that statement in the book's
normalizations, with the sign made explicit, and its proof is nothing but
\eqref{eq:ch20-ibp} together with the relative Hodge index inequality
\eqref{eq:ch20-relative-hodge}.  The Hilbert space \(\mathrm W^1\), the
characterization of finite-energy measures by
\(g_\mu\in L^1(\rd\mu)\), the Dirichlet principle
\eqref{eq:ch26-thuillier-dirichlet-principle} and the fact that
finite-energy measures do not charge polar sets are, in order,
\cite[Propositions~3.5.2, 3.5.7, 3.5.9, 3.5.10 and
Th\'eor\`eme~3.5.11]{Thu05}; the archimedean model of the whole construction
is the space of Green functions of Sobolev regularity, written
\(\mathrm L^2_1(M,\mathbf R)\) on a compact Riemann surface \(M\), that
Bost uses in his potential theory for arithmetic surfaces, and which is
Thuillier's own reference for it.
Normalizations agree on the non-Archimedean side and differ by a factor
\(2\) on the archimedean side; this is audited once, in
\Cref{rem:ch25-thuillier-normalization}, and repeated in
\Cref{rem:ch26-thuillier-dirichlet}.  The capacity half of Thuillier's
Chapter~3, which is what \Cref{chap:bfj-capacity} touches, is presented in
\Cref{sec:ch25-thuillier}; one thing worth carrying across is that he uses
\emph{no} quasicontinuity of subharmonic functions anywhere, whereas the
present chapter integrates unbounded potentials against Monge--Amp\`ere
measures only through \Cref{thm:e1-quasicontinuity}.

\paragraph{What is organized differently here.}
Three choices depart from the arrangement in the sources.  First, the mixed
energies \(I\) and \(J\) have been given a section of their own
(\Cref{sec:ch26-mixed}) and every application of integration by parts has
been labelled with the regularity at which it is legitimate
(\Cref{prop:ch26-ibp-regimes}); in the literature the chain is usually
performed inline, and the distinction between the model, bounded and
finite-energy regimes is often left implicit.  Second, the extension of mixed
products to \(\cE^1\) is proved, not defined: \Cref{sec:ch26-mixed-products}
supplies the locality lemma, the tail estimate and the monotone limit
separately, in accordance with the governing convention that model,
continuous, bounded and finite-energy products are never conflated
(\Cref{sec:ch26-four-products}).  Third, the quantitative structure of
\(\cE^1\) is presented honestly as a \emph{gauge} rather than a metric.  The
counterexample \Cref{ex:ch26-triangle-fails} appears to have no counterpart
in the literature; it is elementary, but it is the reason this chapter does
not speak of a distance.  The corresponding metric theory, where the rooftop
envelope replaces the maximum and the triangle inequality is a theorem,
belongs to the trivially valued setting of
\Cref{chap:bj-strong-topologies} and to \cite[\S12]{BJ22}; the synthetic
account of finite-energy measures in \cite{BJ25FiniteEnergy} isolates what
such a theory needs.

\paragraph{Terminology and scope.}
The term \emph{strong convergence} used here means weak convergence plus
convergence of energy; no later metric-completion theorem is imported into
the original BFJ setting, and \Cref{prop:ch26-no-strong-compactness} records
that the compactness proved in \Cref{prop:e1-compact-superlevels} is weak and
nothing more.  The toric formula \Cref{thm:ch26-toric-energy} is the
non-Archimedean shadow of the classical convex-geometric computation; the
Alexandrov and mixed-volume viewpoint for superforms and supercurrents is
developed in \cite{Ber26Convex}, and the comparison of real and
non-Archimedean Monge--Amp\`ere operators used to justify
\eqref{eq:toric-bfj-alexandrov} is \cite{Vil20} together with
\cite{BGJKK25}.  For the complex counterparts of the finite-energy classes
treated here, the standard references are \cite{GZ07}, \cite{BEGZ10} and,
for the local theory, \cite{Ceg98,Ceg04}; the variational method that this
chapter prepares is \cite{BBGZ13} on the complex side and
\cite[\S\S7--8]{BFJ15} on the non-Archimedean side.

\paragraph{Interfaces.}
The boundary between the products of this chapter and the local
Bedford--Taylor products of the Chambert--Loir--Ducros calculus is drawn in
\Cref{sec:ch16-local-vs-global,sec:ch16-delimitation-table}, and the loci on
which the two agree are listed in \Cref{sec:ch16-transfer}.  The comparison
theorems themselves are \Cref{thm:bfj-cld-ma-common-locus} and
\Cref{chap:compare-ma}.  Energy enters the variational equation in
\Cref{chap:bfj-differentiability}, its Legendre dual in
\Cref{chap:bfj-finite-energy-measures}, and the explicit curve and toric
computations are continued in \Cref{chap:bfj-examples}.
